%%%%%%%%%%%%%%%%%%%%%%%%%%%%%%%%%%%%%%%%%%%%%%%%%%%%%%%%%%%%%%%%%%%%%%
% Lớp 12 | Nguyên hàm, tích phân | Ứng dụng thực tế và hình học của tích phân
% Dạng 2D4?3-5: Bài toán thực tế và ứng dụng thể tích tròn xoay, 𝑆(𝑥)
% Số câu: 70  (sinh tự động bởi tach_cau.py; từ nay thầy sửa trực tiếp file này)
%%%%%%%%%%%%%%%%%%%%%%%%%%%%%%%%%%%%%%%%%%%%%%%%%%%%%%%%%%%%%%%%%%%%%%

%%%%%%%%%%%%%%%% TRẮC NGHIỆM 4 PHƯƠNG ÁN %%%%%%%%%%%%%%%%

% nguồn: Phan1_Lop_12_Dai_so_Chuong_4
\begin{ex}%[2D4H3-5]%[ma:17ab5b10]
	Thể tích hình khối do hình  phẳng giới hạn bởi các đường $y=x^2-4$, $y=2x-4$, $x=0$, $x=2$ quay quanh trục $Ox$ bằng:
	\choice 
		 {$\dfrac{416\pi}{15}$}
		 {$\dfrac{752\pi}{5}$}
		 {$\dfrac{16\pi}{15}$}
		{\True $\dfrac{32\pi}{5}$} 
		\begin{center}
			\begin{tikzpicture}[line cap=butt,line join=miter,>=stealth]
					\tikzset{declare function={f  (\x)=(\x)^2-4;
						g (\x)=2*(\x)-4;}} 
					\begin{scope}[overlay]
						 \path[fill=gray!45] plot[domain=0:2] (\x,{f (\x)})
						plot[domain=2:0] (\x,{g (\x)})--cycle; 
					\end{scope}
					 \draw[->] (-2,0)--(3.5,0)
						 node[shift={(-100:7pt)},font=\normalsize]{$x$};
					 \draw[->] (0,-5)--(0,5.3)
						node[shift={(180:7pt)},font=\normalsize]{$y$}; 
					\draw (5pt,0)  |- (0,5pt)
						(0,0) node[shift={(225:9pt)},font=\normalsize]{$O$}; 
					\begin{scope}
						\clip  (-1,-5) rectangle (3.5,5);
						\draw[samples=100] plot[domain=-1:3]  (\x,{f (\x)});
						 \draw[samples=100] plot[domain=-1:3] (\x,{g (\x)});
						\draw[densely dashed,samples=100] plot[domain=-1:3]  (\x,{-f (\x)});
						\draw[densely dashed,samples=100]  plot[domain=-1:3] (\x,{-g (\x)});
						\draw[densely  dashed,samples=100] (-0.03,0) ellipse (0.4 and 3.97);
					\end{scope}
					 \path (1.65,0) node[above]{$2$};
				\end{tikzpicture}
	\end{center}
	\loigiai{ 
		Thể  tích của hình khối là: $V=\pi \displaystyle\int_0^2\left|\left(x^2-4\right)^2-(2x-4)^2\right|\mathrm{d}x=\dfrac{32\pi}{5}$.
	} 
\end{ex}

% nguồn: Phan1_Lop_12_Dai_so_Chuong_4
\begin{ex}%[2D4H3-5]%[ma:01d2e982]
	Biết một nguyên hàm của hàm số $f(x)=\sin^2x$ là $F(x)=\dfrac{x}{2}-\dfrac{\sin 2x}{4}$. Thể tích của khối tròn xoay giới hạn bởi đồ thị  hàm số $y=\sin x$, trục hoành và hai đường thẳng $x=0$, $x=\pi$ khi quay quanh trục $Ox$ là:
	 \choice
		{$\dfrac{\pi ^2}{4}$} 
		 {\True $\dfrac{\pi ^2}{2}$}
		 {$\dfrac{\pi}{2}$}
		 {$\dfrac{\pi}{4}$}
		\begin{center}
			\begin{tikzpicture}[line cap=butt,line join=miter,>=stealth]
					\tikzset{declare function={f (\x)=sin (\x r); xb=3.14/2; yb=f (xb);}} 
					\begin{scope}
						\clip  (-0.5,-1.5) rectangle (3.5,1.5);
						\fill[gray!35]  (0,0)--plot[domain=0:3.14,samples=300,smooth] (\x,{f (\x)})--cycle;
						\draw plot[domain=-0.5:3.64,samples=300,smooth] (\x,{f  (\x)});
						\draw[dashed] plot[domain=-0.5:3.64,samples=300,smooth] (\x,{-f  (\x)});
					\end{scope}
					 \draw[->] (-0.5,0)--(4,0)
						node[shift={(-100:7pt)},font=\normalsize]{$x$}; 
					\draw[->]  (0,-1.5)--(0,1.6)
						 node[shift={(180:7pt)},font=\normalsize]{$y$};
					 \draw (5pt,0) |- (0,5pt)
						(0,0)  node[shift={(245:13pt)},font=\normalsize]{$O$};
					\draw  (xb,yb) arc (90:270:{0.25*yb} and yb);
					\draw[dashed] (xb,yb) arc (90:-90:{0.25*yb} and yb); 
					\path (xb,yb) node[above,font=\scriptsize]{$y=\sin x$}; 
					 \fill (3.14,0) circle (1pt) node[below,font=\scriptsize]{$\pi$};
				\end{tikzpicture}
	\end{center}
	 \loigiai{
		Thể tích khối tròn xoay là: $V=\pi \displaystyle\int_0^{\pi}(\sin x)^2\mathrm{d}x=\pi  \left(\displaystyle\int_0^{\pi}\dfrac{1-\cos 2x}{2} \mathrm{d}x \right)=\dfrac{\pi ^2}{2}\cdot$
	} 
\end{ex}

% nguồn: Phan1_Lop_12_Dai_so_Chuong_4
\begin{ex}%[2D4V3-5]%[ma:260ea38e]
	 Tính thể tích chứa được của một cái chậu inox to mà khách hàng đặt theo kích thước yêu cầu, biết phần trong của nó có dạng khối tròn xoay được tạo thành khi quay hình phẳng giới hạn bởi đường $y=\sqrt{x}+1$, trục $Ox$ và các đường thẳng $x=0$, $x=2$ quanh trục $Ox$, đơn vị trên trục là decimet (làm tròn kết quả đến hàng phần trăm).
	\choice 
		 {$12{,}12$ $\mathrm{~dm}^3$}
		 {$12{,}21$ $\mathrm{~dm}^3$}
		{$24{,}14$  $\mathrm{~dm}^3$}
		{\True $24{,}41$  $\mathrm{~dm}^3$}
		\begin{center}
			\begin{tikzpicture}[line cap=butt,line join=miter,>=stealth]
					\tikzset{declare function={f  (\x)=sqrt (\x)+1; xa=2; ya=f (xa);}}
					\begin{scope}
						\clip  (-1,-3) rectangle (3,3);
						\fill[gray!35] plot[domain=0:2] (\x,{f  (\x)})--(2,0)--(0,0)--cycle;
						\draw[samples=100] plot[domain=0:3] (\x,{f  (\x)});
						\draw[samples=100,dashed]  plot[domain=0:3] (\x,{-f (\x)});
						\draw (2,0) node[below]{$2$}--(2,ya); 
					\end{scope}
					 \draw[->] (-1,0)--(3.5,0)
						node[shift={(-100:7pt)},font=\normalsize]{$x$}; 
					\draw[->]  (0,-3)--(0,3.5)
						node[shift={(180:7pt)},font=\normalsize]{$y$}; 
					\draw (5pt,0) |- (0,5pt) 
						(0,0) node[shift={(240:11pt)},font=\normalsize]{$O$}; 
					\draw[densely dashed]  (0,0) ellipse ({0.2} and 1)
						({xa},0) ellipse ({0.15*ya} and {ya}); 
					 \path (xa+0.5,ya) node[above left]{$y=\sqrt{x}+1$};
				\end{tikzpicture}
	\end{center}
	\loigiai{ 
		Thể tích của chậu inox  là: $V=\pi \displaystyle\int_0^2\left(\sqrt{x}+1\right)^2\mathrm{d}x=\pi \left(\dfrac{8}{3} \sqrt{2}+4\right)\approx 24,41$ ($\mathrm{dm}^3$).
	} 
\end{ex}

% nguồn: Phan1_Lop_12_Dai_so_Chuong_4
\begin{ex}%[2D4V3-5]%[ma:5919e035]
	Tính thể tích chứa được (dung tích) của một cái chén (bát), biết phần trong của nó có dạng khối tròn xoay được tạo thành khi quay quanh trục $Ox$ hình phẳng giới hạn bởi đường $y=\sqrt{2x}+2$ và trục $Ox$ (như hình  vẽ), bát có độ sâu $5$ cm, đơn vị trên trục là centimet (làm tròn kết quả đến hàng đơn vị).
	 \choice
		{$78$ cm $^3$} 
		{\True $274$ cm $^3$} 
		{$87$  cm $^3$}
		{$247$  cm $^3$}
		\begin{center}
			\resizebox{!}{4cm}{\begin{tikzpicture}[line cap=butt,line join=miter,>=stealth,scale=0.6]
		\tikzset{declare function={f (\x)=sqrt (2*\x)+2; xa=5; ya=f  (xa);}}
		\begin{scope}
			\clip (-1,-ya-1) rectangle  (xa+1,ya+1);
			 \fill[gray!35] plot[domain=0:xa] (\x,{f (\x)})--(xa,0)--(0,0)--cycle;
			\draw[samples=100] plot[domain=0:{xa+1}] (\x,{f  (\x)});
			\draw[samples=100,dashed] plot[domain=0:{xa+1}] (\x,{-f  (\x)});
			\draw  (xa,0)--(xa,ya);
			\path (2.5,0) node[below]{$5$  cm};
		\end{scope}
		\draw[->] (-1,0)--({xa+1.5},0) 
			node[shift={(-100:7pt)},font=\normalsize]{$x$}; 
		\draw[->] (0,-ya-1)--(0,ya+1) 
			 node[shift={(180:7pt)},font=\normalsize]{$y$};
		\draw (5pt,0)  |- (0,5pt)
			(0,0)  node[shift={(120:10pt)},font=\normalsize]{$O$};
		\draw[densely dashed] (0,0)  ellipse ({0.5} and 2)
			({xa},0) ellipse ({0.15*ya} and  {ya});
		 \path (xa+0.5,ya) node[above left]{$y=\sqrt{2x}+2$};
	\end{tikzpicture}}
	\end{center}
	 \loigiai{
		 Thể tích của chén là: $V=\pi \displaystyle\int_0^5\left(\sqrt{2x}+2\right)^2\mathrm{d}x=\pi \left(\dfrac{40}{3} \sqrt{10}+45\right)\approx 274$ (cm$^3$)
	 }
\end{ex}

% nguồn: Phan1_Lop_12_Dai_so_Chuong_4
\begin{ex}%[2D4V3-5]%[ma:fe1a5990]
	Một bác thợ gốm làm một cái chậu trồng  cây, phần trong chậu cây có dạng khối tròn xoay được tạo thành khi quay hình phẳng được tô đậm như hình sau quanh trục $Ox$ (đơn vị trên trục là decimet), biết đường cong trong hình là đồ thị của hàm số $y=\sqrt{x+1}$, đáy chậu và miệng chậu có đường kính lần lượt là $2$ dm và $4$ dm. Dung tích của chậu là bao nhiêu (làm tròn kết quả nếu cần thiết)?
	 \choice
		 {$25{,}13$ lít}
		{\True  $23{,}56$ lít}
		 {$14{,}66$ lít}
		 {$7{,}5$ lít}
		\begin{center}
			\begin{tikzpicture}[line cap=butt,line join=miter,>=stealth]
					 \tikzset{declare function={f (\x)=sqrt (\x+1); xa=3; ya=f (xa);}}
					\begin{scope}
						\clip  (-1,-ya-1) rectangle (xa+1,ya+1);
						\fill[gray!35]  plot[domain=0:xa] (\x,{f (\x)})--(xa,0)--(0,0)--cycle;
						\draw[samples=100] plot[domain=0:{xa+1}] (\x,{f  (\x)});
						\draw[samples=100,dashed] plot[domain=0:{xa+1}] (\x,{-f (\x)}); 
						 \draw (xa,0)--(xa,ya);
					\end{scope}
					\draw[->]  (-1,0)--({xa+1},0)
						node[shift={(-100:7pt)},font=\normalsize]{$x$}; 
					\draw[->]  (0,-ya-1)--(0,ya+1)
						node[shift={(180:7pt)},font=\normalsize]{$y$}; 
					\draw (5pt,0) |- (0,5pt) 
						 (0,0) node[shift={(240:10pt)},font=\normalsize]{$O$};
					\draw[densely dashed] (0,0) ellipse ({0.25} and  1)
						({xa},0) ellipse ({0.15*ya} and  {ya});
					\path  (xa+0.5,ya) node[above left]{$y=\sqrt{x+1}$};
					\draw[dashed] (0,1)--+(-2,0) 
						(0,-1)--+(-2,0); 
					\draw[<->] (-1.25,-1)--+(0,2) node[left,pos=0.5]{$2$  dm};
					\draw[dashed] (3,2)--+(2,0) 
						 (3,-2)--+(2,0);
					\draw[<->] (4.25,-2)--+(0,4) node[right,pos=0.5]{$4$ dm}; 
				\end{tikzpicture}
	\end{center}
	\loigiai{ 
		Với bán kính đáy chậu là $1$ dm  thì $y=1\Rightarrow \sqrt{x+1}=1\Rightarrow x=0$.\\
		Với bán kính  đáy chậu là $2$ dm thì $y=2\Rightarrow \sqrt{x+1}=2\Rightarrow x=3$.\\
		Thể tích của khối chậu là: $V=\pi \displaystyle\int_0^2\left(\sqrt{x+1}\right)^2\mathrm{d}x=\dfrac{15\pi}{2} \approx  23{,}56\mathrm{~dm}^3=23{,}56$ lít.
	 }
\end{ex}

% nguồn: Phan1_Lop_12_Dai_so_Chuong_4
\begin{ex}%[2D4V3-5]%[ma:5d599f0c]
	Từ một quả cầu bằng đá trắng sứ bán kính bằng 1 dm, người ta khoan rút lõi ngay \lq\lq chính giữa\rq\rq\ quả cầu (trục đối xứng của lõi và quả cầu trùng nhau) như hình sau với đường kính mũi khoan là 1 dm được một  vật thể có thể tích $V$ là bao nhiêu (bỏ qua độ dày mũi khoan)?
	\choice 
		{$V=\dfrac{\sqrt{3}}{2} \pi$ dm $^3$} 
		{\True $V=\dfrac{8-3\sqrt{3}}{6}  \pi$ dm $^3$}
		{$V=\dfrac{8+3\sqrt{3}}{6} \pi$ dm  $^3$}
		{$V=\dfrac{3\sqrt{3}}{2} \pi$ dm $^3$} 
		\begin{center}
			\resizebox{!}{5cm}{\begin{tikzpicture}[line cap=round,line join=round]
		\fill (1.8,0.57).. controls (1.87,0.57) and  (1.91,0.53)..
			(1.96,0.45).. controls +(-63:0.1) and +(76:0.3).. (1.95,0.03) -- (1.88,-0.27) -- (1.97,-0.33).. controls +(-36:0.2) and +(114:0.1).. (2.47,-0.82).. controls +(-62:0.1) and +(76:0.3).. (2.48,-1.68).. controls +(-96:0.4) and +(41:0.3).. (1.86,-2.48).. controls +(-140:0.2) and +(75:0.2).. (1.72,-3.05).. controls +(-113:0.1) and +(-18:0.2).. (0.89,-3.25) -- (0.89,-3.25).. controls +(146:0.2) and +(-49:0.2).. (0.43,-2.78).. controls +(-179:0.8) and +(-24:0.8).. (-1.94,-2.4).. controls +(153:0.2) and +(0:0.1).. (-2.35,-2.27)..  controls +(180:0.1) and +(-82:0.1).. (-2.52,-2.12).. controls +(93:0.3) and +(-92:0.3).. (-2.62,-1.11).. controls +(86:0.3) and +(-146:0.4).. (-2.2,-0.41).. controls +(35:0.1) and +(180:0).. (-2,-0.28).. controls +(45:0) and +(-80:0.2).. (-2.05,0.04).. controls +(102:0.3) and +(-119:0.1).. (-2.05,0.47);
		\path[fill=white] (-0.1,3.4) circle (3.4); 
		\path[ball color=gray!5!white,fill  opacity=0.2] (-0.1,3.4) circle (3.4);
		\path[outer color=gray!15,inner color=white,fill opacity=0.3]  (-0.1,3.4) circle (3.4);
	\end{tikzpicture}}~
			\begin{tikzpicture}[line cap=butt,line join=miter]
					\tikzset{declare function={d = 5; r = 1;}} 
					\draw[fill=gray!35] (0,0)--++(d,0) to [out=60,in=-60] ++(0,2*r) 
						--++(-d,0) to [out=-120,in=120] cycle; 
					\draw (0,0) arc  (270:90:{0.1} and {r})
						(d,0) arc (270:90:{0.1}  and {r});
					\draw[dashed] (0,0)  arc (-90:90:{0.1} and {r})
						(d,0) arc (-90:90:{0.1}  and {r});
				\end{tikzpicture}
	\end{center}
	\loigiai{ 
		Gọi $V_1$ là  thể tích của khối trụ và $V_2$ là thể tích của chỏm cầu.\\
		Nửa chiều cao của khối trụ là: $l=\sqrt{1^2-(0,5)^2}=\dfrac{\sqrt{3}}{2}$ nên ta có thể suy ra chiều cao  của chỏm cầu là: $h=1-\dfrac{\sqrt{3}}{2}$.\\
		Thể tích khối  trụ là: $V_1=\pi R^2.H=\pi.R^2.2l=\dfrac{\pi \sqrt{3}}{4}$.\\
		Thể tích chỏm cầu: $V_2=\pi \displaystyle\int_{R-h}^R\left(\sqrt{R^2-x^2} \right)^2 \mathrm{d}x$ $\Leftrightarrow V_2=\pi \displaystyle\int_{R-h}^R\left(R^2-x^2\right)  \mathrm{d}x\Leftrightarrow V_2=\pi \left. \left(R^2x-\dfrac{x^3}{3} \right)\right|_{R-h}^R$
		\[\Leftrightarrow V_2=\pi \left[\left(R^3-\dfrac{R^3}{3} \right)-\left(R^2(R-h)-\dfrac{(R-h)^3}{3} \right)\right]\Leftrightarrow V_2=\pi  h^2\left(R-\dfrac{h}{3}\right).\]
		Thay số ta suy  ra được thể tích của chỏm cầu là $V_2=\pi \left(\dfrac{2}{3}-\dfrac{3\sqrt{3}}{8} \right)$.\\
		Khi đó thể tích của khối cần tìm  là $V=V_1+2V_2=\dfrac{\pi \sqrt{3}}{4}+\pi \left(\dfrac{2}{3}-\dfrac{3\sqrt{3}}{8} \right)=\dfrac{8-3\sqrt{3}}{6} \pi$ dm$^3$.
	 }
\end{ex}


%%%%%%%%%%%%%%%% ĐÚNG/SAI %%%%%%%%%%%%%%%%

% nguồn: Phan1_Lop_12_Dai_so_Chuong_4
\begin{ex}%[2D4H3-5]%[ma:178d1fc5]
	Giả sử chiếc nón rộng vành sau có thể mô hình hóa bằng cách cho hình phẳng $(H)$ giới hạn bởi đồ thị  hàm số $y=\heva{&x^3+1\quad \quad \text{ khi}\quad 0 < x\le 1\\&
	\sqrt{1-x^2}\quad \text{ khi}\quad-1\le x\le 0.}$, trục  $Ox$ và các đường thẳng $x=-1$ và $x=1$ quay quanh trục $Ox$ (đơn vị trên trục là dm).
	\begin{center}
		\begin{tikzpicture}[scale=1,font=\footnotesize,line join=round,line cap=round,>=stealth]
				 \def\xmin{-1.5}\def\xmax{2}\def\ymin{-2.5}\def\ymax{2.5}
				 \fill[color=gray] (-1,0) arc (180:90:1)--(0,0);
				 \fill[color=gray,samples=200,domain=0:1] (1,0) -- plot (\x,{\x*\x*\x+1})
					(1,0)--(0,0)--(0,1); 
				\draw  (\xmin-0.2,0)--(-1,0);
				 \draw[dashed](-1,0)--(1,0);
				\draw[->] (1,0)--(\xmax+0.2,0) node[below] {\footnotesize $x$}; 
				\draw[->]  (0,\ymin-0.2)--(0,\ymax+0.2) node[right] {\footnotesize $y$};
				\draw (0,0) node  [below left] {\footnotesize $O$};
				\draw (-1,0.1)--(-1,-0.1)  node [below left] {\footnotesize $-1$};
				\draw  (1,0.1)--(1,-0.1) node [below] {\footnotesize $1$};
				 \clip (\xmin,\ymin) rectangle (\xmax,\ymax);
				 \draw (-1,0) arc (180:90:1);
				\draw[smooth,samples=200,domain=0:1.2]  plot (\x,{\x*\x*\x+1});
				 \draw[dashed] (-1,0) arc (180:270:1);
				\draw[dashed,smooth,samples=200,domain=0:1.2]  plot (\x,{(-1)*\x*\x*\x-1});
				\draw [dashed] (0,0) ellipse  (0.15 and 1);
				\draw  [dashed] (1,0) ellipse (0.3 and 2);
				\draw  (1,2)--(1,0);
			\end{tikzpicture}
	\end{center}
	\choiceTF 
		{Diện tích hình phẳng $(H)$ được tính theo công thức $S=\displaystyle\int_{-1}^1\left|\sqrt{1-x^2}+x^3+1\right|\,\mathrm{d}x$} 
		{\True Diện tích thiết diện qua trục đối xứng  của khối tròn xoay trên là $\dfrac{\pi+5}{2}$ dm$^2$}
		{\True Công thức tính thể tích khối  tròn xoay trên là $V=\pi\displaystyle\int_0^{-1}\left(x^2-1\right)\,\mathrm{d}x+\pi\displaystyle\int_0^1\left(x^6+2x^3+1\right)\,\mathrm{d}x$}
		{\True  Nếu thể tích của khối tròn xoay có dạng $\dfrac{a\pi}{b}$ với $\dfrac{a}{b}$ là phân số tối giản thì $a+b=139$}
	\loigiai{ 
		\begin{itemchoice}
			\itemch Diện tích hình  phẳng $(H)$ được tính bằng công thức: $$S=\displaystyle\int_{-1}^0\left(\sqrt{1-x^2}\right)dx+\displaystyle\int_0^1\left(x^3+1\right)dx.$$
			\itemch Ta có $\displaystyle\int_0^1\left(x^3+1\right)dx=\dfrac{5}{4}$.\\ 
			Đặt $x=\sin t$, thì $\mathrm{d}x=\cos t\mathrm{d}t$ . Khi  đó
			\allowdisplaybreaks 
			\begin{eqnarray*}
				&&\displaystyle\int_{-1}^0\left(\sqrt{1-x^2}\right)dx=\displaystyle\int_{-\tfrac{\pi}{2}}^0\left(\sqrt{1-\sin ^2x}\cdot\cos  x\right)dt=\displaystyle\int_{-\tfrac{\pi}{2}}^0\cos^2t\,dt\\
				&=&\displaystyle\int_{-\tfrac{\pi}{2}}^0\dfrac{1+\cos 2t}{2}dt=\left(\dfrac{1}{4}+\dfrac{1}{4}\cdot\sin 2t\right)\Bigg|_{-\tfrac{\pi}{2}}^0=\dfrac{\pi}{4}. 
			\end{eqnarray*}
			Suy ra $S=\dfrac{\pi+5}{4}$.\\ 
			Vậy diện tích thiết diện qua trục đối xứng của khối tròn xoay trên là $2S=\dfrac{\pi+5}{2}$ dm$^2$. 
			\itemch Công thức tính thể tích khối tròn xoay trên là  \[V=\pi\displaystyle\int_0^{-1}{\left(\sqrt{1-x^2}\right)^2\,dx}+\pi\displaystyle\int_0^1\left(x^3+1\right)^2\,dx=\pi\displaystyle\int_0^{-1}{\left(x^2-1\right)dx}+\pi\displaystyle\int_0^1\left(x^6+2x^3+1\right)dx.\]
			\itemch $V=\pi\displaystyle\int_0^{-1}{\left(x^2-1\right)dx}+\pi\displaystyle\int_0^1\left(x^6+2x^3+1\right)dx=\dfrac{97\pi}{42}$  nên $a=97$ và $b=42$ $\Rightarrow a+b=139$.
		\end{itemchoice}
	 }
\end{ex}

% nguồn: Phan1_Lop_12_Dai_so_Chuong_4
\begin{ex}%[2D4V3-5]%[ma:9e82d17b]
	Một cái trứng khủng long đồ chơi là  một khối tròn xoay được tạo thành từ $2$ mảnh ghép lại. Biết mảnh trên được tạo thành khi xoay một phần tư đường elip với trục lớn là $8$ và trục nhỏ là $4$ quanh trục $Ox$ và mảnh dưới được tạo thành khi xoay một phần tư đường tròn bán kính $2$ quanh trục $Ox$ như hình sau (bỏ qua độ dày của vỏ trứng).
	\begin{center}
		\begin{tikzpicture}[scale=1,font=\footnotesize,line join=round,line cap=round,>=stealth]
				 \def\xmin{-4.5}\def\xmax{4.5}\def\ymin{-2.5}\def\ymax{2.5}
				\fill[color=gray] (2,0) arc (0:90:2)--(0,0); 
				\fill[color=gray]  (0,2) arc (90:180:4cm and 2cm)--(0,0);
				\draw[->] (\xmin-0.2,0)--(\xmax+0.2,0) node[below]  {\footnotesize $x$};
				\draw[->]  (0,\ymin-0.2)--(0,\ymax+0.2) node[right] {\footnotesize $y$};
				\draw (0,0) node [below left] {\footnotesize  $O$};
				\draw (-4,0.1)--(-4,-0.1) node [below left] {\footnotesize  $-4$};
				\draw (2,0.1)--(2,-0.1) node [below right] {\footnotesize $2$}; 
				\draw (4,0.1)--(4,-0.1) node [below right]  {\footnotesize $4$};
				\draw (0.1,2)--(-0.1,2) node [above left] {\footnotesize  $2$};
				\draw (0.1,-2)--(-0.1,-2) node [below left] {\footnotesize  $-2$};
				 \clip (\xmin,\ymin) rectangle (\xmax,\ymax);
				 \draw (0,2) arc (90:-90:2);
				\draw[dashed] (4,0) arc  (0:-90:4cm and 2cm);
				\draw[dashed] (4,0) arc (0:90:4cm  and 2cm);
				\draw (0,2) arc (90:270:4cm  and 2cm);
				\draw[dashed]  (0,0) ellipse (0.5 and 2);
			\end{tikzpicture}
	\end{center}
	 \choiceTF
		{\True Thể tích phần trong của mảnh trên  được tính bởi $V_1=\dfrac{\pi}{4}\displaystyle\int_{-4}^0\left(16-x^2\right)dx$}
		{\True  Thể tích phần trong của mảnh trên gấp $2$ lần thể tích phần trong của mảnh dưới}
		{\True Thể tích phần trong của quả  trứng đồ chơi này là $16\pi $}
		{Diện tích thiết  diện khi cắt bởi mặt phẳng qua trục của quả trứng là $3\pi $}
	 \loigiai{
		\begin{itemchoice}
			\itemch Mảnh trên được  biểu diễn bởi elip $\dfrac{x^2}{16}+\dfrac{y^2}{4}=1\Rightarrow{y^2}=4\left(1-\dfrac{x^2}{16}\right)$.\\
			Thể  tích phần trong của mảnh trên được tính bởi $V_1=\pi\displaystyle\int_{-4}^0\left[4\left(1-\dfrac{x^2}{16}\right)\right]dx=\dfrac{\pi}{4}\displaystyle\int_{-4}^0\left(16-x^2\right)dx$.
			\itemch  Thể tích phần trong mảnh trên là $V_1=\dfrac{\pi}{4}\displaystyle\int_{-4}^0\left(16-x^2\right)dx=\dfrac{32\pi}{3}$.\\
			Mảnh dưới  được biểu diễn bởi đường tròn $x^2+y^2=4\Rightarrow{y^2}=4-x^2$.\\
			Thể tích phần  trong dưới là $V_2=\pi\displaystyle\int_0^2\left(4-x^2\right)dx=\dfrac{16\pi}{3}$.\\
			Vậy thể tích phần trong của mảnh trên gấp $2$ lần  thể tích phần trong của mảnh dưới.
			\itemch Thể tích phần trong của quả trứng đồ chơi này  là $V=V_1+V_2=\dfrac{32\pi}{3}+\dfrac{16\pi}{3}=16\pi$.
			 \itemch Thiết diện khi cắt bởi mặt phẳng qua trục của quả trứng là một hình tròn có bán kính bằng $2$ nên có diện tích là $4\pi$.
	\end{itemchoice}}
\end{ex}

% nguồn: Phan1_Lop_12_Dai_so_Chuong_4
\begin{ex}%[2D4V3-5]%[ma:48943433]
	Một ly trà sữa dạng hình nón cụt, có đường kính đáy ly $6$ cm, đường kính miệng ly $9$ cm, chiều cao $13{,}4$ cm, ở miệng ly có sử dụng một nắp đậy có hình dạng nửa mặt cầu và ở đỉnh của nửa mặt cầu này có một hình tròn có đường kính $2$ cm để cắm ống hút, mặt phẳng chứa hình tròn này song song với mặt phẳng  chứa miệng ly (tham khảo hình vẽ sau).
	\begin{center}
		{\begin{tikzpicture}[scale=0.5]
		\def\Rtop{4.5} 
		 \def\Rbottom{3}
		\def\Rhole{1} 
		\def\hcone{9.9} 
		 \def\Rcap{4.5}
		\draw (-\Rbottom,0) arc (180:360:\Rbottom cm and 0.6cm); 
		\draw[dashed] (-\Rbottom,0) arc (180:-360:\Rbottom cm and 0.6cm); 
		\draw (-\Rbottom,0) --  (-\Rtop,\hcone);
		\draw (\Rbottom,0) -- (\Rtop,\hcone); 
		\draw (-\Rtop,\hcone) arc  (180:360:\Rtop cm and 0.9cm);
		\draw [dashed](-\Rtop,\hcone) arc (180:0:\Rtop cm and  0.9cm);
		 \draw (-\Rtop,\hcone) arc (180:104:\Rtop);
		\draw (\Rtop,\hcone) arc  (0:78:\Rtop);
		\path (-1,\hcone+\Rcap-0.1)  coordinate (top);
		\draw ([xshift=-\Rhole]top) arc (180:360:\Rhole cm and  0.25cm);
		 \path (-1,\hcone+\Rcap-0.1) coordinate (top);
		\draw ([xshift=-\Rhole]top) arc  (180:-360:\Rhole cm and 0.25cm);
		\draw[<->] (\Rtop+0.3,0)  -- (\Rtop+0.3,\hcone) node[midway,right] {13,4 cm};
		\draw[<->] (-\Rbottom,-0.7) -- (\Rbottom,-0.7) node[midway,below] {6  cm};
		\draw[<->] (-\Rtop,\hcone) --  (\Rtop,\hcone) node[midway,above] {9 cm};
		\draw[<->] (-\Rhole,\hcone+\Rcap+0.5) -- (\Rhole,\hcone+\Rcap+0.5)  node[midway,above] {2 cm};
	\end{tikzpicture}
		}\hspace{0.5cm}{\begin{tikzpicture}[scale=0.5,rotate=-90]
		\draw[->] (5,9.9) -- (-5.5,9.9) node[left]  {\(y\)};
		\draw[->] (0,-1) -- (0,16)  node[below right] {\(x\)};
		 \draw[dashed](0,0) node[below]{$A$}--(-3,0) node[above]{$B$}
			(0,9.9) node[below left]{$O$} 
			 (-4.5,9.9) node[above right]{$C$}
			(-1,14.4)  node[right]{$D$}--(0,14.4);
		 \def\Rtop{4.5}
		 \def\Rbottom{3}
		\def\Rhole{1} 
		 \def\hcone{9.9}
		 \def\Rcap{4.5}
		\draw (-\Rbottom,0)  arc (180:360:\Rbottom cm and 0.6cm);
		\draw[dashed] (-\Rbottom,0) arc  (180:-360:\Rbottom cm and 0.6cm);
		 \draw (-\Rbottom,0) -- (-\Rtop,\hcone);
		 \draw (\Rbottom,0) -- (\Rtop,\hcone);
		\draw (-\Rtop,\hcone) arc (180:360:\Rtop cm  and 0.9cm);
		\draw [dashed](-\Rtop,\hcone) arc (180:0:\Rtop cm and 0.9cm); 
		\draw (-\Rtop,\hcone) arc  (180:104:\Rtop);
		\draw  (\Rtop,\hcone) arc (0:78:\Rtop);
		\path (-1,\hcone+\Rcap-0.1)  coordinate (top);
		\draw ([xshift=-\Rhole]top) arc (180:360:\Rhole cm  and 0.25cm);
		\path (-1,\hcone+\Rcap-0.1)  coordinate (top);
		\draw ([xshift=-\Rhole]top) arc (180:-360:\Rhole  cm and 0.25cm);
	\end{tikzpicture}
		 }
	\end{center}
	Chọn hệ trục $Oxy$ (đơn vị trên trục là centimet) với trục $Ox$ đi qua tâm của $2$ đáy hình nón cụt và gốc tọa độ $O$ trùng với tâm của đáy lớn như hình vẽ trên. 
	\choiceTF 
		{\True Phương trình  đường thẳng $BC$ là $1{,}5x-13,4y+60,3=0$}
		{\True Tọa độ  điểm $D$ là $D\left(\dfrac{\sqrt{77}}{2};1\right)$}
		{\True Thể tích bên trong của ly không bao gồm nắp  là $600$ ml (làm tròn kết quả đến hàng đơn vị)}
		{Thể tích bên trong của ly bao gồm cả thể tích  của nắp là $780$ ml (làm tròn kết quả đến hàng đơn vị)}
	 \loigiai{
		\begin{itemchoice}
			 \itemch Ta có $B=\left(-13{,}4; 3\right)$ và $C=\left(0; 4{,}5\right)$ nên $\overrightarrow{BC}=\left(13{,}4; 1{,}5\right)$ nên đường thẳng $BC$ có một vectơ pháp tuyến là $\overrightarrow{n}_{BC}=\left(-1{,}5; 13{,}4\right)$.\\
			Phương trình đường  thẳng $BC$ là $-1{,}5\left(x+13{,}4\right)+13{,}4\left(y-3\right)=0$ $\Leftrightarrow 1{,}5x-13{,}4y+60{,}3=0$.
			\itemch Hoành độ của  điểm $D$ là $x_D=\sqrt{\left(4{,}5\right)^2-1}=\dfrac{\sqrt{77}}{2}\Rightarrow D\left(\dfrac{\sqrt{77}}{2}; 1\right)$.
			 \itemch Thể tích bên trong của ly không bao gồm nắp là \[V_1=\pi{\displaystyle\int_{-13{,}4}^0\left(\dfrac{1{,}5x+60{,}3}{13{,}4}\right)^2}dx\approx 600 \,\text{ cm}^3=600 \,\text{ ml}.\]
			\itemch Thể tích bên trong của ly bao gồm cả thể tích của nắp là: $V=V_1+\left(V_2-V_3\right)$  trong đó $V_2$ là nửa thể tích của khối cầu và $V_3$ là thể tích của chỏm cầu.\\
			Nửa thể tích khối cầu là:  $V_2=\dfrac{1}{2}\cdot\dfrac{4}{3}\cdot\pi\cdot R^3=\dfrac{243\pi}{4}$.\\
			 Thể tích chỏm cầu
			 \allowdisplaybreaks
			\begin{eqnarray*}
				 V_3&=&\pi\displaystyle\int_{R-h}^R\left(\sqrt{R^2-x^2}\right)^2\,\mathrm{d}x=\pi\displaystyle\int_{R-h}^R\left(R^2-x^2\right)\,\mathrm{d}x\\
				 &=&\pi\left(R^2x-\dfrac{x^3}{3}\right)\Big|_{R-h}^R=\pi\left[\left(R^3-\dfrac{R^3}{3}\right)-\left(R^2\left(R-h\right)-\dfrac{\left(R-h\right)^3}{3}\right)\right]=\pi{h^2}\left(R-\dfrac{h}{3}\right).
			\end{eqnarray*}
			Thay số ta có $h=4{,}5-\dfrac{\sqrt{77}}{2}$; $R=4{,}5$  $\Rightarrow{V_3}=\dfrac{\left(729-83\sqrt{77}\right)\pi}{12}$.\\
			Thể tích  bên trong của ly bao gồm cả thể tích của nắp là
			\[V=600+\left[\dfrac{243\pi}{4}-\dfrac{\left(729-83\sqrt{77}\right)\pi}{12}\right]\approx  790\, \text{ ml}.\]
		\end{itemchoice}
	} 
\end{ex}

% nguồn: Phan1_Lop_12_Dai_so_Chuong_4
\begin{ex}%[2D4V3-5]%[ma:a69bb52f]
	Người ta chế tác một giọt nước bằng  thủy tinh. Biết giọt nước thủy tinh này là vật thể tròn xoay khi xoay hình phẳng giới hạn bởi đồ thị của hàm số \[f(x)=\heva{&\sqrt{4-x^2}\quad\quad\quad\quad\quad\left(-2\le x\le 0{,}6\right)\\&
	-\dfrac{\sqrt{91}}{20}x+\dfrac{23\sqrt{91}}{100}\quad\left(0{,}6 < x\le 4{,}6\right).}\] và trục $Ox$ quanh trục $Ox$ (đơn vị trên trục là centimet). Xét tính đúng  sai của các khẳng định sau:
	\begin{center}
		\begin{tikzpicture}[scale=1,font=\footnotesize,line join=round,line cap=round,>=stealth]
				\def\xmin{-2.5}\def\xmax{5}\def\ymin{-2.5}\def\ymax{2.5} 
				\fill[color=gray]  (-2,0) arc (180:72.4:2)--(0.6,0);
				\fill[color=gray,samples=200,domain=0.6:4.6] (0.6,0) -- plot (\x,{-0.476*\x+2.19}); 
				 \draw[->](\xmin-0.2,0)--(\xmax+0.2,0) node[below] {\footnotesize $x$};
				\draw[->] (0,\ymin-0.2)--(0,\ymax+0.2) node[right] {\footnotesize  $y$};
				\draw (0,0) node [below left]  {\footnotesize $O$};
				 \draw[dashed] (0.6,0) node [below] {\footnotesize $0.6$}--(0.6,1.907);
				\clip (\xmin,\ymin) rectangle (\xmax,\ymax); 
				\draw (-2,0) arc (180:72.4:2); 
				\draw[smooth,samples=200,domain=0.6:4.6]  plot (\x,{-0.476*\x+2.19});
				\draw[dashed] (-2,0) arc  (180:287.5:2);
				\draw[dashed,smooth,samples=200,domain=0.6:4.6] plot (\x,{+0.476*\x-2.19}); 
				 \draw (-2,0) node[below left]{$-2$}
					(0,2) node[above  left]{$2$};
			\end{tikzpicture}
	\end{center}
	\choiceTF 
		{\True  Hàm số $y=f(x)$ liên tục tại $x=0{,}6$}
		{\True Diện tích mặt cắt của giọt nước thủy tinh khi cắt bởi mặt phẳng qua trục  được tính bởi công thức $S=2\displaystyle\int_{-2}^{4,6}f(x)\,\mathrm{d}x$ cm$^2$}
		{Thể tích của giọt nước thủy tinh này lớn hơn $40$  cm$^3$}
		{\True Biết khối lượng riêng của thủy tinh là $\rho=2{,}6$ g/cm$^3$, khối lượng của giọt nước thủy tinh  này là $102{,}22$ g (làm tròn kết quả đến hàng phần trăm)}
	\loigiai{ 
		\begin{itemchoice}
			\itemch  Ta có $\lim\limits_{x\to 0{,}6^-}f(x)=\sqrt{4-\left(0{,}6\right)^2}=\sqrt{4-0{,}36}=\sqrt{3{,}64}\approx 1{,}907$.\\ và $\lim\limits_{x\to 0{,}6^+}f(x)=-\dfrac{\sqrt{91}}{20}\cdot0{,}6+\dfrac{23\sqrt{91}}{100}\approx 1{,}907$.\\
			 Do đó, $\lim\limits_{x\to 0{,}6^-}f(x)= \lim\limits_{x\to 0{,}6^+}f(x)=f(0{,}6)$.\\
			Vậy hàm số $y=f(x)$ liên tục tại  $x=0{,}6$.
			 \itemch Diện tích mặt cắt của giọt nước thủy tinh khi cắt bởi mặt phẳng qua trục được tính bởi công thức $S=2\displaystyle\int_{-2}^{4,6}f(x)\,\mathrm{d}x$ cm$^2$.
			 \itemch Thể tích của giọt nước thủy tinh này là\\
			 $V=V_1+V_2=\pi\displaystyle\int_{-2}^{0{,}6}{\left(\sqrt{4-x^2}\right)^2\,\mathrm{d}x+}\,\pi\displaystyle\int_{0{,}6}^{4{,}6}{\left(-\dfrac{\sqrt{91}}{20}x+\dfrac{23\sqrt{91}}{100}\right)^2\,\mathrm{d}x}=\dfrac{4\,693\pi}{375}\approx 39{,}32$ cm$^3$.
			\itemch Khối lượng của giọt  nước thủy tinh này là $m=\rho \cdot V=2{,}6\cdot\dfrac{4\,693\pi}{375}\approx 102{,}22$ g.
		\end{itemchoice}
	} 
\end{ex}

% nguồn: Phan1_Lop_12_Dai_so_Chuong_4
\begin{ex}%[2D4V3-5]%[ma:d895107c]
	Một cái bể nước có dạng khối chóp cụt tứ giác đều  ngược với cạnh đáy bằng $3\sqrt{2}$ dm và chiều cao bằng $6$ dm (tham khảo hình vẽ bên – các kích thước được nêu ra là phần bên trong hình). Nước được bơm vào bể với tốc độ không đổi là $2$ lít/phút và ban đầu bể không chứa nước (các kết quả bên dưới được làm tròn đến hai chữ số thập phân sau dấu phẩy).
	\begin{center}
		\begin{tikzpicture}[scale=0.8,font=\footnotesize,line join=round,line cap=round,>=stealth]
				\path 
					 (0,1.5) coordinate (A)
					++(-140:1.5)  coordinate (B)
					++(0:2.5)  coordinate (C)
					($(A)+(C)-(B)$) coordinate  (D)
					($(A)!1/2!(C)$) coordinate (O) 
					 ($(O)+(0,-3.5)$) coordinate (S);
				\path  ($(S)!0.7!(A)$) coordinate (M);
				\path ($(S)!0.7!(B)$) coordinate (N); 
				\path ($(S)!0.7!(C)$) coordinate (P); 
				\path  ($(S)!0.7!(D)$) coordinate (Q);
				 \fill[cyan!30,opacity=0.7] (M)--(N)--(P)--(Q)--cycle;
				 \fill[cyan!30,opacity=0.5] (N)--(P)--(S)--cycle;
				\fill[cyan!30,opacity=0.5]  (P)--(Q)--(S)--cycle;
				\draw[teal]  (A)--(B)--(C)--(D)--cycle;
				\draw[teal,dashed] (S)--(A) 
					 (S)--(C)
					 (Q)--(M)--(N);
				\draw[teal] (S)--(D) 
					(S)--(C) 
					(S)--(B) 
					(N)--(P)--(Q); 
				\draw[teal]  (M)--(S)
					 (N)--(S)
					(P)--(S) 
					(Q)--(S); 
				 \draw[teal,dashed] (A)--(S);
			\end{tikzpicture}
	\end{center}
	\choiceTF 
		{\True Bể  nước được bơm đầy sau $18$ phút}
		{\True Tốc độ dâng lên của nước là $0{,}23$ dm/phút khi thể tích nước bằng  $\dfrac{1}{3}$ thể tích bể}
		{\True Khi mực nước cách miệng bể $0{,}5$ dm, người ta ngừng bơm và bắt đầu xả ra với lưu lượng tốc độ giảm  chiều cao của mực nước trong bể theo thời gian $t$ (phút) được mô hình hóa bởi hàm số $h'(t)=\dfrac{1}{350}-\dfrac{193}{700}t$ (dm/phút). Sau $5$ phút, thể tích nước trong bể là $11{,}97$ dm$^3$}
		{Cùng điều kiện ở c) thì sau  $23{,}59$ phút nước trong bể vừa được xả hết}
	\loigiai{ 
		\begin{itemchoice}
			\itemch 
			Thể tích chậu  nước hình chóp tứ giác đều là
			 \[
			 V_{\text{ chậu}}=\dfrac{1}{3}\cdot\left(3\sqrt{2}\right)^2\cdot 6=36\ \mathrm{dm}^3=36 \, \text{ (lít).}
			\] 
			Thời gian  bơm nước đầy bể là $\dfrac{36}{2}=18$ phút.
			 \itemch
			Gọi $V(t)$, $h(t)$ lần  lượt là thể tích và chiều cao của nước sau $t$ phút.\\
			Ta có 
			\[ 
			 \dfrac{V(t)}{V_{\text{ chậu}}}=\left(\dfrac{h(t)}{6}\right)^3\quad\Rightarrow\quad V(t)=\left(\dfrac{h(t)}{6}\right)^3\cdot 36
			\] 
			hay 
			\[ 
			V(t)=\dfrac{\left( h(t)\right) ^3}{6}.  \tag{*}
			 \]
			Đạo hàm hai vế của  (*) theo biến $t$ ta được
			 \[
			 \dfrac{\mathrm{d}V(t)}{\mathrm{d}t}=\dfrac{(h(t))^2}{2}\cdot \dfrac{\mathrm{d}h(t)}{\mathrm{d}t}. \tag{**}
			\] 
			Thời điểm thể tích nước bằng $\dfrac{1}{3}$  thể tích chậu thì
			\[ 
			 \dfrac{(h(t))^3}{6}=\dfrac{36}{3}\Rightarrow (h(t))^3=72\Rightarrow h(t)\approx 4{,}16\ \mathrm{dm}.
			\] 
			Thay $\dfrac{\mathrm{d}V(t)}{\mathrm{d}t}=2$  dm/phút và $h(t)\approx 4{,}16$ vào (**) thì
			\[ 
			2=\dfrac{(4{,}16)^2}{2}\cdot  \dfrac{\mathrm{d}h}{\mathrm{d}t}\Rightarrow \dfrac{\mathrm{d}h(t)}{\mathrm{d}t}\approx 0{,}23 \,\text{ dm/phút}.
			 \]
			\itemch 
			Mực nước cách miệng bể $0{,}5$  dm nên chiều cao ban đầu bằng $5{,}5$ dm.\\
			Chiều  cao của nước trong chậu sau $5$ phút là
			 \[
			h(5)=5{,}5+\displaystyle\int_0^5\left(\dfrac{1}{350}t-\dfrac{193}{700}\right) \mathrm{\,d}t=\dfrac{291}{70}\approx 4{,}16\ \mathrm{dm}. 
			 \]
			Thể tích nước còn lại trong  bể là
			\[ 
			V(5)=\left(\dfrac{h(5)}{6}\right)^3\cdot  36\approx 11{,}97\ \mathrm{dm}^3.
			 \]
			\itemch 
			Thời gian để  mực nước trong chậu đang là $5{,}5$ dm trở về $0$ dm thoả mãn
			\[ 
			5{,}5+\displaystyle\int_0^t\left(\dfrac{1}{350}t-\dfrac{193}{700}\right)\mathrm{\,d}t=0\Rightarrow t\approx 22{,}59\  \text{ phút}.
			 \]
			Với khoảng thời gian $22{,}59$ phút thì nước trong  bể vừa được xả hết.
		\end{itemchoice}
	} 
\end{ex}


%%%%%%%%%%%%%%%% TRẢ LỜI NGẮN %%%%%%%%%%%%%%%%

% nguồn: Phan2_Lop_12_Dai_so_Chuong_4
\begin{bt}%[2D4N3-5]%[ma:a20057ac]
	Cho $(H)$ là hình phẳng giới hạn  bởi đồ thị hàm số $y=\sqrt{x}$, trục hoành và hai đường thẳng $x=1$, $x=3$. Tính thể tích khối tròn xoay tạo thành khi quay $(H)$ xung quanh trục hoành (làm tròn kết quả đến hàng phần chục).
	\shortans{12,6} 
	\loigiai{ 
		Thể tích khối tròn xoay là $V=\pi\displaystyle \int_{1}^{3}\left(\sqrt{x} \right)^2\mathrm{\,d}x\approx12{,}6$. 
	 }
\end{bt}

% nguồn: Phan2_Lop_12_Dai_so_Chuong_4
\begin{bt}%[2D4H3-5]%[ma:55ecb6c5]
	Xét phần bên trong của một thùng rượu: là một khối  tròn xoay có $2$ đáy là hình tròn bằng nhau và có chiều cao là $60$ cm, đường cong (bên trong) của thùng là một cung tròn của đường tròn bán kính là $36$ cm. Thùng rượu này chứa được tối đa bao nhiêu lít rượu? (kết quả làm tròn đến hàng đơn vị)
	\begin{center}
		\begin{tikzpicture}[scale=0.4,font=\footnotesize,line join=round,line cap=round,>=stealth]
				\begin{scope}[rotate=90]
					\fill[orange!70!black] (-5,-3) arc (-90:-270: 0.4 and 3)--(-5,3)--plot[domain=-5:5](\x,{-1/25*(\x)^2+4})--(5,3) arc (90:-90: 0.4 and  3)--(5,-3)--plot[domain=5:-5](\x,{1/25*(\x)^2-4})--cycle;
					\fill[orange!50!black] (5,0) ellipse  (0.4 and 3);
					\fill[orange!80!black]  (5,0) ellipse (0.3 and 2.8);
					\draw[color=orange!50!black,line width=3pt]  (4.9,3) arc (90:270: 0.4 and 3);
					 \draw[color=orange!50!black,line width=3pt] (4,3.3) arc (90:270: 0.4 and 3.3);
					\draw[color=orange!50!black,line width=3pt] (2.5,3.7) arc (90:270: 0.4 and 3.7); 
					\draw[color=orange!50!black,line width=3pt] (-2.5,3.7) arc (90:270:  0.4 and 3.7);
					\draw[color=orange!50!black,line width=3pt] (-4,3.3) arc (90:270: 0.4  and 3.3);
					\draw[color=orange!50!black,line width=3pt] (-4.9,3) arc (90:270: 0.4 and  3);
				\end{scope}
			\end{tikzpicture}
	\end{center}
	 \shortans{188}
	 \loigiai{
		Phương trình đường  tròn là $x^2+y^2=36^2\Rightarrow{y^2}=36^2-x^2=1296-x^2$.\\
		Thể tích  của thùng rượu là $V=\pi\displaystyle\int_{-30}^{30}y^2\,\mathrm{d}x=\pi\displaystyle\int_{-30}^{30}{\left(1296-x^2\right)\,\mathrm{d}x}=59760\pi$ cm$^3$$\approx187{,}7$ dm$^3$.\\
		Vậy  thùng rượu chứa được tối ta $187{,}7$ lít rượu.
	 }
\end{bt}

% nguồn: Phan2_Lop_12_Dai_so_Chuong_4
\begin{bt}%[2D4H3-5]%[ma:85bbcee1]
	Một cái chậu cao $16$ cm. Khi đổ nước vào chậu, nếu độ cao của nước là $x$ (cm) $(0 \leq x \leq 16)$ thì mặt thoáng của nước là hình tròn bán  kính $r=9+\sqrt{x}$
	 (cm). Tính dung tích của chậu (kết quả làm tròn đến hàng phần trăm của lít).
	 \shortans{6,89}
	 \loigiai{
		Ta có \begin{eqnarray*}
		V=& & \pi \displaystyle\int_{0}^{16}(9+\sqrt{x})^2 \mathrm{\,d}x\\ 
		=&&\pi \displaystyle\int_0^{16}(81+18 \sqrt{x}+x)  \mathrm{\,d}x=\left.\pi\left(81 x+12x\sqrt{x}+\dfrac{x^2}{2}\right)\right|_0^{16}=2192 \pi \mathrm{~cm}^3.\\
	\end{eqnarray*}
	Suy  ra $V=2{,}192 \pi$ $\ell$ $\approx 6{,}89$ $\ell$.
 }
\end{bt}

% nguồn: Phan2_Lop_12_Dai_so_Chuong_4
\begin{bt}%[2D4V3-5]%[ma:5eab5713]
	Một thùng chứa rượu có dạng hình tròn xoay với đường sinh là một cung của đường tròn, hai đáy là hai hình tròn bằng nhau và có đường kính bằng $20$ cm, khoảng cách giữa hai đáy bằng $32$ cm, đường kính lớn nhất của thân thùng là $24$ cm (biết các kích thước đã trừ lớp vỏ thùng, tham khảo hình vẽ). Hỏi phần không gian bên trong thùng chứa rượu đỏ có thể tích bằng bao nhiêu lít (kết quả lấy đến chữ số hàng phần chục và không  làm tròn)?
	\begin{center}
		\begin{tikzpicture}
				\begin{scope}[scale=.65,line join=round,line cap=round,>=stealth,line width=1pt,y=0.9cm,rotate=90]
					\draw[left color=brown!60!black,right color=brown!60!black,middle color=brown!30!orange,opacity=0.5] plot[variable=\t,domain=-0.7:0.7,smooth] ({3*cos (\t r)},{5*sin  (\t r)}) arc (0:-180:2.29 cm and 0.5cm)--plot[variable=\t,domain=0.7:-0.7,smooth] ({-3*cos (\t r)},{5*sin (\t r)}) arc (180:360:2.29 cm and 0.5cm)--cycle;
					\draw[left color=brown!60!black,right color=brown!60!black,middle color=brown!30!orange,opacity=0.5] plot[variable=\t,domain=-0.7:0.7,smooth] ({3*cos (\t r)},{5*sin (\t r)}) arc (0:-180:2.29 cm and 0.5cm)--plot[variable=\t,domain=0.7:-0.7,smooth] ({-3*cos (\t r)},{5*sin (\t r)}) arc (180:360:2.29  cm and 0.5cm)--cycle;
					\shade[left  color=brown!40!black,right color=brown!40!black,middle color=brown!50!yellow,opacity=0.5] (0,3.22) ellipse (2.29 and 0.5);
					 \draw (0,3) ellipse (2.29 and 0.5);
					\draw (0,3.22) ellipse (2.29 and 0.5); 
					\draw[left color=brown!60!black,right  color=brown!60!black,middle color=brown!30!orange](-2.29,-3.22) arc (180:360:2.29 cm and 0.5 cm)--(2.48,-2.82) arc (360:180:2.48 cm and 0.5 cm)--cycle;
					\draw[left  color=brown!60!black,right color=brown!60!black,middle color=brown!50!orange](-2.9,-1.29) arc (180:360:2.9 cm and 0.5 cm)--(2.96,-0.83) arc (360:180:2.96cm and 0.5 cm)--cycle;
					\draw[left color=brown!60!black,right color=brown!60!black,middle color=brown!50!orange](-2.94,0.99) arc (180:360:2.94 cm and 0.5  cm)--(2.89,1.35) arc (360:180:2.89cm and 0.5 cm)--cycle;
					\draw[left color=brown!60!black,right color=brown!60!black,middle color=brown!50!orange](-2.45,2.9) arc (180:360:2.45 cm and 0.5 cm)--(2.29,3.22) arc (360:180:2.29cm and  0.5 cm)--cycle;
				\end{scope}
				\begin{scope}[xshift=6cm,scale=0.7,font=\footnotesize,line join=round,line cap=round,>=stealth,yscale=0.5]
					\tikzset{declare function={tl=0.3;  xmin=-4; xmax=4; ymin=-5; ymax=5;
						f  (\x)=sqrt (25-(\x)^2);}}
					 \draw[smooth,samples=100,domain=-3:3] plot (\x,{f (\x)})
						 plot ({\x,-f (\x)});
					\draw[dashed,<->] ({-f (0)+2},0) -- ({f (0)-2},0); \draw (2,0) node[below]{$32$\,cm}; 
					 \draw[<->] (-3,{-f (-3)}) -- (-3,{f (-3)}) node[sloped,midway,above]{$20$\,cm};
					\draw[<->] (3,{-f (3)}) -- (3,{f  (3)});
					\draw[<->] (0,{-f (0)}) --  (0,{f (0)}) node[sloped,midway,above]{$24$\,cm};
				\end{scope}
			\end{tikzpicture}
	\end{center}
	\shortans{12,9} 
	\loigiai{ 
		\begin{center}
			\begin{tikzpicture}[scale=1,font=\footnotesize,line join=round,line cap=round,>=stealth,yscale=0.5]
					\tikzset{declare function={tl=0.3; xmin=-4; xmax=4; ymin=-5; ymax=5; 
						f  (\x)=sqrt (25-(\x)^2);}}
					\draw[smooth,samples=100,domain=-4:4]  plot (\x,{f (\x)})
						 plot ({\x,-f (\x)});
					\draw[->]  (-5,0) -- (5,0) node[above]{$x$};
					 \draw[->] (0,-6) -- (0,6) node[right]{$y$};
					\draw (-4,{-f (-4)}) --  (-4,{f (-4)});
					 \draw (4,{-f (4)}) -- (4,{f (4)});
					\path 
						(0,5)  coordinate (A)
						(0,0) coordinate (O) 
						(4,0)  coordinate (D)
						(-4,{f (-4)}) coordinate (M) 
						(4,{f  (4)}) coordinate (N);
					 \foreach \i/\g in {A/45,M/90,N/0,O/45}
					\fill (\i) circle (1pt)+(\g:4mm) node[scale=1]{$\i$}; 
					\draw (-4,0)  node[above right] {$-16$};
					\draw (4,0)  node[above left] {$16$};
					\draw (0,5) node[above left] {$12$}; 
					\draw (0,{f (4)}) node[above left]  {$10$};
					\draw[dashed] (M)--(N); 
					 \fill (-4,0) circle (1pt);
					 \fill (4,0) circle (1pt);
					\fill  (0,{f (4)}) circle (1pt);
				\end{tikzpicture}
		\end{center}
		 Chọn hệ toạ độ $Oxy$ như hình vẽ.\\
		Gọi  $(C)$ là đường tròn đi qua $3$ điểm $A(0;12)$, $M(-16;10)$, $N(16;10)$.\\
		Giả sử phương trình của $(C)$ có dạng  $x^2+y^2-2ax-2by+c=0.$\\
		Vì $(C)$ đi qua $3$ điểm $A(0;12)$,  $M(-16;10)$, $N(16;10)$ nên ta có
		\allowdisplaybreaks 
		\begin{eqnarray*}
			 \heva{&0^2+12^2-2a\cdot0-2b\cdot12+c=0\\&(-16)^2+10^2-2a\cdot(-16)-2b\cdot10+c=0\\&16^2+10^2-2a\cdot16-2b\cdot10+c=0}\Leftrightarrow\heva{&24b-c=144\\&32a-20b+c=-356\\&32a+20b-c=356}\Leftrightarrow\heva{&a=0\\&b=-53\\&c=-1\,416.}
		\end{eqnarray*}
		Suy  ra phương trình đường tròn $(C)$ là
		 \[x^2+y^2+106y-1\,416=0\Leftrightarrow x^2+(y+53)^2=4\,225.
		\] 
		 Ta có
		\[x^2+(y+53)^2=4\,225\Leftrightarrow\hoac{&y=-53+\sqrt{4\,225-x^2}\\&y=-53-\sqrt{4\,225-x^2}.} 
		 \]
		Gọi $(C_1)$ là cung của đường tròn $(C)$ đi qua $3$  điểm $A$, $M$, $N$.\\
		Ta có $(C_1)$  là đồ thị hàm số $y=-53+\sqrt{4\,225-x^2}=f(x)$.\\
		Gọi $(H)$  là hình phẳng giới hạn bởi các đồ thị $y=f(x)$, $Ox$, $x=-16$, $x=16$.\\
		Quay hình phẳng $(H)$ quanh trục $Ox$ ta được vật  thể tròn xoay có thể tích là
		\begin{align*}
			V=\pi\displaystyle\int_{-16}^{16}\left(-53+\sqrt{4\,225-x^2}\right)^2\mathrm{d}x\approx12\,957{,}7\text{ \,cm}^3\approx12{,}9577\text{ \,dm}^3. 
		\end{align*}
		Vậy thể tích của thùng chứa rượu xấp xỉ $12{,}9$ lít. 
	} 
\end{bt}

% nguồn: Phan2_Lop_12_Dai_so_Chuong_4
\begin{bt}%[2D4V3-5]%[ma:51added2]
	Một chi  tiết máy được thiết kế bằng cách quay hình phẳng giới hạn bởi đồ thị hàm số $y=x^2-4x+5$, trục hoành và các đường thẳng $x=1$, $x=4$ quanh trục hoành. Biết đơn vị trên các trục tọa độ là cen-ti-mét. Thể tích của chi tiết máy đó (kết quả làm tròn đến hàng đơn vị) bằng bao nhiêu cen-ti-mét khối?
	 \shortans{49}
	 \loigiai{
	$V = \pi \displaystyle \int_{1}^{4} \left(x^2 - 4x + 5\right)^2 \mathrm{\,d}x  = \dfrac{78\pi}{5} \approx 49 \left(\text{cm}^3\right).$
	 }
\end{bt}

% nguồn: Phan2_Lop_12_Dai_so_Chuong_4
\begin{bt}%[2D4V3-5]%[ma:d2af4811]
	Một quả  bóng bầu dục theo quy định được sử dụng trong giải bầu dục quốc gia có kích thước $28$ (cm) từ đầu này đến đầu kia và đường kính $17$ (cm) ở phần dày nhất (quy định có phép thay đổi một chút về các kích thước này) (Nguồn: NFL).
	Hình dạng của một quả bóng bầu dục có kích thước nói trên có thể được tạo  thành khi quay phần diện tích hình phẳng giới hạn bởi đồ thị hàm số $f(x)=ax^2+bx+c$, trục hoành và các đường thẳng $x=-4$; $x=24$, trong đó $x$ tính bằng centimet (cm). Thể tích (đơn~vị: cm$^3$, kết quả làm tròn đến hàng đơn~vị) của quả bóng bầu dục có kích thước nói trên bằng bao nhiêu?\shortans{3390}
	\loigiai{\begin{center}
		\begin{tikzpicture}[scale=.3,>=stealth,font=\footnotesize,line join=round,line cap=round]
				 \def\a{-17/392} \def\b{85/98} \def\c{204/49}
				\def\xmin{-4.5}  \def\xmax{25}
				 \def\ymin{-1.5} \def\ymax{9.5}
				\draw[color=gray!50,dashed] (\xmin,\ymin) grid (\xmax,\ymax); 
				\draw[->] (\xmin,0)--(\xmax,0) node [below]{$x$}; 
				\draw[->] (0,\ymin)--(0,\ymax) node  [left]{$y$};
				\node  at (0,0) [below left]{$O$};
				\clip (\xmin+0.1,\ymin+0.1)  rectangle (\xmax-0.5,\ymax-0.1);
				\draw[smooth,samples=300,domain=-4.2:24.2]  plot (\x,{\a*(\x)^2+\b*(\x)+\c});
				\draw[dashed](10,0)  node[below]{$10$}|-(0,17/2) node[left]{$\tfrac{17}{2}$};
				\draw (24,0) node[below]{$24$}circle  (1pt)
					(-4,0) node[below]{$-4$}circle  (1pt);
			\end{tikzpicture}
	\end{center}
	Ta thấy đồ thị hàm số $f(x)=ax^2+bx+c$ đi  qua các điểm $(-4;0)$, $(24;0)$ và $\left(10;\dfrac{17}{2}\right)$ nên ta có hệ phương trình\\
	\[\left\{\begin{aligned}
	&(-4)^2a+(-4)b+c=0\\  &24^2a+24b+c=0\\ &10^2a+10b+c=\dfrac{17}{2}
	\end{aligned}
	\Leftrightarrow \left\{\begin{aligned}
	&a=-\dfrac{17}{392} \\ &b=\dfrac{85}{98}  \\ &c=\dfrac{204}{49}.
	\end{aligned}
	\right.\right.\\ 
	 \]
	Vậy  thể tích của quả bóng bầu dục trên là\\
	\[V=\pi\displaystyle\int^{24}_{-4}{\left(-\dfrac{17}{392}x^2+\dfrac{85}{98}x+\dfrac{204}{49}\right)}^2  \,\mathrm{d}x=\dfrac{16284}{15}\pi\approx 3390\left(\text{cm}^3\right).
	\] 
	 }
\end{bt}

% nguồn: Phan2_Lop_12_Dai_so_Chuong_4
\begin{bt}%[2D4V3-5]%[ma:05e74d6f]
	Một ly trà sữa dạng hình nón cụt, có đường kính đáy ly $6$ cm, đường kính miệng ly $9$~cm, chiều cao $13{,}4$ cm, ở miệng ly có sử dụng một nắp đậy có hình dạng nửa mặt cầu và ở đỉnh của nửa mặt cầu này có một hình tròn có đường kính $2$ cm để cắm ống hút, mặt phẳng chứa hình tròn này song song với mặt phẳng chứa miệng ly (tham khảo hình vẽ sau). Chọn hệ trục $Oxy$ (đơn vị trên trục là centimet) với trục $Ox$ đi qua tâm của $2$ đáy hình nón cụt và gốc tọa độ $O$ trùng với tâm của đáy lớn như hình vẽ trên. Thể tích bên trong của ly bao gồm cả thể tích của nắp là bao nhiêu centimét khối \textit{(không làm tròn kết quả các phép tính trung gian, chỉ làm tròn kết quả cuối cùng đến hàng đơn vị)}? 
	 \shortans{791}
	\begin{center}
		{\begin{tikzpicture}[font=\footnotesize,line join=round,line cap=round,>=stealth,scale=0.3]
		\def\Rtop{4.5} 
		\def\Rbottom{3} 
		 \def\Rhole{1}
		 \def\hcone{9.9}
		 \def\Rcap{4.5}
		\draw (-\Rbottom,0) arc (180:360:\Rbottom  cm and 0.6cm);
		\draw[dashed]  (-\Rbottom,0) arc (180:-360:\Rbottom cm and 0.6cm);
		\draw (-\Rbottom,0) -- (-\Rtop,\hcone); 
		\draw  (\Rbottom,0) -- (\Rtop,\hcone);
		\draw (-\Rtop,\hcone) arc (180:360:\Rtop cm and 0.9cm); 
		\draw  [dashed](-\Rtop,\hcone) arc (180:0:\Rtop cm and 0.9cm);
		 \draw (-\Rtop,\hcone) arc (180:104:\Rtop);
		\draw  (\Rtop,\hcone) arc (0:78:\Rtop);
		 \path (-1,\hcone+\Rcap-0.1) coordinate (top);
		 \draw ([xshift=-\Rhole]top) arc (180:360:\Rhole cm and 0.25cm);
		 \path (-1,\hcone+\Rcap-0.1) coordinate (top);
		\draw ([xshift=-\Rhole]top) arc  (180:-360:\Rhole cm and 0.25cm);
		\draw[<->] (\Rtop+0.3,0) -- (\Rtop+0.3,\hcone)  node[midway,right] {$13{,}4$ cm};
		\draw[<->] (-\Rbottom,-0.7) -- (\Rbottom,-0.7) node[midway,below]  {$6$ cm};
		\draw[<->] (-\Rtop,\hcone) -- (\Rtop,\hcone) node[midway,above=3mm] {$9$  cm};
		 \draw[<->] (-\Rhole,\hcone+\Rcap+0.5) -- (\Rhole,\hcone+\Rcap+0.5) node[midway,above] {$2$ cm};
	\end{tikzpicture}}
		\hspace{0.5cm} 
		{\begin{tikzpicture}[font=\footnotesize,line join=round,line cap=round,>=stealth,scale=0.3,rotate=-90]
		\draw[->] (5,9.9) -- (-5.5,9.9) node[left]  {$y$};
		\draw[->]  (0,-1) -- (0,16) node[below right] {$x$};
		\draw[dashed](0,0) node[below]{$A$}--(-3,0) node[above]{$B$} 
			 (0,9.9) node[below left=2mm]{$O$}
			(-4.5,9.9) node[above right]{$C$} 
			 (-1,14.4) node[right]{$D$}--(0,14.4);
		\def\Rtop{4.5} 
		\def\Rbottom{3} 
		\def\Rhole{1} 
		\def\hcone{9.9} 
		\def\Rcap{4.5} 
		\draw (-\Rbottom,0) arc (180:360:\Rbottom cm and 0.6cm); 
		\draw[dashed] (-\Rbottom,0) arc (180:-360:\Rbottom cm  and 0.6cm);
		\draw (-\Rbottom,0) -- (-\Rtop,\hcone); 
		 \draw (\Rbottom,0) -- (\Rtop,\hcone);
		\draw  (-\Rtop,\hcone) arc (180:360:\Rtop cm and 0.9cm);
		\draw [dashed](-\Rtop,\hcone) arc  (180:0:\Rtop cm and 0.9cm);
		\draw (-\Rtop,\hcone)  arc (180:104:\Rtop);
		 \draw (\Rtop,\hcone) arc (0:78:\Rtop);
		 \path (-1,\hcone+\Rcap-0.1) coordinate (top);
		\draw ([xshift=-\Rhole]top) arc (180:360:\Rhole cm and  0.25cm);
		\path (-1,\hcone+\Rcap-0.1) coordinate (top); 
		\draw ([xshift=-\Rhole]top) arc (180:-360:\Rhole  cm and 0.25cm);
	\end{tikzpicture}}
	\end{center}
	\loigiai{ 
		Khi đó toạ độ các điểm  $A$, $B$, $C$ lần lượt là $A(-13{,}4; 0)$; $B(-13{,}4; 3)$; $C(0; 4{,}5)$.\\
		 Gọi phương trình đường thẳng đi qua $BC$ có dạng $y=ax+b$.\\
		Ta có  hệ phương trình $\heva{&-13{,}4a+b=3\\&0a+b=4{,}5} \Leftrightarrow\heva{&a=\dfrac{15}{134} \\&b=\dfrac{9}{2}.}$\\
		Phương  trình đường thẳng $BC$ là $y=\dfrac{15}{134} x+\dfrac{9}{2}$.\\
		Khi đó thể tích của phần thân ly trà sữa chưa bao gồm nắp  là
		 \[
		V_1=\pi \displaystyle\int_{-13{,}4}^0\left(\dfrac{15}{134} x+\dfrac{9}{2}\right)^2\mathrm{\,d}  x=\dfrac{3\,819\pi}{20}.
		\] 
		Điểm $C$, $D$ thuộc đường tròn  tâm $O(0; 0)$, bán kính $R=4{,}5$ có phương trình $x^2+y^2=\dfrac{81}{4}$.\\
		Thay $y=1$, ta được  $x=\sqrt{\dfrac{81}{4}-1}=\dfrac{\sqrt{77}}{2} \Rightarrow D\left(\dfrac{\sqrt{77}}{2}; 1\right)$. \\
		 Suy ra $y=\sqrt{\dfrac{81}{4}-x^2}$.\\
		Khi đó thể tích nắp của ly trà sữa là  $V_2=\pi \displaystyle\int_0^{\frac{\sqrt{77}}{2}}\left(\sqrt{\dfrac{81}{4}-x^2}\right)^2\mathrm{\,d}x=\dfrac{83\pi \sqrt{77}}{12}$.\\
		Từ $(1)$ và $(2)$ ta suy ra thể tích bên trong của ly bao gồm cả thể tích của  nắp là
		\[ 
		V=V_1+V_2=\dfrac{3\,819\pi}{20}+\dfrac{83\pi  \sqrt{77}}{12} \approx 791\,\left(\mathrm{cm}^3\right).
		\] 
	 }
\end{bt}

% nguồn: Phan2_Lop_12_Dai_so_Chuong_4
\begin{bt}%[2D4V3-5]%[ma:ba8e466b]
	\immini{Chuẩn bị cho đêm hội diễn văn nghệ chào đón năm mới, nhóm bạn Lan đã làm một chiếc mũ \lq\lq cách điệu\rq\rq~cho ông già Noel có dáng một  khối tròn xoay. Mặt cắt qua trục của chiếc mũ như hình vẽ bên. Biết rằng $OO'=6$\,(cm), $OA=10$\,(cm), $OB=20$\,(cm), đường cong $AB$ là một phần của parabol có đỉnh là điểm $B$. Tính thể tích của chiếc mũ (\textit{làm tròn kết quả đến hàng đơn vị}).
	 }
	{\begin{tikzpicture}[line cap=butt,line join=miter,>=stealth,scale=0.8]
		\path (2.5,0) coordinate (A) 
			(0,5) coordinate  (B)
			(0,-1.5) coordinate (O') 
			(-2.5,0)  coordinate (C)
			 (2.5,-1.5) coordinate (D)
			(-2.5,-1.5) coordinate (E) 
			 (0,0) coordinate (O);
		\draw (C)--(A)--(D)--(E)--cycle (O')--(B); 
	\draw (C) to[bend right=20] (B);
\draw (A) to[bend left=20] (B);
\foreach \t/\g in {B/0,O'/30,O/30,A/30}{ 
	\draw[fill=white] (\t) circle (1pt)  node[shift={(\g:9pt)},font=\scriptsize]{$\t$};
 }
\end{tikzpicture}}
	\shortans{3142} 
	 \loigiai{
	\immini{Chọn  hệ trục tọa độ như hình vẽ bên.
		\\ 
		Parabol $(P)$ đi qua $2$ điểm $A$, $B$ và có  $B$ là đỉnh.\\ Suy ra $(P)$ đi qua $A(20;10)$ và $B(0;0)$.\\
		Ta có $(P)\colon  y=ax^2+bx+c$.\\
		Vì  $(P)$ qua $B(0;0)\Rightarrow c=0$.\\
		Vì $(P)$ qua $A(20;10)$ nên $10=20^{2}a+20b$.  \hfill(1)\\
		Mà $B$ là  đỉnh của parabol $(P)$ nên $\dfrac{-b}{2a}=0\Rightarrow b=0$.\\
		Thay $b=0$ vào  (1) ta được $a=\dfrac{1}{40}$.\\
		Khi đó  $(P)\colon y=\dfrac{x^2}{40}$.\\
		$V=\pi\displaystyle\int_{0}^{20}\left(\dfrac{x^2}{40}\right)^2\,\mathrm{d}x+\pi\displaystyle\int_{20}^{26}10^2\,\mathrm{d}x=1\,000\pi\approx 3\,142$ (cm$^3$).\\ 
		Vậy thể tích của chiếc mũ là $3\,142$ (cm$^3$).}{\begin{tikzpicture}[scale=.8,line cap=butt,line join=miter,>=stealth]
		\draw[->] (-1,0)--(7.5,0) 
			 node[shift={(-100:7pt)},font=\normalsize]{$x$};
		\draw[->] (0,-3)--(0,5) 
			node[shift={(180:7pt)},font=\normalsize]{$y$}; 
		\path (0,0) coordinate  (B)
			(5,0) coordinate (O) node[above  right,font=\scriptsize]{$20$}
			(6.5,0) coordinate (O') node[above  right,font=\scriptsize]{$26$}
			(5,2.5) coordinate (A) node[above right,font=\scriptsize]{$A (20;  10)$}
			(6.5,2.5) coordinate  (C) node[above right,font=\scriptsize]{$C (26; 10)$}
			(5,-2.5) coordinate  (D)
			(6.5,-2.5)  coordinate (E);
		\draw (A)--(C)--(E)--(D)--cycle; 
	\draw (B) to[bend right=20] (A);
\draw (B) to[bend left=20] (D);
\draw[fill=white] (O) circle  (1pt);
\draw[fill=white] (O')  circle (1pt);
\draw[fill=white] (A) circle  (1pt);
\draw[fill=white] (C) circle (1pt); 
\foreach \t/\g in {B/120}{ 
	\draw[fill=white] (\t) circle (1pt)  node[shift={(\g:9pt)},font=\scriptsize]{$\t$};
} 
\end{tikzpicture}}
	} 
\end{bt}

% nguồn: Phan2_Lop_12_Dai_so_Chuong_4
\begin{bt}%[2D4V3-5]%[ma:3b979c8a]
	Một bồn hình trụ đang chứa dầu, được đặt nằm ngang, có chiều dài bồn là $5$ m, có bán kính đáy $1$ m.  Người ta đã rút dầu trong bồn tương ứng với $0{,}5$ m của đường kính đáy. Tính thể tích của lượng dầu còn lại trong bồn (theo đơn vị m$^3$).
	\begin{center}
		\begin{tikzpicture}[>=stealth,scale=1,rotate=-90]
				\coordinate (O')  at (0,0) circle (1pt) node[left]{$O'$};
				\coordinate  (D) at (-120:2 and 2/2);
				\coordinate  (C) at (140:2 and 2/2);
				\coordinate (A) at  ($(D) + (0,5)$);
				\coordinate (B)  at ($(C) + (0,5)$);
				 \draw (A) node[below right]{$A$}--(D) node[left]{$D$};
				 \draw[dashed](C) node[left]{$C$}--(D);
				 \draw[dashed](C)--(B) node[right]{$B$};
				\draw (A)--(B); 
				\draw[dashed](O')--(0,5)  node[right]{$O$};
				\draw[dashed](-2,0) arc (180:0:2 and 1); 
				\draw  (-2,0) arc (-180:0:2 and 1);
				\draw (-2,5) arc (180:0:2  and 1);
				\draw (-2,5) arc (-180:0:2 and  1);
				 \draw (-2,0)--(-2,5)
					(2,0)--(2,5) 
					(-2,5)  node[above]{$H$}--(0,5);
			\end{tikzpicture}
	\end{center}
	\shortans{12,6} 
	\loigiai{ 
		Dung tích của  bồn (hình trụ) đựng dầu là $V=\pi{r^2}h=\pi{1^2}\cdot5=5\pi $ m$^3$.\\
		Thể  tích phần đã rút dầu ra (phần trên mặt phẳng $\left(ABCD\right)$) là $V_1=\left(\dfrac{\pi}{3}-\dfrac{\sqrt 3}{4}\right)\cdot5$ m$^3$.\\
		Vậy thể tích lượng dầu còn lại là $V_2=V-V_1\approx 12{,}6$ m$^3$. 
	} 
\end{bt}

% nguồn: Phan2_Lop_12_Dai_so_Chuong_4
\begin{bt}%[2D4V3-5]%[ma:be234609]
	Một chiếc bát thủy tinh có bề dày của phần xung quanh là một khối tròn xoay, khi xoay hình phẳng $D$ quanh một đường thẳng $a$ bất kì nào đó mà khi gắn hệ trục tọa độ $Oxy$ (đơn vị trên trục là decimet) vào hình phẳng $D$ tại một vị trí thích hợp, thì đường thẳng $a$ sẽ trùng với trục $Ox$. Khi đó hình phẳng  $D$ được giới hạn bởi các đồ thị hàm số $y = x + \dfrac{1}{x}$, $y = x$ và hai đường thẳng $x=1$, $x=4$ như hình dưới đây.
	\begin{center}
		\begin{tikzpicture}[line join=round,line cap=round,>=stealth]
				 \tikzset{every node/.style={scale=0.9}}
				\draw[->] (-1.1,0)--(6.1,0) node[below left]  {$x$};
				\draw[->] (0,-5.1)--(0,5.1)  node[below left] {$y$};
				\draw (0,0)  node [below left] {$O$};
				\foreach  \x/\nx in {1/1,2/2,3/3,4/4}
				\draw[thin] (\x,1pt)--(\x,-1pt) node [below] {$\nx$}; 
				\foreach \y/\ny in {1/1,2/2} 
				\draw[thin] (1pt,\y)--(-1pt,\y)  node [left] {$\ny$};
				\draw[dashed,thin] (0.01,-5)--(0.01,5); 
				\draw (4,0)--(4,4); 
				\begin{scope}
					\clip (-1,-5) rectangle (6,5); 
					 \draw[pattern=north west lines,pattern color=cyan!65] (1,0)--plot[domain=1:4] (\x,{(1*((\x)^2) +0*(\x)+1)/(1*(\x)+0)})--(4,4)--(1,1)--cycle;
					\draw[samples=100,domain=0.01:6,smooth,variable=\x]  plot (\x,{(1*((\x)^2) +0*(\x)+1)/(1*(\x)+0)});
					 \draw[dashed,samples=100,domain=1:6,smooth,variable=\x] plot (\x,{(-1*((\x)^2)-0*(\x)-1)/(1*(\x)+0)});
					\draw (0,0)--(6.1,6.1); 
					\draw[dashed,samples=100,domain=0:360,smooth,variable=\t] 
						plot ({1 + 0.5*cos (\t)},{2*sin (\t)}); 
					\draw[dashed,samples=100,domain=0:360,smooth,variable=\t] 
						plot ({1 + 0.25*cos  (\t)},{1*sin (\t)});
					 \draw[dashed,samples=100,domain=0:360,smooth,variable=\t]
						plot ({4 +  0.5*cos (\t)},{4.25*sin (\t)});
					 \draw[dashed,samples=100,domain=0:360,smooth,variable=\t]
						plot  ({4 + 0.4*cos (\t)},{4*sin (\t)});
					\draw[dashed] (0,0)--(6.1,-6.1); 
				\end{scope}
			\end{tikzpicture}
	\end{center}
	Thể tích của bề dày  chiếc bát thủy tinh đó bằng bao~nhiêu decimet khối? (làm tròn kết quả đến hàng phần mười).
	\shortans{18} 
	\loigiai{ 
		Hình phẳng $D$ giới hạn bởi các đường $y = f(x) = x + \dfrac{1}{x}$, $y  = g(x) = x $, $x=1$, $x=4$.\\
		Khi quay hình phẳng $D$ quanh trục $Ox$, ta được khối tròn  xoay là phần bề dày của chiếc bát.\\
		Thể tích của khối tròn xoay 
		$$V = \pi \displaystyle \int_{1}^{4} \left[f^2(x) - g^2(x) \right] \mathrm{\,d}x= \pi \displaystyle\int_{1}^{4} \left(2 + \dfrac{1}{x^2}\right) \mathrm{\,d}x= \pi \left[2x  - \dfrac{1}{x} \right]\bigg|_{1}^{4}= \dfrac{27\pi}{4}.$$
		Tính giá trị xấp  xỉ: $V = \dfrac{27\pi}{4} \approx 21{,}2$ (decimet khối).\\
		Làm tròn kết quả đến hàng phần  mười (một chữ số thập phân), ta được $21{,}2$. \\
		Vậy thể tích của bề dày chiếc  bát thủy tinh đó bằng $21{,}2$ decimet khối.
	 }
\end{bt}

% nguồn: Phan2_Lop_12_Dai_so_Chuong_4
\begin{bt}%[2D4V3-5]%[ma:d22c4908]
	Khuôn viên nhà bạn Thùy Dương có dạng nửa hình tròn, trên đó người thiết kế phần để trồng hoa có dạng của một cánh hoa hình parabol có đỉnh trùng với tâm và có trục đối xứng vuông góc với  đường kính của nửa hình tròn, hai đầu mút của cánh hoa nằm trên nửa đường tròn và cách nhau một khoảng bằng $4$\,m. Phần còn lại của khuôn viên dành để trồng cỏ Nhung Nhật. Biết các kích thước cho như hình vẽ, chi phí trồng hoa và cỏ Nhung Nhật tương ứng là $250\,000$\,đồng/m$2$ và $150\,000$\,đồng/m$^2$. Hỏi chi phí để trồng hoa và trồng cỏ Nhung Nhật trong khuôn viên đó hết bao nhiêu triệu đồng (làm tròn kết quả đến hàng phần chục)?
	\begin{center}
		\begin{tikzpicture}[line join=round,line cap=round,>=stealth]
				 \tikzset{every node/.style={scale=0.9}}
				\begin{scope}
					\clip (-4.5,-0.1) rectangle (4.5,4.5); 
					 \fill[gray!40]plot[samples=200,domain=-2:2,smooth,variable=\x] (\x,{sqrt (20-(\x)^2)})--plot[samples=200,domain=2:-2,smooth,variable=\x] (\x,{(\x)^2})--cycle;

					 \draw (4.472135955,0) arc (0:180:4.472135955);
					 \draw plot[samples=200,domain=-2:2,smooth,variable=\x] (\x,{(\x)^2});
					 \draw (-4.472135955,0)--(4.472135955,0);
				\end{scope}
				 \draw[dashed](-2,0)--(-2,4) node[pos=0.5,left]{$4$\,m}
					(-2,4)--(2,4)  node[pos=0.5,above]{$4$\,m}
					(2,0)--(2,4)  node[pos=0.5,right]{$4$\,m};
			\end{tikzpicture}
	\end{center}
	 \shortans{5,9}
	 \loigiai{
		Chọn hệ trục tọa độ $Oxy$ như  hình vẽ.
		\begin{center}
			\begin{tikzpicture}[line join=round,line cap=round,>=stealth]
					 \tikzset{every node/.style={scale=0.9}}
					\begin{scope}
						\clip (-5.5,-1)  rectangle (5.5,5.5);
						\fill[gray!40]plot[samples=200,domain=-2:2,smooth,variable=\x] (\x,{sqrt (20-(\x)^2)})--plot[samples=200,domain=2:-2,smooth,variable=\x] (\x,{(\x)^2})--cycle; 

						\draw (4.472135955,0)  arc (0:180:4.472135955);
						\draw  plot[samples=200,domain=-2:2,smooth,variable=\x] (\x,{(\x)^2});
					\end{scope}
					\draw[dashed](-2,0)--(-2,4) node[pos=0.5,left]{$4$\,m} 
						(-2,4)--(2,4) node[pos=0.6,above]{$4$\,m} 
						 (2,0)--(2,4) node[pos=0.5,right]{$4$\,m};
					 \draw[->] (-5.6,0)--(5.6,0) node[below left] {$x$};
					\draw[->] (0,-1.1)--(0,5.6) node[below  left] {$y$};
					\draw (0,0) node [below left] {$O$}--(2,4); 
					 \draw (2,4) circle (1.0pt) node [above right] {$A$}
						 (2,0) circle (1.0pt) node[below]{$B$};
				\end{tikzpicture}
		\end{center}
		Xét parabol có đỉnh là gốc tọa độ  và đi qua điểm $(2;4)$:\\
		Phương trình $(P)$ có  dạng $y=ax^2$.\\
		Vì $A(2;4) \in(P)$ nên $4=a\cdot 2^2\Leftrightarrow  a=1$.\\
		Suy ra  $(P)\colon y=x^2$.\\
		Xét $\triangle OAB$ vuông tại $B$, theo định  lý Pythagore ta có
		\[OA^2=AB^2+OB^2=4^2+2^2=16+4=20\Rightarrow  R=OA=2\sqrt 5.\]
		 Phương trình đường tròn $(C)$ có tâm $O$, bán kính $R=2\sqrt 5$ là $x^2+y^2=20$.\\
		 Suy ra $|y|=\sqrt{20-x^2}$.\\
		Diện  tích phần trồng hoa là $S_1=\displaystyle\int_{-2}^2\left(\sqrt{20-x^2}-x^2\right)\mathrm{\,d}x$.\\
		Diện tích trồng cỏ  Nhung Nhật là $S_2=\dfrac{\pi R^2}{2}-S_1=10\pi-S_1$.\\
		Tổng chi phí để trồng hoa và trồng cỏ của Nhung Nhật  trong khuôn viên đó là
		\allowdisplaybreaks 
		\begin{eqnarray*}
			250\,000\cdot S_1+150\,000\cdot  S_2&=&250\,000\cdot S_1+150\,000 \cdot\left(10 \pi-S_1\right)\\
			&\approx&5\,906\,351\text{ (đồng)}\\ 
			&\approx&5{,}9\text{  (triệu đồng).}
		\end{eqnarray*}
		Vậy chi phí để trồng hoa và trồng cỏ Nhung Nhật trong khuôn viên nhà bạn  Thùy Dương gần bằng $5{,}9$ triệu đồng.
	} 
\end{bt}

% nguồn: Phan2_Lop_12_Dai_so_Chuong_4
\begin{bt}%[2D4V3-5]%[ma:94d2ff20]
	\immini[thm]{Viên gạch men dùng để lát nền nhà là một hình vuông có cạnh bằng $80$\,cm (xem hình bên dưới). Mỗi viên gạch có $4$ bông hoa, mỗi bông hoa gồm $4$ cánh hoa. Mỗi cánh hoa (phần màu xanh) là phần giao nhau của hai hình tròn có cùng bán kính và khoảng cách giữa hai tâm là $20\sqrt 2$\,cm. Ước tính ở công đoạn tráng men, phần màu xanh có chi phí $50$ nghìn đồng trên một mét vuông, còn phần màu trắng có chi phí $30$ nghìn đồng trên một mét vuông. Tính chi  phí (đơn vị: tỉ đồng) của công đoạn tráng men này, khi cơ sở dự định sản xuất $100\,000$ viên gạch như thế (làm tròn kết quả đến hàng phần trăm).
	}{ 
	\begin{tikzpicture}[line join=round,line cap=round,>=stealth]
			\begin{scope}
				\foreach \x/\y in {0/0,4/0,0/4,4/4}{ 
					 \foreach \i/\j/\k/\l in {0/90/-90/-180,-90/0/-180/-270,-270/-360/180/270,0/-90/90/180}{
					\fill[green!60!black](\x,\y) arc  (\i:\j:2)--(\x,\y) arc (\k:\l:2)--cycle;
					\draw (\x,\y) arc (\i:\j:2) 
					(\x,\y) arc  (\k:\l:2);
				 }
				 }
			\end{scope}
			\draw (-2,-2) rectangle  (6,6);
		\end{tikzpicture}
	 }
	 \shortans{2,65}
	 \loigiai{
	\immini{Xét $\dfrac{1}{16}$ viên gạch men (vùng hình vuông  chứa một cánh hoa).
		Vì phần cánh hoa (màu xanh) là phần giao nhau của các đường tròn có khoảng cách giữa các tâm là $20\sqrt 2$ nên một cánh hoa là phần giao của hai phần tư hình tròn tâm $A$, $B$ (hình vẽ), có  bán kính $r=20$.\\
		Mỗi phần  tư hình tròn tâm $A$, $B$ có diện tích $S=\dfrac{1}{4}\pi \cdot 20^2$.
		}{ 
		\begin{tikzpicture}[line join=round,line cap=round,>=stealth]
				\begin{scope}
					 \fill[green!60!black](0,0) arc (-90:0:2)--(0,0) arc (-180:-270:2)--cycle;
					 \draw (0,0) arc (-90:0:2)
						(0,0)  arc (-180:-270:2);
				\end{scope}
				 \draw (0,0) rectangle (2,2);
				\draw (0,2) circle (1.0pt)  node[above left]{$A$}
					 (2,0) circle (1.0pt) node[below right]{$B$};
			\end{tikzpicture}
		 }\noindent
		Diện tích hình  vuông có đường chéo $AB$ là $S_v=20^2$.\\
		Suy ra diện  tích của một cánh hoa là \[2S-S_v=\dfrac{1}{2}\pi \cdot 20^2-20^2=\left(\dfrac{\pi}{2}-1\right) \cdot 20^2.\]
		Do  đó
		\begin{itemize}
			\item Diện tích phần màu  xanh là $S_1=16\left(\dfrac{\pi}{2}-1\right) \cdot 20^2$\,(cm$^2$).
			\item Diện tích phần màu trắng  $S_2=80^2-S_1$\,(cm$^2$).
		\end{itemize}
		Số tiền tráng  men một viên gạch là $T=5S_1+3S_2=26\,506,19298$ đồng.\\
		Để sản xuất $100\,000$ viên gạch thì số tiền tráng men là \[100\,000T=2\,650\,619\,298\text{  (đồng)}\approx 2{,}65\text{ (tỉ đồng).}\]
	} 
\end{bt}

% nguồn: Phan2_Lop_12_Dai_so_Chuong_4
\begin{bt}%[2D4V3-5]%[ma:1e7ee87b]
	\immini[thm]{Mái  chùa đối xứng qua trục, có biên dạng $f(x)=6 \cdot \mathrm{e}^{-0,5x}+1$, $0 \le x \le 6$.
	Cổng chùa dạng parabol, rộng  $4$ m và cao $3$ m.
	Mỗi cửa  sổ kích thước $1\text{ m} \times 0{,}5$m.
	Mặt tiền phía trước  của mặt sẽ được sơn màu vàng dày $0{,}1$ mm.
	Hỏi cần bao nhiêu lít nước sơn để sơn mặt tiền của chùa (không sơn các cửa)? (làm tròn kết quả đến hàng phần mười).}{\begin{tikzpicture}[scale=.7,>=stealth,font=\footnotesize,line join=round,line cap=round]
		\draw[domain=-5:5,smooth,variable=\x] 
			plot  ({\x},{6*exp (-0.5*abs (\x))+1});
		\fill[yellow!70,draw] 
			plot[domain=-5:5,samples=200] ({\x},{6*exp (-0.5*abs (\x))+1}) 
			-- (5,0) -- (-5,0)  -- cycle;
		 \fill[green!40,draw]
			plot[domain=-2:2,samples=100]  ({\x},{-0.75*(\x)^2+3})
			-- (2,0) -- (-2,0) -- cycle; 
		 \fill[pink!70,draw] (-4,0) rectangle (-3,1.5);
		\fill[pink!70,draw] (3,0) rectangle (4,1.5); 
		\fill[pink!70,draw] (-0.5,3.5) rectangle  (0.5,5);
		\draw[->,red] (-5.5,0)--(5.5,0)  node[right]{$x$};
		\draw[->,red]  (0,-0.5)--(0,7) node[above]{$y$};
	\end{tikzpicture}
	} 
	\shortans{2,5} 
	\loigiai{ 
		Ta cần diện tích mặt tiền trừ diện tích cổng và ba  cửa sổ.\\
		Diện tích mặt tiền $A_1 = 2 \displaystyle\int_{0}^{6} (6\mathrm{e}^{-0,5x}+1)\mathrm{\,d}x =  2 \left(-12\mathrm{e}^{-0,5x}+x\right)\bigg|_0^6 \approx 34{,}81\left(\text{m}^2\right)$.\\
		 Phương trình, diện tích cổng.\\
		 Giả sử $g(x)=ax^2+b$.\\
		Theo  giả thiết có $g(0)=3 \Rightarrow b=3$; $g(2)=0 \Rightarrow a=-\dfrac{3}{4}$.\\
		Vậy  $g(x)=-\dfrac{3}{4}x^2+3$.\\
		Diện tích cổng là $A_2 = 2 \displaystyle\int_{0}^{2} \left(-\dfrac{3}{4}x^2+3\right)\mathrm{\,d}x = 2 \left(-\dfrac{1}{4}x^3+3x\right)\bigg|_{0}^{2} =  8\left(\text{m}^2\right)$.\\
		Ba cửa sổ, mỗi cái $0{,}5$ m$^2 \Rightarrow  1{,}5$ m$^2$.\\
		Vậy diện tích cần sơn vàng $A = A_1 - A_2 -  1{,}5 = 34{,}81 - 8 - 1{,}5 = 25{,}31\left(\text{m}^2\right)$.\\
		Ta có $25{,}31 \text{ m}^2 = 2\,531 \left(\text{dm}^2\right)$, dày $0{,}1 \text{ mm} =  0{,}001$ (dm).\\
		Thể tích nước sơn  bằng
	Vậy $V = A \cdot \text{ dày} = 2\,531 \cdot 0{,}001 = 2{,}531 \left(\text{dm}^3\right) \approx 2{,}5  \left(\text{lít}\right)$.}
\end{bt}

% nguồn: Phan2_Lop_12_Dai_so_Chuong_4
\begin{bt}%[2D4V3-5]%[ma:babe66fc]
	Nếu cắt chậu nước có hình dạng như hình bên bằng mặt phẳng song song và cách mặt đáy $x$ cm $(0 \leq x \leq 16)$ thì mặt cắt là hình tròn có bán kính $R=10+\sqrt{x}$ cm. Tìm $x$ (đơn vị cm, làm tròn kết quả đến hàng phần trăm) để dung  tích nước trong chậu bằng $\dfrac{1}{2}$ thể tích của chậu?
	\begin{center}
		\begin{tikzpicture}[scale=1,font=\footnotesize,line join=round,line cap=round,>=stealth]
				\pgfmathsetmacro{\r}{2*sqrt (2)} 
				\fill[color=green!20] (-2.6,-1)--(2,-2) arc (0:180:2  and 0.3)
					 (2.6,-1)--(2.63,-1) arc (-20:-45:{\r})--(2,-2) arc (0:-180:2 and 0.3)--(-2.63,-1) arc (200:225:{\r})--(-2.6,-1);
				 \draw [densely dashed,fill=green!20] (2,-2) arc (0:180:2 and 0.3); \draw[fill=green!20] (2,-2) arc (0:-180:2 and 0.3);
				\draw [densely dashed,fill=green!20] (2.65,-1) arc  (0:180:2.65 and 0.3);
				\draw[fill=green!20] (2.65,-1) arc (0:-180:2.65 and 0.3); 
				\draw (0:\r) arc (0:180:{\r} and  0.4);
				\draw (0:\r) arc (0:-180:{\r} and  0.4);
				 \draw (-2.82,0) arc (180:225:{\r});
				\draw  (2.82,0) arc (0:-45:{\r});
				\draw[<->]  (-5,0)--(-5,-2);
				\draw[<->] (-3.5,-1)--(-3.5,-2); 
				\draw[dashed] (0,0)--(0,-2) 
					(0,0)--(-5,0) 
					 (0,-2)--(-5,-2)
					 (0,-1)--(-3.5,-1);
				\draw (-5,-0.5) node[right]{$16$ cm} 
					(-4.9,-1.5) node[right]{$x$  cm}
					(-0.25,-0.75)  node[left]{$10+\sqrt{x}$ cm};
			\end{tikzpicture}
	\end{center}
	 \shortans{8,94}
	\loigiai{ 
		Diện tích mặt cắt  của chậu là $S(x)=\pi \cdot(10+\sqrt{x})^{2}=\pi \cdot(100+x+20 \sqrt{x})$\\
		Thể tích của chậu là $V=\displaystyle\int_{0}^{16} S(x) \mathrm{\,d} x=\left.\pi  \cdot\left(100 x+\dfrac{x^{2}}{2}+20 \cdot \dfrac{2}{3} \cdot \sqrt{x^{3}}\right)\right|_{0}^{16}=\dfrac{7744}{3} \pi$.\\
		Dung  tích nước trong chậu là
		 \[V_{\text {nước}}=\displaystyle\int_{0}^{x} S(x) \mathrm{\,d} x=\pi \cdot\left(100 x+\dfrac{x^{2}}{2}+\dfrac{40}{3} \cdot \sqrt{x^{3}}\right).\]
		Để dung tích  nước trong chậu bằng $\dfrac{1}{2}$ thể tích của chậu thì
		\[\pi \cdot  \left(100 x+\dfrac{x^{2}}{2}+\dfrac{40}{3} \cdot \sqrt{x^{3}}\right)=\dfrac{3872}{3} \pi \Leftrightarrow 100 x+\dfrac{x^{2}}{2}+\dfrac{40}{3} \cdot \sqrt{x^{3}}=\dfrac{3872}{3} \Leftrightarrow x \approx 8{,}94.\]
	 }
\end{bt}

% nguồn: Phan2_Lop_12_Dai_so_Chuong_4
\begin{bt}%[2D4V3-5]%[ma:6b09f436]
	Một nhà kho có dạng khối hộp chữ nhật cùng phần mái phía trên được thiết kế dạng vòm cong. Chiều dài của xưởng là $40$ m, chiều rộng $10$ m, chiều cao $5$ m; đường cắt của mái vòm khi cắt bởi mặt phẳng  vuông góc với chiều dài của kho là một cung tròn, điểm cao nhất của mái vòm có chiều cao là $6$ m (hình vẽ minh họa).
	\begin{center}
		\begin{tikzpicture}[line join=round,line cap=round,thin,scale=0.6]
				\coordinate (A)  at (0,0);
				 \coordinate (B) at (7,0);
				 \coordinate (D) at (-4,4);
				\coordinate (C)  at ($(B)+(D)-(A)$);
				\coordinate (M) at (0,5.1); 
				\coordinate (N)  at (7,5);
				 \coordinate (Q) at ($(D)+(0,5.1)$);
				 \coordinate (P) at ($(C)+(0,5)$);
				\coordinate (K) at ($(C)!0.5!(D)$); 
				\fill[brown!20] (A)--(B)--(C)--(D)--cycle; 
				\draw[teal,dashed]  (P)--(C);
				\draw[teal,dashed]  (D)--(C)--(B);
				 \draw[teal] (D)--(A)--(B)--(N)--(P);
				 \draw[teal] (M)--(A);
				 \draw[teal] (M)--(Q)--(D);
				 \draw[teal] (N) arc (66.99:109.5:9.7);
				\draw[teal]  (P) arc (66.99:110:9.6);
				\coordinate (K_front)  at (3.5,5.8);
				\coordinate  (K_back) at ($(Q)!0.5!(P)+(0,0.7)$);
				\draw[teal] (K_front)--(K_back); 
				\draw[teal] (A) -- (M); 
				\node[sloped,below=2mm] at ($(A)!0.7!(D)$) {$40$ m}; 
				\node[sloped,below=2mm] at ($(B)!0.5!(A)$)  {$10$ m};
				 \coordinate (H_back) at ($(Q)!0.5!(P)$);
				\fill  (1,7.3) node[left]{$6$ m};
				 \fill (-4,7) node[left]{$5$ m};
				\draw[teal,dashed] (K)--($(H_back)+(0,0.7)$); 
			\end{tikzpicture}
	\end{center}
	Để tính toán công suất cho hệ thống làm lạnh nhà kho, người ta cần ước lượng thể tích phần không gian bên trong nhà kho. Em hãy cho biết thể tích không gian bên trong nhà kho là bao nhiêu  m$^{3}$? (\textit{kết quả làm tròn đến hàng đơn vị}).
	\shortans{2269} 
	 \loigiai{
		Gắn  hệ trục $Oxy$ như hình vẽ. Phần mái của nhà kho là một phần của đường tròn $(C)$ có phương trình \[x^{2}+y^{2}+ax+by+c=0.\]
		\begin{center}
			\begin{tikzpicture}[line join=round,line cap=round,thin,scale=0.4]
					\def\r{sqrt  (50)}
					\coordinate  (O) at (0,0);
					\coordinate (A) at  (-5,0);
					\coordinate (D)  at (5,0);
					\coordinate (B) at  (-5,5);
					\coordinate (C) at (5,5); 
					 \coordinate (H) at (0,6);
					\draw[teal,->,blue] (-8,0) -- (8,0) node[below] {$x$}; 
					\draw[teal,->,blue] (0,-8) --  (0,9) node[left] {$y$};
					\draw[teal,green!60!black]  (O) circle (\r);
					\draw[teal] (A) -- (B) -- (C)  -- (D);
					\node[right,above=2pt] at (-4.5,0) {$-5$}; 
					\node[above=2pt] at (5.5,0) {$5$}; 
					 \node[above=2pt] at (-0.2,5) {$5$};
					 \node[above=2pt] at (-0.5,6.7) {$6$};
					\node[right,above=2pt] at (0.3,7) {$H$}; 
					\foreach \i/\g in {O/225,A/-90,B/135,C/45,D/-90}{ 
						 \draw[teal,fill=white] (\i) circle (1.5pt);
						\node at ($(\i)+(\g:8mm)$)  {$\i$};
					} 
				\end{tikzpicture}
		\end{center}
		Vì $(C)$ đi qua các điểm $B(-5;5)$, $C(5;5)$, $H(0;6)$ nên ta có hệ phương trình 
		\allowdisplaybreaks 
		\begin{eqnarray*}
			&&\heva{& 25+25-5a+5b+c=0  \\& 25+25+5a+5b+c=0 \\& 36+6b+c=0}\Leftrightarrow\heva{&
			-5a+5b+c = -50  \\&
			 5a +5b+c=-50 \\&
			6b+c=-36} 
			\Leftrightarrow 
			 \heva{&
			a = 0 \\& 
			 b=14 \\&
			c=-120.}\\ 
			&\Rightarrow& (C)\colon  x^{2}+y^{2}+14y-120=0 \Leftrightarrow x^{2}+(y+7)^{2}=169.\\ &\Leftrightarrow &\hoac{&y=\sqrt{169-x^{2}}-7 \\& y=-\sqrt{169-x^{2}}-7.}
		\end{eqnarray*}
		Vậy thể tích không gian bên trong nhà kho là \[V=S \cdot h=\left[\displaystyle\int_{-5}^{5}\left(\sqrt{169-x^{2}}-7\right)\mathrm{\,d}x\right]  \cdot 40\approx 2\,269 \text{ (m$^3$)}.\]
	 }
\end{bt}

% nguồn: Phan2_Lop_12_Dai_so_Chuong_4
\begin{bt}%[2D4V3-5]%[ma:7d7ac239]
	Nhân dịp đi dã ngoại, lớp 12A dự kiến dựng một cái trại, các khung xương có dạng hình parabol như hình vẽ. Nền của lều trại là một hình chữ nhật có kích thước bề ngang $3$ mét, chiều dài $5$ mét, đỉnh trại cách nền $3$ mét. Thể tích phần không gian bên  trong lều trại bằng bao nhiêu mét khối?
	\begin{center}
		\begin{tikzpicture}[scale=0.65,font=\footnotesize,line join=round,line cap=round,>=stealth]
				\def\f{-7/6*(\x)^2+5/6*(\x)+3} 
				\def\g{-7/6*(\x-6)^2+5/6*(\x-6)+5} 
				 \def\s{-7/6*(\x-3)^2+5/6*(\x-3)+4}
				\coordinate  (A) at (-1,1);
				\coordinate (B)  at (2,0);
				\coordinate (C) at (5,3); 
				\coordinate (D) at  (8,2);
				\coordinate (I) at (5/14,529/168); 
				\coordinate (J) at  (89/14,865/168);
				\draw[smooth,samples=200,domain=-1:2] plot  (\x,{\f});
				 \draw[smooth,samples=200,domain=5:8] plot (\x,{\g});
				\draw  (A)--(B)--(D)
					(I)--(J); 
				 \draw[dashed] (A)--(C)--(D);
			\end{tikzpicture}
	\end{center}
	 \shortans{30}
	 \loigiai{
		\begin{center}
			\begin{tikzpicture}[scale=0.65,font=\footnotesize,line join=round,line cap=round,>=stealth]
					\def\f{-7/6*(\x)^2+5/6*(\x)+3} 
					\def\g{-7/6*(\x-6)^2+5/6*(\x-6)+5} 
					 \def\s{-7/6*(\x-3)^2+5/6*(\x-3)+4}
					 \coordinate (A) at (-1,1);
					\coordinate (B)  at (2,0);
					\coordinate  (C) at (5,3);
					\coordinate (D)  at (8,2);
					 \coordinate (I) at (5/14,529/168);
					\coordinate (J)  at (91/14,865/168);
					\draw[smooth,samples=200,domain=-1:2] plot (\x,{\f}); 
					\draw[smooth,samples=200,domain=5:8] plot (\x,{\g}); 
					\draw  (A)--(B)--(D)
						 (I)--(J);
					\coordinate (M)  at ($(A)!3/5!(B)$);
					\coordinate (N)  at ($(C)!3/5!(D)$);
					\coordinate  (Q) at ($(M)+(90:2.5)$);
					 \coordinate (P) at ($(N)+(90:2.5)$);
					\coordinate (K)  at ($(C)!7/5!(D)$);
					\coordinate (H) at  ($(C)!1/2!(D)$);
					 \coordinate (Y) at ($(H)!3/2!(J)$);
					 \draw[dashed] (A)--(C)--(D)
						(H)--(J); 
					 \draw[dashed,fill=blue,opacity=0.2] (M)--(N)--(P)--(Q)--(M);
					\draw[->] (D)--(K) node  [below]{$x$};
					 \draw[->] (J)--(Y) node [right]{$y$};
					\draw[fill=white] (3.5,2.5) node[rotate=0]{\scriptsize$S  (x)$} circle (14pt);
				\end{tikzpicture}
		\end{center}
		Chọn hệ trục tọa độ như hình vẽ, hình dạng khung trại  là parabol $y=f(x)=ax^2+bx+c$.\\
		Vì đỉnh trại  cao $3$ m và bề ngang rộng $3$ m nên parabol đi qua điểm $(0;3)$ và $\left(-\dfrac{3}{2};0 \right)$ nên ta có
		 \[\heva{
		& b=0 \\ 
		&  3=c \\
		 & 0=a\cdot {{\left( \dfrac{3}{2} \right)}^2}+c}
		\Leftrightarrow  \heva{
		 & b=0 \\
		& a=-\dfrac{4}{3}  \\
		& c=3}\Rightarrow y=f(x)=-\dfrac{4}{3}x^2+3.\] 
		Ta có $f(x)=-\dfrac{4}{3}x^2+3 \geq 0,\, \forall  x \in \left[-\dfrac{3}{2}; \dfrac{3}{2}\right]$.\\
		Mỗi mặt phẳng vuông góc $Ox$ tại điểm có hoành độ $x$ , $-\dfrac{3}{2}\le x\le \dfrac{3}{2}$ cắt khối chóp theo mặt cắt là hình chữ nhật có độ dài các cạnh  lần lượt là $5$ và $ f(x)$, có diện tích $S(x)=5\cdot f(x)$, với $-\dfrac{3}{2}\le x\le \dfrac{3}{2}$.\\
		Vậy thể tích phần  không gian trong trại là \[V=\displaystyle\int_{-\tfrac{3}{2}}^{\tfrac{3}{2}}{5\cdot f(x)}\mathrm{\,d}x=5\displaystyle\int_{-\tfrac{3}{2}}^{\tfrac{3}{2}}{\left( -\frac{4}{3}x^2+3 \right)}\mathrm{\,d}x=30\,\left( \text{ m}^3\right) .\]
	} 
\end{bt}

% nguồn: Phan2_Lop_12_Dai_so_Chuong_4
\begin{bt}%[2D4V3-5]%[ma:146cfa65]
	 \immini{
	Cho đường tròn tâm $O$, bán kính bằng $2$. Trên đường tròn, lấy $6$ điểm chia đều đường tròn, lần lượt là $A$, $B$, $C$, $D$, $E$, $F$. Vẽ cung tròn $C_1$ tiếp xúc với hai đoạn thẳng $OA$ và $OB$, đi qua hai điểm $A$ và $B$. Tương tự, vẽ cung tròn $C_2$ tiếp xúc với  hai đoạn thẳng $OB$ và $OC$, đi qua hai điểm $B$ và $C$. Tiếp tục bằng cách tương tự, ta vẽ các cung tròn $C_3$, $C_4$, $C_5$, $C_6$ tiếp xúc với các cặp đoạn thẳng liên tiếp tạo bởi các điểm trên. Diện tích phần được bao bởi $6$ cung tròn $C_1$ đến $C_6$ là bao nhiêu? (làm tròn kết quả đến hàng phần trăm).
	}{ 
	\begin{tikzpicture}[scale=1,font=\footnotesize,line join=round,line cap=round,>=stealth]
			 \def\a{sqrt (3)/2} \def\b{0} \def\c{sqrt (3)/2}
			 \def\xmin{-3} \def\xmax{3}
			 \def\ymin{-3} \def\ymax{3}
			 \path
				(0,0) coordinate  (O)
				(1,1.732)  coordinate (A)
				(-1,1.732) coordinate (B) 
				 (-2,0) coordinate (C)
				 (-1,-1.732) coordinate (D)
				 (1,-1.732) coordinate (E)
				(2,0) coordinate (F); 
			\clip (\xmin+0.1,\ymin+0.1) rectangle  (\xmax-0.1,\ymax-0.2);
			\draw (O)  circle (2);
			\draw[smooth,samples=300]  plot (\x,{\a*(\x)^2+\b*(\x)+\c});
			 \draw[smooth,samples=300,rotate=60] plot (\x,{\a*(\x)^2+\b*(\x)+\c});
			\draw[smooth,samples=300,rotate=-60]  plot (\x,{\a*(\x)^2+\b*(\x)+\c});
			\draw[smooth,samples=300,rotate=180] plot  (\x,{\a*(\x)^2+\b*(\x)+\c});
			 \draw[smooth,samples=300,rotate=120] plot (\x,{\a*(\x)^2+\b*(\x)+\c});
			\draw[smooth,samples=300,rotate=-120] plot (\x,{\a*(\x)^2+\b*(\x)+\c}); 
			\fill[color=gray!60]  (O)-- plot[domain=-1:1](\x,{\a*(\x)^2+\b*(\x)+\c})--cycle;
			\fill[color=gray!60,rotate=60] (O)-- plot[domain=-1:1](\x,{\a*(\x)^2+\b*(\x)+\c})--cycle; 
			\fill[color=gray!60,rotate=-60] (O)--  plot[domain=-1:1](\x,{\a*(\x)^2+\b*(\x)+\c})--cycle;
			\fill[color=gray!60,rotate=120] (O)-- plot[domain=-1:1](\x,{\a*(\x)^2+\b*(\x)+\c})--cycle; 
			\fill[color=gray!60,rotate=-120] (O)--  plot[domain=-1:1](\x,{\a*(\x)^2+\b*(\x)+\c})--cycle;
			\fill[color=gray!60,rotate=180]  (O)-- plot[domain=-1:1](\x,{\a*(\x)^2+\b*(\x)+\c})--cycle;
			\draw  ($(A)!3.2!(O)$)--($(O)!2.3!(A)$)
				 ($(B)!3.2!(O)$)--($(O)!2.3!(B)$);
			\draw[->]  (\xmin,0)--(\xmax,0) node [below]{$x$};
			 \draw[->] (0,\ymin)--(0,\ymax) node [left]{$y$};
			\node at  (0,0) [below left]{$O$};
			\foreach \p/\q in  {A/20,B/70,C/130,D/-70,E/-10,F/40}
			\fill  (\p) circle (1.0pt)
				($(\p)+(\q:3mm)$)  node{$\p$};
		\end{tikzpicture}
	 }
	\shortans{3,46} 
	\loigiai{ 
		Gắn hệ trục tọa độ như hình  vẽ.
		\begin{center}
			\begin{tikzpicture}[scale=1,font=\footnotesize,line join=round,line cap=round,>=stealth]
					\def\a{sqrt (3)/2} \def\b{0} \def\c{sqrt  (3)/2}
					 \def\xmin{-3} \def\xmax{3}
					\def\ymin{-3} \def\ymax{3} 
					 \path
						(0,0)  coordinate (O)
						 (1,1.732) coordinate (A)
						(-1,1.732) coordinate  (B)
						(-2,0) coordinate  (C)
						 (-1,-1.732) coordinate (D)
						 (1,-1.732) coordinate (E)
						 (2,0) coordinate (F);
					\draw[->] (\xmin,0)--(\xmax,0) node [below]{$x$}; 
					\draw[->]  (0,\ymin)--(0,\ymax) node [left]{$y$};
					\clip (\xmin+0.1,\ymin+0.1)  rectangle (\xmax-0.1,\ymax-0.1);
					\draw (O) circle  (2);
					\draw[smooth,samples=300] plot  (\x,{\a*(\x)^2+\b*(\x)+\c});
					\draw[smooth,samples=300,rotate=60] plot (\x,{\a*(\x)^2+\b*(\x)+\c}); 
					\draw[smooth,samples=300,rotate=-60] plot (\x,{\a*(\x)^2+\b*(\x)+\c}); 
					 \draw[smooth,samples=300,rotate=180] plot (\x,{\a*(\x)^2+\b*(\x)+\c});
					\draw[smooth,samples=300,rotate=120] plot (\x,{\a*(\x)^2+\b*(\x)+\c}); 
					\draw[smooth,samples=300,rotate=-120] plot (\x,{\a*(\x)^2+\b*(\x)+\c}); 
					\fill[color=gray!60] (O)--  plot[domain=-1:1](\x,{\a*(\x)^2+\b*(\x)+\c})--cycle;
					\fill[color=gray!60,rotate=60]  (O)-- plot[domain=-1:1](\x,{\a*(\x)^2+\b*(\x)+\c})--cycle;
					\fill[color=gray!60,rotate=-60]  (O)-- plot[domain=-1:1](\x,{\a*(\x)^2+\b*(\x)+\c})--cycle;
					 \fill[color=gray!60,rotate=120] (O)-- plot[domain=-1:1](\x,{\a*(\x)^2+\b*(\x)+\c})--cycle;
					\fill[color=gray!60,rotate=-120]  (O)-- plot[domain=-1:1](\x,{\a*(\x)^2+\b*(\x)+\c})--cycle;
					\fill[color=gray!60,rotate=180] (O)-- plot[domain=-1:1](\x,{\a*(\x)^2+\b*(\x)+\c})--cycle; 
					 \draw ($(A)!3.2!(O)$)--($(O)!2.3!(A)$)
						($(B)!3.2!(O)$)--($(O)!2.3!(B)$); 
					\draw[->] (\xmin,0)--(\xmax,0)  node [below]{$x$};
					\draw[->] (0,\ymin)--(0,\ymax) node [left]{$y$}; 
					\node at (0,0)  [below left]{$O$};
					\foreach \p/\q in  {A/20,B/70,C/130,D/-70,E/-10,F/40}
					 \fill (\p) circle (1.0pt)
						($(\p)+(\q:3mm)$) node{$\p$}; 
				\end{tikzpicture}
		\end{center}
		Ta dễ có được góc $\widehat{AOF}=\dfrac{\pi}{3}$ suy ra đường thẳng $O  A\colon y=\sqrt{3} x$ do đó điểm $A(1;\sqrt{3})$.\\
		Để tiện tính toán: Ta xét riêng diện tích tô đậm được tạo bởi cung tròn  $C_1$ với hai đoạn thẳng $OA$; $OB$.
		\begin{center}
			\begin{tikzpicture}[scale=1,font=\footnotesize,line join=round,line cap=round,>=stealth]
					\def\a{sqrt (3)/2}  \def\b{0} \def\c{sqrt (3)/2}
					 \def\xmin{-3} \def\xmax{3}
					\def\ymin{-1} \def\ymax{5} 
					\path 
						 (0,0) coordinate (O)
						(1,1.732)  coordinate (A)
						(-1,1.732) coordinate  (B)
						(0,0.866) coordinate  (I);
					\draw[color=gray!50,dashed] (\xmin,\ymin)  grid (\xmax,\ymax);
					\draw[->]  (\xmin,0)--(\xmax,0) node [below]{$x$};
					\draw[->] (0,\ymin)--(0,\ymax) node  [left]{$y$};
					\node at (0,0) [below  left]{$O$};
					\clip  (\xmin+0.1,\ymin+0.1) rectangle (\xmax-0.1,\ymax-0.1);
					\draw[smooth,samples=300]  plot (\x,{\a*(\x)^2+\b*(\x)+\c});
					\draw (0,0)--($(0,0)!3!(1,1.732)$) 
						(0,0)--($(0,0)!3!(-1,1.732)$); 
					\fill[pattern=north  east lines,opacity=0.8] (O)--(I)-- plot[domain=0:1](\x,{0.866*(\x)^2+0.866})--cycle;
					\node  at (A)[below right]{$\left(1; \sqrt{3}\right)$};
					\foreach \p/\q in {A/150,B/190,I/100} 
					 \fill (\p) circle (1.0pt)
						($(\p)+(\q:2.5mm)$) node{$\p$}; 
				\end{tikzpicture}
		\end{center}
		Dễ thấy được đường cong $C_1$ là đồ thị hàm số parabol đối xứng qua trục $O y$. Do đó hàm số parapol đó có dạng: $y=ax^2+b$ và  đồ thị hàm số này đi qua điểm $A$, $B$.\\
		Ta có $\heva{&a+b=\sqrt{3}\\&  2 a=\sqrt{3}} \Rightarrow\heva{&a=\dfrac{\sqrt{3}}{2}\\& b=\dfrac{\sqrt{3}}{2}.}$\\
		Vậy $(C_1)\colon y=\dfrac{\sqrt{3}}{2}  x^2+\dfrac{\sqrt{3}}{2}$.\\
		Khi đó diện tích tô đậm được tạo bởi cung tròn $C_1$ với hai đoạn thẳng $OA$; $OB$  là
		$$S_1=2 \displaystyle\int_0^1\left(\dfrac{\sqrt{3}}{2} x^2-\sqrt{3} x+\dfrac{\sqrt{3}}{2}\right)  \mathrm{\,d}x.$$
		Suy ra diện tích phần được bao bởi $6$ cung tròn $C_1$ đến  $C_6$ là
		$$S=6 S_1=12  \displaystyle\int_0^1\left(\dfrac{\sqrt{3}}{2}x^2-\sqrt{3}x +\dfrac{\sqrt{3}}{2}\right) \mathrm{\,d}x=2 \sqrt{3} \approx 3{,}46.$$
	} 
\end{bt}

% nguồn: Phan2_Lop_12_Dai_so_Chuong_4
\begin{bt}%[2D4V3-5]%[ma:04cf80fe]
	 \immini{
	Một vật trang trí có dạng một  khối tròn xoay được tạo thành khi quay miền $(R)$ (phần được tô màu trong hình vẽ bên) quanh trục $AB$. Miền ($R$) được giới hạn bởi các cạnh $AB$, $AD$ của hình vuông $ABCD$ và các cung phần tư của các đường tròn bán kính bằng $1$ cm với tâm lần lượt là trung điểm của các cạnh $AD$, $AB$. Tính thể tích của vật trang trí đó, làm tròn kết quả đến hàng phần mười.
	 }{
	\begin{tikzpicture}[scale=0.8,font=\footnotesize,line join=round,line cap=round,>=stealth]
			 \coordinate (A) at (0,0);
			\coordinate  (B) at (4,0);
			\coordinate  (C) at (4,4);
			\coordinate  (D) at (0,4);
			\draw (A)--(B)--(C)--(D)--cycle; 
			\draw[pattern=north east  lines] (D) arc (90:0:2) arc (90:0:2)--(B)--(A);
			\path node at  (A) [above left]{$A$}
				node at (B) [below] {$B$} 
				node at (C) [above  right] {$C$}
				node at (D) [above left]  {$D$};
		\end{tikzpicture}
	 }
	\shortans{12,3} 
	 \loigiai{
		Chọn trục $Ox$ chứa điểm $B$, trục $Oy$ chứa điểm $D$,  và gốc tọa độ $O$ trùng điểm $A$ (như hình vẽ).
		\begin{center}
			\begin{tikzpicture}[scale=0.8,font=\footnotesize,line join=round,line cap=round,>=stealth]
					\coordinate (A)  at (0,0);
					 \coordinate (B) at (4,0);
					\coordinate (C) at (4,4); 
					 \coordinate (D) at (0,4);
					 \coordinate (F) at (2,0);
					\coordinate (E) at (0,2); 
					\draw (A)--(B)--(C)--(D)--cycle; 
					\draw (D) arc  (90:-90:2)
						(A)  arc (180:0:2);
					\path node at (A) [above left]{$A\equiv  O$}
						node at  (B) [below] {$B$}
						 node at (C) [above right] {$C$}
						node at (D) [left] {$D$}; 
					\draw[->] (B)--(6,0)  node[below right]{$x$};
					\draw[->]  (D)--(0,5.5) node[left]{$y$};
					\node at (2,0)[below]{$F$}; 
					\fill (F)  circle (1.0pt) node[above]{$1$};
					\node at (0,2)[left]{$E$}; 
					 \fill (E) circle (1.0pt) node[right]{$1$}
						(B)  circle (1.0pt)
						(C) circle (1.0pt) 
						 (D) circle (1.0pt)
						(A)  circle (1.0pt);
					\node  at (B)[above right]{$2$};
					\node at (D)[above right]{$2$}; 
				\end{tikzpicture}
		\end{center}
		Gọi $E$, $F$ lần lượt là trung điểm của $AD$, $AB$. Khi  đó $E(0; 1)$, $F(1; 0)$.\\
		+ Phương trình đường tròn có tâm $E(0; 1)$ và đường kính $AD=2$  là $x^2+(y-1)^2=1$.\\
		Suy ra phương trình cung trên của  đường tròn tâm $E$ là $y=1+\sqrt{1-x^2}$.\\
		+ Phương trình đường  tròn có tâm $F(1; 0)$ và đường kính $AB=2$ là $(x-1)^2+y^2=1$.\\
		Suy ra  phương trình cung trên của đường tròn tâm $F$ là $y=\sqrt{1-(x-1)^2}$.\\
		Vậy, thể tích  vật trang trí là
	$$V=\pi \displaystyle\int_0^1\left(1+\sqrt{1-x^2}\right)^2 \mathrm{\,d}x+\pi \displaystyle\int_1^2\left(1-(x-1)^2\right) \mathrm{\,d}x \approx  12{,}3\left(\text{cm}^3\right).$$}
\end{bt}

% nguồn: Phan2_Lop_12_Dai_so_Chuong_4
\begin{bt}%[2D4V3-5]%[ma:f2e5ed17]
	Khi sử dụng phần mềm mô  phỏng để thiết kế một chậu cây, người ta quay hình phẳng giới hạn bởi đồ thị hàm số $y= \sqrt{x} + 2$, trục hoành và hai đường thẳng $x=0$, $x=4$ quanh trục hoành. Biết đơn vị trên các trục tọa độ là đề-xi-mét. Thể tích của chậu cây (kết quả làm tròn đến hàng phần mười) bằng bao nhiêu đề-xi-mét khối?
	 \shortans{142,4}
	\loigiai{ 
	Thể tích  của chậu cây là
	\begin{flalign*}
		V &=  \pi \displaystyle \int_{0}^{4} \left(\sqrt{x} + 2\right)^2 \mathrm{\,d}x = \int_{0}^{4} \left(x + 4\sqrt{x} + 4\right) \mathrm{\,d}x &\\
		&= \pi  \left(\dfrac{x^2}{2} + \dfrac{8}{3}x\sqrt{x} + 4x\right)\bigg|_{0}^{4} = \dfrac{136\pi}{3} \approx 142{,}4 \left(\text{dm}^3\right).
	\end{flalign*}
	} 
\end{bt}

% nguồn: Phan2_Lop_12_Dai_so_Chuong_4
\begin{bt}%[2D4V3-5]%[ma:af264116]
	Trang trí một sân hình chữ nhật kích thướt $28$\,m$\times 16$\,m, trong đó hai Parabol $(P_1)$ đối xứng với $(P_2)$ qua đường thẳng đi qua hai trung điểm của chiều dài sân (hình vẽ), khoảng cách giữa hai đỉnh Parabol bằng $4$\,m. Chi phí trang trí cho phần hoa văn là $180$ ngàn đồng trên một mét vuông, phần trắng là $160$\,ngàn đồng trên một mét vuông. Tổng chi phí trang trí cho sân là bao nhiêu triệu đồng? (làm tròn kết quả đến hàng  phần mười)
	\begin{center}
		\begin{tikzpicture}[line join=round,line cap=round,>=stealth,scale=0.4]
				\tikzset{every  node/.style={scale=0.9}}
				\begin{scope}[rotate=90]
					\draw  (-8,-14) rectangle (8,14);
					\fill[pattern=bricks]plot[samples=200,domain=-8:8,smooth,variable=\x] (\x,{0.25*(\x)^2-2})--plot[samples=200,domain=8:-8,smooth,variable=\x] (\x,{-0.25*(\x)^2+2})--cycle; 
					\draw  plot[samples=200,domain=-8:8,smooth,variable=\x] (\x,{0.25*(\x)^2-2}) node[below right=0.3]{$(P_1)$};
					 \draw plot[samples=200,domain=-8:8,smooth,variable=\x] (\x,{-0.25*(\x)^2+2}) node[below left=0.3]{$(P_2)$};
				\end{scope}
			\end{tikzpicture}
	\end{center}
	\shortans{74,4} 
	\loigiai{ 
		\begin{center}
			\begin{tikzpicture}[line join=round,line cap=round,>=stealth,scale=0.4]
					 \tikzset{every node/.style={scale=0.9}}
					\begin{scope}
						\clip (-9,-14.5) rectangle  (9,14.5);
						 \draw (-8,-14) rectangle (8,14);
						\fill[gray!40]plot[samples=200,domain=-8:8,smooth,variable=\x] (\x,{0.25*(\x)^2-2})--plot[samples=200,domain=8:-8,smooth,variable=\x] (\x,{-0.25*(\x)^2+2})--cycle; 
						\draw plot[samples=200,domain=-8:8,smooth,variable=\x] (\x,{0.25*(\x)^2-2}) node[above]{$f  (x)$};
						\draw  plot[samples=200,domain=-8:8,smooth,variable=\x] (\x,{-0.25*(\x)^2+2}) node[below]{$g (x)$};
					\end{scope}
					 \draw[->] (-9.1,0)--(9.1,0) node[below] {$x$};
					\draw[->] (0,-14.6)--(0,15)  node[left] {$y$};
					\draw (0,0)  node [below left] {$O$};
					\draw[thin]  (-8,1pt)--(-8,-1pt) node [below left] {$-8$};
					\draw[thin]  (8,1pt)--(8,-1pt) node [below right] {$8$};
					\foreach \y/\ny in  {-14/-14,-2/-2}
					\draw[thin] (1pt,\y)--(-1pt,\y) node [below  left] {$\ny$};
					 \foreach \y/\ny in {2/2,14/14}
					\draw[thin] (1pt,\y)--(-1pt,\y) node [above left]  {$\ny$};
				\end{tikzpicture}
		\end{center}
		Đặt  hình chữ nhật vào hệ tọa độ $Oxy$ sao cho tâm hình chữ nhật trùng với góc tọa độ, chiều dài của hình chữ nhật vuông góc với trục $Ox$.\\
		Theo đề bài ta có Parabol $(P_1)$ đi  qua các điểm $A(0;-2)$, $B(-8;14)$, $C(8;14)$.\\
		Từ đây ta  có hệ phương trình \[\heva{&c=-2\\&64a-8b=16\\&64a+8b=16}\Leftrightarrow \heva{&c=-2\\&b=0\\&a=\dfrac{1}{4}.}\]
		Suy ra  $(P_1)\colon y=\dfrac{1}{4}x^2-2$.\\
		Suy ra $(P_2)\colon y=-\dfrac{1}{4}x^2+2$.\\ 
		Gọi $S_1$, $S_2$ lần lượt là diện tích phần hoa văn và diện tích  phần trắng.\\
		Ta có 
		\begin{itemize}
			\item $S_1=2\displaystyle\int_{-8}^{-2\sqrt 2}\left|\dfrac{1}{2}x^2-4\right|\mathrm{\,d}x+\displaystyle\int_{-2\sqrt 2}^{2 \sqrt  2}\left|\dfrac{1}{2}x^2-4\right|\mathrm{\,d}x$.
			\item  $S_2=16\cdot 28-S_1=448-S_1$.
		\end{itemize}
		Tổng  chi phí trang trí là $180\,000S_1+160\,000S_2=74{,}4$ (triệu đồng).
	} 
\end{bt}

% nguồn: Phan2_Lop_12_Dai_so_Chuong_4
\begin{bt}%[2D4V3-5]%[ma:3a6d9c84]
	Cho đường tròn tâm $O$, bán kính bằng $2$. Trên đường tròn, lấy $6$ điểm chia đều đường tròn,  lần lượt là $A$, $B$, $C$, $D$, $E$, $F$. Vẽ parabol $C_1$ tiếp xúc với hai đoạn thẳng $OA$ và $OB$, đi qua hai điểm $A$ và $B$. Tương tự, vẽ parabol $C_2$ tiếp xúc với hai đoạn thẳng $OB$ và $OC$, đi qua hai điểm $B$ và $C$. Tiếp tục bằng cách tương tự, ta vẽ các parabol $C_3$, $C_4$, $C_5$, $C_6$ tiếp xúc với các cặp đoạn thẳng liên tiếp tạo bởi các điểm trên.
	Hỏi: Diện tích phần được bao bởi $6$ parabol $C_1$ đến $C_6$ là bao nhiêu? (làm tròn kết quả đến hàng phần trăm)\shortans{3,46}\begin{center}
	\begin{tikzpicture}[line cap=butt,line join=miter,>=stealth]
			\tikzset{declare function={a=sqrt (3)/2;  b=sqrt (3)/2; f (\x)=a*(\x)^2+b;}}
			\begin{scope}
				 \clip (-3,-3) rectangle (3,3);
				 \foreach[count=\i] \goc/\t in {0/A,60/B,120/C,180/D,240/E,300/F}{
					\begin{scope}[rotate=\goc]
					\fill[cyan!35,samples=100] (0,0)--plot[domain=-1:1] (\x,{f (\x)})--cycle; 
					\draw[samples=100] plot[domain=-2:2] (\x,{f (\x)}); 
					 \draw (240:5)--(60:5);
					\end{scope}
					\path ({68+\goc}:2.2)  node{$\t$}
					 ({70+\goc}:2.8) node{$C_\i$};
				} 
			\end{scope}
			 \draw[->] (-3,0)--(3,0)
				 node[shift={(-100:7pt)},font=\normalsize]{$x$};
			 \draw[->] (0,-3)--(0,3)
				node[shift={(180:7pt)},font=\normalsize]{$y$}; 
			\draw (5pt,0)  |- (0,5pt)
				 (0,0) node[shift={(220:9pt)},font=\normalsize]{$O$};
			\draw (0,0) circle (2); 
		\end{tikzpicture}
\end{center}
\loigiai{ 
	Gắn hệ  trục tọa độ như hình vẽ.
	\begin{center}
		\begin{tikzpicture}[line cap=butt,line join=miter,>=stealth]
				\tikzset{declare function={a=sqrt (3)/2; b=sqrt (3)/2; c=sqrt (3);  f (\x)=a*(\x)^2+b;}}
				\begin{scope}
					\clip (-3,-1)  rectangle (3,3);
					\begin{scope}
						\fill[cyan!35,samples=100] (0,0)--plot[domain=-1:1]  (\x,{f (\x)})--cycle;
						\draw[samples=100] plot[domain=-2:2] (\x,{f  (\x)});
						\draw  (240:0)--(60:5);
						\draw (300:0)--(120:5); 
					\end{scope}
				\end{scope}
				 \draw[->] (-3,0)--(3,0)
					node[shift={(-100:7pt)},font=\normalsize]{$x$}; 
				\draw[->]  (0,-1)--(0,3)
					 node[shift={(180:7pt)},font=\normalsize]{$y$};
				\draw (5pt,0)  |- (0,5pt)
					(0,0)  node[shift={(220:9pt)},font=\normalsize]{$O$};
				\draw (1,c) node[right]{$A$}circle (1pt) 
					 (-1,c) node[right]{$B$}circle (1pt);
			\end{tikzpicture}
	\end{center}
	Ta dễ có được góc $\widehat{AOF}=\dfrac{\pi}{3}$ suy ra đường thẳng $OA\colon y=\sqrt{3}x$ do đó  điểm $A(1;\sqrt{3})$.\\
	Để tiện tính toán:  Ta xét riêng diện tích tô đậm được tạo bởi parabol $C_1$ với hai đoạn thẳng $OA$; $OB$.\\
	Dễ thấy được parabol $C_1$ là đồ thị hàm số parabol đối xứng qua trục $Oy$.  Do đó hàm số parapol đó có dạng $y=ax^2+b$ và đồ thị hàm số này đi qua điểm $A$ và nhận $OA\colon y=\sqrt3x$ làm tiếp tuyến nên $y'(1)=\sqrt3\Rightarrow 2a=\sqrt3$.\\
	Ta có $\left\{\begin{aligned}
	&a+b=\sqrt{3} \\ 
	 &2a=\sqrt{3}
	\end{aligned}
	\Rightarrow \left\{\begin{aligned}
	&a=\dfrac{\sqrt{3}}{2}  \\
	 & b=\dfrac{\sqrt{3}}{2}.
	\end{aligned}
	\right.\right.$ \\ 
	Vậy  $\left(C_1\right)\colon y=\dfrac{\sqrt{3}}{2}x^2+\dfrac{\sqrt{3}}{2}$.\\
	Khi đó diện tích tô đậm được tạo bởi cung tròn $C_1$ với hai đoạn  thẳng $OA$; $OB$ là\\
	\[S_1=2\displaystyle\int^1_0\left(\dfrac{\sqrt{3}}{2}x^2-\sqrt{3}x+\dfrac{\sqrt{3}}{2}\right)\, \mathrm{d}x. 
	\] 
	Suy ra diện tích  phần được bao bởi $6$ cung tròn $C_1$ đến $C_6$ là\\
	\[S=6S_1=12\displaystyle\int^1_0\left(\dfrac{\sqrt{3}}{2}x^2-\sqrt{3}x+\dfrac{\sqrt{3}}{2}\right)\,\mathrm{d}x=2\sqrt{3}\approx  3{,}46.
	\] 
	} 
\end{bt}

% nguồn: Phan2_Lop_12_Dai_so_Chuong_4
\begin{bt}%[2D4V3-5]%[ma:9f8bd13e]
	Một vật trang trí có dạng một khối tròn xoay được tạo thành khi quay miền $(R)$ (phần màu xám trong hình vẽ) quanh trục $M N$. Biết rằng $ABCD$ là hình chữ nhật với $AB=6$ cm, $AD=10$ cm và $M$, $N$ lần lượt là trung điểm của $AB$, $CD$. Hai đường cong là đường Elip có hình chữ nhật cơ sở là $ABCD$ và đường tròn tiếp xúc với hai cạnh $AD$ và $BC$ (tham khảo hình vẽ). Tính thể tích của vật trang trí đó (kết quả làm tròn đến hàng đơn vị theo đơn  vị cm$^{3}$).
	\begin{center}
		\begin{tikzpicture}[scale=1,font=\footnotesize,line join=round,line cap=round,>=stealth]
				\fill[gray!50,even odd rule] (0,0) ellipse (5 and 3) 
					(0,0) circle  (3);
				\draw (-5,-3)  rectangle (5,3);
				\draw (0,0) ellipse (5 and 3); 
				\draw (0,0) circle  (3);
				\foreach \x/\label in  {5/N}{
					\draw (\x,2pt) -- (\x,-2pt); 
					\node[below right] at (\x,0) {$\label$}; 
				} 
				\foreach \x/\label  in {-5/M}{
					\draw  (\x,2pt) -- (\x,-2pt);
					\node[below left] at  (\x,0) {$\label$};
				} 
				\node[above left]  at (-5,3) {$A$};
				\node[below left] at (-5,-3) {$B$}; 
				\node[below right]  at (5,-3) {$C$};
				\node[above right] at  (5,3) {$D$};
			\end{tikzpicture}
	\end{center}
	 \shortans{75}
	\loigiai{ 
		Gắn vào hình chữ nhật $ABCD$ hệ trục tọa độ $Oxy$ với gốc $O$ trùng với giao hai đường chéo, các điểm $M$,  $N$ nằm trên $Ox$ (tham khảo hình vẽ).
		\begin{center}
			\begin{tikzpicture}[scale=1,font=\footnotesize,line join=round,line cap=round,>=stealth]
					\fill[gray!50,even odd rule]  (0,0) ellipse (5 and 3)
						(0,0)  circle (3);
					 \draw[->] (-6,0) -- (6,0) node[anchor=west] {$x$};
					\draw[->] (0,-4) --  (0,4) node[anchor=south] {$y$};
					 \draw[red] (-5,-3) rectangle (5,3);
					\draw (0,0) ellipse (5 and 3); 
					\draw  (0,0) circle (3);
					\foreach \x/\label in {5/N}{ 
						\draw (\x,2pt) --  (\x,-2pt);
						\node[below right] at (\x,0) {$\label$}; 
					} 
					\foreach \x/\label  in {-5/M}{
						\draw  (\x,2pt) -- (\x,-2pt);
						\node[below left] at (\x,0) {$\label$}; 
					 }
					\foreach \x  in {-3,3}{
						\draw (\x,2pt) --  (\x,-2pt);
						\node[above right] at (\x,0)  {$\x$};
					 }
					 \foreach \y in {-3,3}{
						\draw  (2pt,\y) -- (-2pt,\y);
						\node[above left] at (0,\y)  {$\y$};
					} 
					\node[above left] at (-5,3)  {$A$};
					\node[below left] at (-5,-3)  {$B$};
					\node[below right] at  (5,-3) {$C$};
					 \node[above right] at (5,3) {$D$};
					\node[above  right] at (0,0) {$O$};
				\end{tikzpicture}
		\end{center}
		Do $M  N$ là trục đối xứng của miền $(R)$ nên khối tròn xoay được tạo thành khi quay miền $(R)$ quanh trục $MN$ trùng với khối tròn xoay được tạo thành khi quay phần $(R)$ nằm phía trên $MN$ quanh $MN$.\\
		Gọi $V_{1}$, $V_{2}$, $V$ tương ứng là thể tích của các khối tròn xoay được tạo thành khi quay Elips ($E$), hình  tròn $(C)$ và miền $(R)$ quanh $MN$ (hay quanh $Ox$), ta có $V=V_{1}-V_{2}$.\\
		Elip $(E)$  có các bán trục $a=\dfrac{1}{2} AD=5$, $b=\dfrac{1}{2} A B=3$ nên có phương trình $\dfrac{x^2}{25}+\dfrac{y^2}{9}=1$.\\
		Ta có $V_{1}=\pi \displaystyle\int_{-5}^{5} 9\left(1-\dfrac{x^2}{25}\right) \mathrm{\,d} x=\left.\pi\left(9 x-\dfrac{9}{25} x^{3}\right)\right|_{-5}^{5}=60 \pi \left(cm^{3}\right)$.\\ 
		Hình tròn $(C)$ có tâm tại $O$, bán kính $R=3$ khi quay quanh $M N$ tạo thành khối cầu tâm $O$, bán kính  bằng $3$. Do đó $V_{2}=\dfrac{4}{3} \pi \cdot 3^{3}=36 \pi\left(\mathrm{~cm}^{3}\right)$.\\
		Vậy thể tích của vật trang trí đó là: $V=V_{1}-V_{2}=24  \pi \simeq 75\left(\mathrm{~cm}^{3}\right)$.\\
	} 
\end{bt}

% nguồn: Phan2_Lop_12_Dai_so_Chuong_4
\begin{bt}%[2D4V3-5]%[ma:62c84419]
	\immini{Một vật có kích thước và hình dáng như hình vẽ dưới đây. Đáy là hình tròn giới hạn bởi đường tròn $x^2+y^2=16$ , cắt vật bởi các mặt phẳng  vuông góc với trục Ox ta được thiết diện là tam giác đều. Khi đó thể tích của vật thể có dạng $\dfrac{a\sqrt 3}{b}$ với $\dfrac{a}{b}$ là phân số tối giản. Tính $S=a+b$.}
	{\tdplotsetmaincoords{70}{40} 
	\begin{tikzpicture}[>=stealth,tdplot_main_coords]
			\draw[canvas is yx plane at z=0]  (0,0) circle[radius=2cm];
			 \draw[->] (-3,0,0)--(3,0,0) node[below]{$x$};
			\draw[->] (0,-3,0)--(0,3,0) node[right]{$y$}; 
			\draw[->]  (0,0,-1)--(0,0,4) node[left]{$z$};
			\foreach  \i in {-12,-11,...,12}{
				 \pgfmathsetmacro\x{\i/6}
				\pgfmathsetmacro\y{sqrt  (4-(\x)^2)}
				\pgfmathsetmacro\z{\y*sqrt  (3)}
				\coordinate (A\i) at (\x,\y,0); 
				 \coordinate (B\i) at (\x,-\y,0);
				\coordinate (C\i)  at (\x,0,\z);
				\draw  (A\i)--(B\i)--(C\i)--cycle;}
			\draw plot  [smooth,tension=.5] coordinates{(C-12)
				(C-11) 
				(C-10) 
				(C-9) 
				(C-8) 
				(C-7) 
				(C-6) 
				(C-5) 
				 (C-4)
				(C-3) 
				(C-2) 
				 (C-1)
				(C0) 
				 (C1)
				 (C2)
				 (C3)
				 (C4)
				(C5) 
				(C6) 
				(C7) 
				(C8) 
				(C9) 
				(C10) 
				(C11) 
				(C12)}; 
		\end{tikzpicture}}
	 \shortans{259}
	\loigiai{ 
		Ta có  $x^2+y^2=16\Leftrightarrow y=\pm\sqrt{16-x^2}$.\\
		Suy ra thiết diện  là tam giác đều có cạnh bằng $2\sqrt{16-x^2}$.\\
		Do đó diện tích thiết  diện là $S(x)=\left(2\sqrt{16-x^2}\right)^2\cdot\dfrac{\sqrt 3}{4}=\sqrt 3\left(16-x^2\right)$.\\
		Thể  tích cần tìm là $V=\displaystyle\int_{-4}^4S(x)\,\mathrm{d}x=\sqrt 3\displaystyle\int_{-4}^4\left(16-x^2\right)\mathrm{d}x=\dfrac{256\sqrt 3}{3}$.\\
		 Suy ra $a=256$, $b=3$. Do đó $S=a+b=259$.
	} 
\end{bt}

% nguồn: Phan2_Lop_12_Dai_so_Chuong_4
\begin{bt}%[2D4V3-5]%[ma:72ba0e3e]
	\immini[thm]{Trong các làng nghề chế tạo đá mỹ nghệ, người ta làm được nhiều đồ  trang sức bằng đá rất đẹp. Một người thợ thủ công chế tạo vòng đeo tay hình tròn bằng đá xanh, biết rằng bán kính vòng ngoài của vòng là $7$\,(cm), bán kính vòng trong là $5$\,(cm) (hình vẽ). Tính thể tích của chiếc vòng (theo đơn vị đo là\,cm$^3$ và làm tròn đến hàng đơn vị).}{

	\begin{tikzpicture}
			\begin{axis}[
					 axis equal image,
					hide  axis,
					z buffer  = sort,
					view =  {122}{30},
					scale =  1]
				\addplot3[ 
					surf, 
					 shader = interp,
					samples  = 30,
					samples y = 50, 
					domain  = 0:2*pi,
					domain  y = 0:2*pi,
					variable  = \u,
					variable y = \v,]( 
					{(3+cos (deg  (\u)))*cos (deg (\v))},
					 {(3+cos (deg (\u)))*sin (deg (\v))},
					{sin  (deg (\u))}
					); 
			\end{axis}
		\end{tikzpicture}}
	 \shortans{118}
	\loigiai{ 
	\immini{Gắn hệ trục tọa độ như  hình vẽ.\\
		Gốc tọa độ là tâm của chiếc vòng, trục $Ox$ là trục đối xứng của vòng, trục  $Oy$ đi qua tâm một thiết diện cắt ngang của vòng.\\
		Khi đó phương trình  đường tròn thiết diện là
		\[\left(C\right)\colon x^2+\left(y-6\right)^2 = 1  \Leftrightarrow y = 6 \pm \sqrt{1-x^2}.\]
		Vòng được tạo ra bằng cách quay đường tròn $\left(C\right)$ quanh trục  $Ox$.\\
		 Thể tích của chiếc vòng là
		\[V=\pi\left[\displaystyle\int_{-1}^{1}  \left(6+\sqrt{1-x^2}\right)^2\,\mathrm{d}x- \displaystyle\int_{-1}^{1}\left(6-\sqrt{1-x^2}\right)^2\, \mathrm{d}x\right]=\pi\displaystyle\int_{-1}^{1} 24\sqrt{1-x^2}\,\mathrm{d}x=12\pi^2\approx 118\,\,\left(\text{cm}^3\right).\]}
		{ 
		\begin{tikzpicture}[scale=.6,line cap=butt,line join=miter,>=stealth]
				\draw[gray!55]  (-3,-1) grid (5,8);
				 \draw[->] (-3,0)--(5,0)
					node[shift={(-100:7pt)},font=\normalsize]{$x$}; 
				\draw[->] (0,-1)--(0,8) 
					node[shift={(180:7pt)},font=\normalsize]{$y$}; 
				\draw (5pt,0) |- (0,5pt) 
					(0,0) node[shift={(225:9pt)},font=\normalsize]{$O$}; 
				\draw (0,6) circle (1cm); 
				\foreach \x in {-3,-2,-1,1,2,3,4}{ 
					\draw (\x,1.5pt)--(\x,-1.5pt)  node[font=\scriptsize,shift={(-90:5pt)}]{\x};
				 }
				\foreach \y in {-1,1,2,3,4,5,6}{ 
					\draw  (1.5pt,\y)--(-1.5pt,\y) node[font=\scriptsize,shift={(180:5pt)}]{\y};
				 }
			\end{tikzpicture}
		 }
	 }
\end{bt}

% nguồn: Phan2_Lop_12_Dai_so_Chuong_4
\begin{bt}%[2D4V3-5]%[ma:427f2cc4]
	Bình chứa nhiên liệu của một chiếc máy bay phản lực được mô tả trong không gian là một khối tròn xoay thu được khi quay hình phẳng giới hạn bởi đồ thị hàm số $y=\dfrac{1}{8}x^2\sqrt{2-x}$ với $(0\le x\le 2)$ và trục $Ox$ quanh trục $Ox$ ($x$ tính theo đơn vị  m). Thể tích của bình chứa bằng $\dfrac{a\pi}{b}$ m$^3$ với $\dfrac{a}{b}$ là phân số tối giản. Khi đó $a+b$ bằng
	\shortans{31} 
	 \loigiai{
		Thể  tích của bình chứa là
		$$V=\displaystyle\int_0^2\pi\left(\dfrac{1}{8}x^2\sqrt{2-x}\right)^2\mathrm{\,d}x=\dfrac{\pi}{30}\ \left(\mathrm{m}^3\right).$$ 
		Do đó, $a=1$ và $b=30$  nên $a+b=31$.
	 }
\end{bt}

% nguồn: Phan2_Lop_12_Dai_so_Chuong_4
\begin{bt}%[2D4V3-5]%[ma:27bbe339]
	\immini{Một lều cắm trại có dạng như hình vẽ dưới, khung lều được tạo thành  từ hai parabol giống nhau có chung đỉnh $O$ và thuộc hai mặt phẳng vuông góc nhau (một parabol đi qua $A$, $O$, $C$ và một parabol đi qua $B$, $D$, $O$). Bốn chân tạo thành hình vuông $ABCD$ có cạnh là $2\sqrt{2}$\,(m), chiều cao tính từ đỉnh lều là $2$\,(m). Biết mặt cắt của lều khi cắt bởi một mặt phẳng song song với mặt phẳng $(ABCD)$ luôn là một hình vuông. Tính thể tích của lều (đơn vị là\,m$^3$).}
	{\begin{tikzpicture}[scale=1,>=stealth,line join=round,line cap=round,font=\footnotesize]
		 \def\g{-45}
		 \def\r{.95}
		\def\xmax{3} 
		 \def\ymax{2}
		\def\zmax{3.5} 
		\pgfmathsetmacro\ra{sqrt  (1-(\r)^2)}
		\pgfmathsetmacro\a{cos  (\g)}
		\pgfmathsetmacro\b{sin (\g)} 
		 \pgfmathsetmacro\ax{\ra*cos (\g)/sqrt ((\ra*\a)^2+(\b)^2)}
		 \pgfmathsetmacro\bx{\ra*sin (\g)/sqrt ((\ra*\a)^2+(\b)^2)}
		 \def\f{-.5*(\x)^2+2*\r}
		 \path
			 (0,0) coordinate (O)
			 +({\ax},{\bx}) coordinate (i)
			 +({\bx/\ra},{-\ax*\ra}) coordinate (k)
			+(0,{\r}) coordinate (j) 
			+($2*(j)$) coordinate  (P);
		\path 
			(0,0)  coordinate (O)
			[samples at={2}] plot ({\ax*\x},{\bx*\x+\r*\f}) coordinate  (A)
			[samples at={-2}] plot ({\ax*\x},{\bx*\x+\r*\f}) coordinate (C) 
			[samples at={2}]  plot ({\bx/\ra*\x},{-\ax*\ra*\x+\r*\f}) coordinate (D)
			 [samples at={-2}] plot ({\bx/\ra*\x},{-\ax*\ra*\x+\r*\f}) coordinate (B);
		\draw [smooth,samples=100]  plot [domain=-.65:2] ({\ax*\x},{\bx*\x+\r*\f})
			plot [domain=2:-2] ({\bx/\ra*\x},{-\ax*\ra*\x+\r*\f}) 
			 (D)--(A)--(B);
		\draw [dash pattern=on 1pt off 2pt,samples=100] plot  [domain=-.65:-2] ({\ax*\x},{\bx*\x+\r*\f});
		\draw [dash pattern=on 1pt off 2pt] (D)--(C)--(B); 
		\foreach \x/\y in  {A/-90,B/0,C/180,D/180}{
			\draw (\x) circle (.5pt)  ++(\y:0.2) node {$\x$};
		 }
		 \draw (P) circle (1pt) node [above] {$O$};
	\end{tikzpicture}
	} 
	 \shortans{8}
	\loigiai{ 
	\immini{Ta có  $AB=2\sqrt{2}\Rightarrow AC=4$.\\
		Chọn  hệ trục $Oxy$ như hình vẽ.\\
		Khi  đó, mặt cắt tại $x=t$ là hình vuông có diện tích $S(t) =\dfrac{1}{2}\cdot\left(2\sqrt{2t}\right)^2=4t$.\\
		 Vậy thể tích của lều là
		 $V=\displaystyle\int_0^2 S(t)\,\mathrm{d}t
		=\displaystyle\int_0^2  4t\,\mathrm{d}t=8$ (m$^3$).}{
		\begin{tikzpicture}[scale=1.5,>=stealth,line join=round,line cap=round,font=\footnotesize]
				 \def\g{-45}
				\def\r{.95} 
				 \def\xmax{3}
				\def\ymax{3} 
				\def\zmax{2.5} 
				 \pgfmathsetmacro\ra{sqrt (1-(\r)^2)}
				\pgfmathsetmacro\a{cos (\g)} 
				 \pgfmathsetmacro\b{sin (\g)}
				 \pgfmathsetmacro\ax{\ra*cos (\g)/sqrt ((\ra*\a)^2+(\b)^2)}
				 \pgfmathsetmacro\bx{\ra*sin (\g)/sqrt ((\ra*\a)^2+(\b)^2)}
				 \def\f{-.5*(\x)^2+2*\r}
				 \path
					(0,0) coordinate  (O)
					 +({\ax},{\bx}) coordinate (i)
					+({\bx/\ra},{-\ax*\ra}) coordinate  (k)
					+(0,{\r}) coordinate (j) 
					+($2*(j)$)  coordinate (P)
					 [samples at={2}] plot ({\ax*\x},{\bx*\x+\r*\f}) coordinate (A)
					 [samples at={-2}] plot ({\ax*\x},{\bx*\x+\r*\f}) coordinate (C)
					[samples at={2}]  plot ({\bx/\ra*\x},{-\ax*\ra*\x+\r*\f}) coordinate (D)
					[samples at={-2}] plot  ({\bx/\ra*\x},{-\ax*\ra*\x+\r*\f}) coordinate (B)
					 ($(A)!.5!(C)$) coordinate (I);
				\foreach \i in {1.3}{ 
					\fill [samples at={\i},cyan!90!white!60]  plot ({\bx/\ra*\x},{-\ax*\ra*\x+\r*\f})--plot ({\ax*\x},{\bx*\x+\r*\f})-- plot ({-\bx/\ra*\x},{\ax*\ra*\x+\r*\f})
					--plot ({-\ax*\x},{-\bx*\x+\r*\f})--cycle; 
				} 
				\foreach \i in {1.3}{ 
					\draw [samples at={\i},gray] plot ({\bx/\ra*\x},{-\ax*\ra*\x+\r*\f}) coordinate  (D')--plot ({\ax*\x},{\bx*\x+\r*\f})--plot ({-\bx/\ra*\x},{\ax*\ra*\x+\r*\f}) coordinate (B');
					\draw [samples at={\i},gray,dash pattern=on 1pt off 2pt] plot ({-\bx/\ra*\x},{\ax*\ra*\x+\r*\f})--plot ({-\ax*\x},{-\bx*\x+\r*\f})--plot  ({\bx/\ra*\x},{-\ax*\ra*\x+\r*\f});
					\path ($(B')!.5!(D')$) coordinate (I'); 
				} 
				 \fill [cyan!90!white!60] (A)--(B)--(C)--(D)--cycle;
				\draw  [smooth,samples=100]
					plot [domain=-.8:2]  ({\ax*\x},{\bx*\x+\r*\f})
					 plot [domain=2:-2] ({\bx/\ra*\x},{-\ax*\ra*\x+\r*\f})
					 (D)--(A)--(B);
				\draw [dash pattern=on 1pt off 2pt,samples=100,gray] plot [domain=-.8:-2] ({\ax*\x},{\bx*\x+\r*\f}); 
				\draw [dash pattern=on 1pt off 2pt,gray] (D)--(C)--(B)--cycle (A)--(C) 
					(D')--(B')--++($(P)-(I')$) 
					 (B)--++($(P)-(I)$);
				\draw [->]  ($(P)+1.25*(D)-1.25*(O)$)--($(P)+1.25*(B)-1.25*(O)$) node [above] {$y$};
				\draw [->,dash pattern=on  1pt off 2pt] (P)--++($-.8*\ymax*(j)$) node [left] {$x$};
				\draw [->] ($(P)+(C)-(O)$)--++($(A)-(C)$) node  [above right] {$z$};
				\foreach \x/\y in  {A/-90,B/0,C/180,D/180}{
					\draw (\x) circle (.5pt) ++(\y:0.2) node  {$\x$};
				} 
				\draw (P) circle (.5pt)  node [above] {$O$}
					($(B')!.5!(D')$)  circle (.5pt) node [above left] {$t$}
					(B')++($(P)-(I')$) circle (.5pt) node  [above] {$\sqrt{2t}$}
					 (I) circle (.5pt) ++(-130:.15) node {$2$}
					(B)++($(P)-(I)$) circle (.5pt) node [above] {$2$}; 
				\pic [draw,angle radius=0.1cm] {right  angle=P--I'--D'};
				 \pic [draw,angle radius=0.1cm] {right angle=P--I--D};
			\end{tikzpicture}
		} 
	} 
\end{bt}

% nguồn: Phan2_Lop_12_Dai_so_Chuong_4
\begin{bt}%[2D4V3-5]%[ma:7f19beb6]
	Họa sĩ thiết kế một micro có dạng khối tròn xoay, mặt cắt đứng chứa trục của khối tròn xoay có dạng như hình sau, trong đó $OA=OB=OI=2$ cm, $MC=MD=1$ cm, đường thẳng $OM$ là đường trung trực của đoạn thẳng $CD$,  $OM=20$ cm, $\widehat{AOB}=90^\circ $. Thể tích của micro này là bao nhiêu cm$^3$? (Làm tròn kết quả đến hàng đơn vị).
	\begin{center}
		\resizebox{0.8\linewidth}{!}{\begin{tikzpicture}
		\coordinate (A) at  (1.41,1.41);
		\coordinate (B) at (1.41,-1.41); 
		\coordinate  (O) at (0,0);
		\coordinate (C)  at (15,1);
		\coordinate (D)  at (15,-1);
		\coordinate  (M) at (15,0)
			(0,0) node[below  left]{$O$};
		\draw[fill]  (0,0) circle (1pt)
			(2,0) node [below  right] {$I$}
			 (15,0) node [right] {$M$};
		\draw (1.41,1.41)  arc (45:315:2);
		\draw[dashed](1.41,-1.41) arc  (-45:45:2)
			 (A)--(O)--(B)
			(O)--(M); 
		\draw  (A)--(C)--(D)--(B);
		\fill 
			(A) node[shift={(70:11pt)}]{$A$}circle  (1pt)
			(B) node[shift={(-70:11pt)}]{$B$}circle  (1pt)
			(C) node[shift={(90:11pt)}]{$C$}circle  (1pt)
			(D)  node[shift={(-90:11pt)}]{$D$}circle (1pt);
	\end{tikzpicture}}
	\end{center}
	\shortans{117} 
	 \loigiai{
		\begin{center}
			\resizebox{0.8\linewidth}{!}{\begin{tikzpicture}
		\coordinate (A)  at (1.41,1.41);
		\coordinate (B) at (1.41,-1.41); 
		 \coordinate (O) at (0,0);
		 \coordinate (C) at (15,1);
		 \coordinate (D) at (15,-1);
		\coordinate (M) at (15,0) 
			(0,0) node[below left]{$O$}; 
		\path (2,0) coordinate  (I)
			 ($(O)!(A)!(I)$) coordinate (H);
		\draw[fill]  (0,0) circle (1pt)
			 (2,0) node [below right] {$I$}
			node [above  right] {$2$}
			(15,0) node [below right] {$M$} 
			node [above right]  {$20$}
			 (-0.35,-2.3) node{$-2$}
			(-0.25,2.2) node{$2$}; 
		\draw  (1.41,1.41) arc (45:315:2);
		\draw[dashed](1.41,-1.41) arc (-45:45:2) 
			(A)--(O)--(B) 
			(O)--(M) 
			(A)--(B); 
		\draw  (A)--(C)--(D)--(B);
		 \draw[fill] (A) node[shift={(70:11pt)}]{$A$}circle (1pt)
			(B) node[shift={(-70:11pt)}]{$B$}circle  (1pt)
			(C) node[shift={(90:11pt)}]{$C$}circle (1pt) 
			(D) node[shift={(-90:11pt)}]{$D$}circle (1pt) 
			(H)  node[shift={(45:11pt)}]{$H$}circle (1pt);
		\draw[->] (0,-3)--(0,3) node[left]{$y$}; 
		 \draw[->] (-2.5,0)--(16,0) node[below]{$x$};
	\end{tikzpicture}}
		\end{center}
		Đặt hệ trục $Oxy$ như hình vẽ, đơn vị mỗi trục là  $1$ cm.\\
		Gọi  $K(-2;0)$ và $H$ là giao điểm của $AB$ và $OI$.\\
		Ta có tam  giác $OAB$ là tam giác vuông cân, nên điểm $A(\sqrt{2};\sqrt{2})$.\\
		Gọi $y=f(x)$ là đồ thị bao gồm cung nhỏ $KA$ và  đoạn thẳng $AC$. Thể tích micro bằng
		\begin{eqnarray*}
			V=\displaystyle  \pi\int_{-2}^{20}{f(x)\mathrm{\,d}x}=\pi\int_{-2}^{\sqrt{2}}{f(x)\mathrm{\,d}x}+\pi\int_{\sqrt{2}}^{20}{f(x)}^2\mathrm{\,d}x
		\end{eqnarray*}
		Phương trình đường tròn là $x^2+y^2=4$, hàm số có đồ thị cung nhỏ $KA$  là $y=\sqrt{4-x^2}$, $x \in [-2;\sqrt{2}]$.\\
		Vì $A\left(\sqrt{2};\sqrt{2}\right)$, $C(20;1)$  nên phương trình đường thẳng $AC$ là
		\begin{eqnarray*}	\left(AC\right)\colon\dfrac{x-20}{20-\sqrt{2}}=\dfrac{y-1}{1-\sqrt{2}}\Leftrightarrow y=\dfrac{\left(1-\sqrt{2}\right)\left(x-20\right)}{20-\sqrt{2}}+1.
		\end{eqnarray*}
		Thể tích  micro là
		\begin{eqnarray*}
			 V=\displaystyle\pi\int_{-2}^{\sqrt{2}}{\left(4-x^2\right)\mathrm{\,d}x}+\pi\int_{\sqrt{2}}^{20}{{\left(\dfrac{\left(1-\sqrt{2}\right)\left(x-20\right)}{20-\sqrt{2}}+1\right)}^2\mathrm{\,d}x}\approx 117\left(\text{cm}^3\right).
	\end{eqnarray*}}
\end{bt}

% nguồn: Phan2_Lop_12_Dai_so_Chuong_4
\begin{bt}%[2D4V3-5]%[ma:3b516d97]
	Để đặt được một vật trang trí trên mặt bàn, người ta thiết kế  một chân đế như sau. Lấy một
	khối gỗ có dạng khối chóp cụt tứ giác đều với độ  dài hai cạnh đáy lần lượt bằng $7{,}4$ cm và $10{,}4$ cm,
	bề dày của khối gỗ bằng $1{,}5$  cm. Sau đó khoét bỏ đi một phần của khối gỗ sao cho phần đó có dạng
	 vật thể $H$, ở đó $H$ nhận được bằng cách cắt khối cầu bán kính $5{,}5$ cm bởi một mặt phẳng cắt mà
	mặt cắt là hình tròn bán  kính $3{,}5$ cm (xem hình dưới).
	Thể tích của khối chân đế bằng bao nhiêu centimét khối (không làm tròn kết  quả các phép tính trung
	 gian, chỉ làm tròn kết quả cuối cùng đến hàng phần mười)?
	\begin{center}
		\begin{tikzpicture}[scale=0.6,declare function={gocx=80; goc=-150; a=5; b=a/2; h=4;}]
				\begin{scope}[overlay]
					\path (0,0) coordinate (A)--+(gocx:h) coordinate  (S)
						(a,0)  coordinate (B)
						(goc:b) coordinate  (D)
						+(a,0) coordinate  (C)
						($(A)!.5!(C)$) coordinate  (O)
						 ($(A)!.4!(S)$) coordinate (A')
						($(B)!.4!(S)$)  coordinate (B')
						($(C)!.4!(S)$) coordinate  (C')
						($(D)!.4!(S)$) coordinate  (D')
						($(O)+(0,0.45*h)$) coordinate (O'); 
				\end{scope}
				\draw[dashed]  (A')--(A)
					(D)--(A)--(B); 
				\draw  (C)--(C')
					(D')--(D)--(C)--(B)--(B')--(C')--(D')--(A')--(B'); 
				\path (A')--(B')--([turn]90:7pt) coordinate  (xt)
					 ($(A')-(B')+(xt)$) coordinate (yt);
				 \draw[>=stealth,|<->|] (xt)--(yt) node[fill=white,inner sep=0pt,font=\scriptsize,pos=0.5,sloped]{$7{,}4$ cm};
				 \path (D)--(C)--([turn]-90:7pt) coordinate (xt)
					 ($(D)-(C)+(xt)$) coordinate (yt);
				\draw[>=stealth,|<->|]  (xt)--(yt) node[fill=white,inner sep=0pt,font=\scriptsize,pos=0.5,sloped]{$10{,}4$ cm};
				\draw[>=stealth,|<->|,dashed]  (O)--(O') node[left,pos=0.5,font=\scriptsize]{$1{,}5$ cm};
			\end{tikzpicture}
		\begin{tikzpicture}[scale=0.7]
				\tikzset{declare function={ 
					R=3;  d=-0.65*R; theta=75; goct = acos (-1.1*d/R);
					r=sqrt (R*R-d*d); betacrit=asin (d/r*cot  (theta)); tgoc=330;},
					samples=200,smooth} 
				\tdplotsetmaincoords{theta}0 
				\path[ball color=blue] (0,0)  coordinate (O) circle (R);
				\path[fill=white] ({-goct}:R) arc  ({-goct}:{180+goct}:R)--cycle;
				\draw (O) circle (R); 
				\begin{scope}[tdplot_main_coords,line join=round]
					\path[inner color=blue,outer  color=blue!75!black] plot[variable=\t,domain=0:360]({r*cos (\t)},{r*sin (\t)},d);
					 \draw[dashed,dash pattern=on 2pt off 1.5pt] plot[variable=\t,domain=0:180]({R*cos (\t)},{R*sin (\t)},0);
					\draw plot[variable=\t,domain=180:360]({R*cos (\t)},{R*sin  (\t)},0);
				\end{scope}
				\draw[latex-,cyan]  (-95:{0.95*R})--+(-150:1) node[left]{$H$};
			\end{tikzpicture}
		\begin{tikzpicture}[scale=0.6]
				\tikzset{declare  function={
					R=3.5; d=2; theta=70; goct = acos (d/R); 
					r=sqrt (R*R-d*d);  betacrit=asin (d/r*cot (theta)); tgoc=330;},samples=200,smooth}
				\tdplotsetmaincoords{theta}0 
				\begin{scope}[rotate=-90]
					\draw (0,0) coordinate (O) circle  (R);
					\begin{scope}[tdplot_main_coords,line join=round]
						\draw[dashed] 
						plot[variable=\t,domain=betacrit:180-betacrit]({r*cos (\t)},{r*sin (\t)},d) 
						plot[variable=\t,domain=0:180]({R*cos (\t)},{R*sin  (\t)},0);
						\draw  plot[variable=\t,domain=betacrit:-180-betacrit]({r*cos (\t)},{r*sin (\t)},d);
						 \draw plot[variable=\t,domain=180:360]({R*cos (\t)},{R*sin (\t)},0);
					\end{scope}
				\end{scope}
				\draw[-stealth]  (-1.2*R,0)--(-R,0)
					(R,0)--(1.25*R,0)  node[below]{$x$};
				\draw[-stealth] (0,-1.2*R)--(0,-R) 
					 (0,R)--(0,1.25*R) node[right]{$y$};
				\draw[dashed,shorten <=2pt,shorten >=2pt]  (-R,0)--(R,0);
				\draw[dashed,shorten <=2pt,shorten >=2pt] (0,-R)  coordinate (By)--(0,R) coordinate (Ay) node[above left=-3pt,font=\footnotesize]{$5{,}5$};
				\draw[dashed] (-goct:R) -- (goct:R)  coordinate (M)--($(By)!(M)!(Ay)$) node[left=-3pt,font=\footnotesize]{$3{,}5$};
				\fill (0,0) circle (1pt) node[below left]{$O$}; 
				\draw[latex-] (-5:{0.95*R})--+(-30:1)  node[below]{$H$};
			\end{tikzpicture}
	\end{center}
	 \shortans{94,7}
	\loigiai{ 
		Thể tích khối gỗ ban đầu  là
		\[ 
		V_{\text{  gỗ}} = \dfrac{1{,}5}{3} \left( 10{,}4^2 + 7{,}4^2 + \sqrt{10{,}4^2 \cdot 7{,}4^2} \right) = 119{,}94 \quad (\text{cm}^3).
		\] 
		Để tìm thể tích khối chân đế, ta cần tìm thể tích vật thể $H$. Gọi $O$ là tâm của khối cầu và $R = 5,5 \, \text{ cm}$ là bán  kính khối cầu. Khối cầu được tạo thành bằng cách quay hình tròn $(C)$ có tâm $O(0; 0)$, bán kính $R$ quanh trục hoành $Ox$. Ta có bán kính mặt cắt
		 \[
		OI =  \sqrt{R^2 - IA^2} = \sqrt{5{,}5^2 - 3{,}5^2} = 3\sqrt{2} \quad (\text{cm}).
		 \]
		Đường  tròn $(C)$ có phương trình
		\[ 
		x^2 + y^2 = (5{,}5)^2 \Rightarrow y^2 =  30{,}25 - x^2.
		 \]
		Thể tích vật thể $H$  là
		 \[
		V_H = \pi \int_{3\sqrt{2}}^{5{,}5} \left( 30{,}25 - x^2 \right) \mathrm{\,d}x = \pi \left( 30{,}25x - \dfrac{x^3}{3} \right) \Biggr|_{3\sqrt{2}}^{5{,}5} = \pi  \left( \dfrac{1331}{12} - \dfrac{291\sqrt{2}}{4} \right).
		 \]
		Vậy thể tích khối gỗ  cần tìm là
		\[ 
		V  = V_{\text{ gỗ}} - V_H \approx 94{,}7 \quad (\text{cm}^3).
		 \]
	 }
\end{bt}

% nguồn: Phan2_Lop_12_Dai_so_Chuong_4
\begin{bt}%[2D4V3-5]%[ma:7c435954]
	Hình elip được ứng dụng nhiều trong thực tiễn, đặc biệt là kiến trúc, xây dựng, thiết bị nội thất,... Mặt trong (lọt lòng) và ngoài (phủ bì) của một bồn rửa (lavabo) bằng sứ có hình dạng là một nửa khối tròn xoay khi quay quanh một trục của $2$ elip có chung các trục đối xứng (hình minh họa). Thông số kĩ thuật mặt trên của bồn rửa dài $\times $ rộng là  $660\times 380$ mm (phủ bì) và elip (lọt lòng) có trục lớn, trục nhỏ ít hơn elip phủ bì một khoảng $40$ mm. Tính thể tích chứa nước của bồn rửa (đơn vị: lít) (làm tròn kết quả đến hàng phần mười).
	\begin{center}
		\begin{tikzpicture}[scale=1]
				\draw (0,0) ellipse (3.3 and  1.9);
				\draw (0,0)  ellipse (3 and 1.6);
				\draw (-3.70,0) -- (3.70,0); 
				\node at (0,0.2)  {trục xoay};
				\draw[dashed] (0,-1.9) --  (4.10,-1.9);
				\draw[dashed] (0,1.9) --  (4.10,1.9);
				\draw[<->]  (3.6,1.9) -- (3.6,-1.9);
				\node[rotate=90] at  (4.10,0) {380 mm};
				\draw[dashed] (3.30,0) --  (3.3,-2.5);
				\draw[dashed]  (-3.30,0) -- (-3.3,-2.5);
				\draw[<->] (-3.3,-2.5)  -- (3.3,-2.5);
				\node at  (0,-2.8) {660 mm};
				\draw (0,1.6) -- (0,1.9); 
				\draw[dashed,->] (-2.75,1.85) -- (0,1.75) 
					(-2.75,1.85) -- (-3.15,0); 
				 \draw[dashed,->] (-2.75,1.85) -- (0,1.75);
				\node  at (-2.75,2.3) {$\Delta = 20$ mm};
			\end{tikzpicture}
	\end{center}
	\shortans{18,8} 
	 \loigiai{
		Chọn hệ trục tọa độ $Oxy$ thích hợp với đơn vị trên trục là decimet. 
		\begin{center}
			\begin{tikzpicture}[scale=1]
					 \draw (0,0) ellipse (3 and 1.6);
					\draw[->] (-3.70,0) -- (4,0) node[above]{$x$}; 
					\draw[->] (0,-2) -- (0,2) node[above]{$y$}; 
					\draw (0,0)  node[below left]{$O$}
						(-3.1,0) node[below  left]{$-3{,}1$}
						 (3.1,0) node[below right]{$3{,}1$}
						(0,-1.7) node[below  left]{$-1{,}7$}
						(0,1.7) node[above left]{$1{,}7$}; 
				\end{tikzpicture}
		\end{center}
		Phương  trình elip lọt lòng $(E)\colon\dfrac{x^2}{3{,}1^2}+\dfrac{y^2}{1{,}7^2}=1\Leftrightarrow y=\pm 1{,}7\sqrt{1-\dfrac{x^2}{3{,}1^2}}$.\\
		Thể tích chứa nước của bồn rửa $V=\dfrac{1}{2}\cdot\pi\displaystyle\int_{-3{,}1}^{3{,}1}{1{,}7^2\left(1-\dfrac{x^2}{3{,}1^2}\right)dx}\approx 18{,}8$ lít. 
	} 
\end{bt}

% nguồn: Phan2_Lop_12_Dai_so_Chuong_4
\begin{bt}%[2D4V3-5]%[ma:b0c94353]
	Một nhà máy nhiệt điện sử dụng $90$ máng parabol thu nhiệt năng lượng mặt trời có cùng kích thước, bề mặt cong  đều nhau. Mỗi máng có chiều rộng $2$\,(m), bề dày của khối silic làm mặt máng là $2$\,(dm), chiều dài $3$\,(m). Đặt máng tiếp giáp mặt đất có điểm cao nhất của khối silic làm mặt máng so với mặt đất là $5$\,(dm). Khi đó thể tích của khối silic làm $90$ mặt máng bằng bao nhiêu m$^3$?
	 \shortans{108}
	 \loigiai{
	\immini{Xét hệ trục $Oxy$ như  hình vẽ.\\ Khi đó đường parabol tương ứng với vành trên và vành dưới của máng lần lượt là $(P_1)\colon f(x)=ax^2+b$ và $(P_2)\colon g(x)=mx^2+n$.\\
		Vì $(P_1)$ đi qua các  điểm $(0;0{,}5)$ và $(-1{,}2;0)$ nên
		$f(x)=-\dfrac{25}{72}x^2+0{,}5$.\\ 
		Lại có $(P_2)$ đi  qua các điểm $(0;0{,}3)$ và $(-1;0)$ nên
		$g(x)=-\dfrac{3}{10}x^2+0{,}3$.\\ 
		 }{
		\begin{tikzpicture}[scale=2,line cap=butt,line join=miter,>=stealth]
				 \draw[->] (-2,0)--(2,0)
					 node[shift={(-100:7pt)},font=\normalsize]{$x$};
				 \draw[->] (0,-0.75)--(0,1.5)
					 node[shift={(180:7pt)},font=\normalsize]{$y$};
				\draw (0,0)  node[shift={(225:9pt)},font=\normalsize]{$O$};
				\foreach \x in  {-1,1}{
					\fill (\x,0) circle (.5pt)  node[font=\scriptsize,shift={(-90:5pt)}]{\x};
				} 
				\fill (-1.2,0) circle (.5pt) node[font=\scriptsize,shift={(150:8pt)}]{$-1{,}2$}; 
				\fill  (1.2,0) circle (.5pt) node[font=\scriptsize,shift={(40:7pt)}]{$1{,}2$};
				 \fill (0,0.3) circle (.5pt) node[font=\scriptsize,shift={(220:10pt)}]{$0{,}3$};
				\fill (0,0.5) circle (.5pt) node[font=\scriptsize,shift={(150:10pt)}]{$0{,}5$}; 
				\tikzset{declare function={f  (\x)=-25*(\x)^2/72+1/2; g (\x)=-3*(\x)^2/10+3/10;}}
				\begin{scope}
					\draw[samples=100]  plot[domain=-1.5:1.5] (\x,{f (\x)});
					\draw[samples=100]  plot[domain=-1.5:1.5] (\x,{g (\x)});
				\end{scope}
			\end{tikzpicture}
		 }
		\noindent Diện tích mặt cắt của  máng $S=\displaystyle\int_{-1{,}2}^{1{,}2}f(x)\,\mathrm{d}x-\displaystyle\int_{-1{,}2}^{1{,}2}g(x)\,\mathrm{d}x=0{,}4$.\\
		Thể tích khối silic là $V=90\cdot  3\cdot 0{,}4=108$ (m$^3$).
	} 
\end{bt}

% nguồn: Phan2_Lop_12_Dai_so_Chuong_4
\begin{bt}%[2D4V3-5]%[ma:27526091]
	Hình $1$ là một tác phẩm dự  thi của nhà thiết kế sân khấu trong một cuộc thi thiết kế sân khấu ngoài trời tổ chức tại một quảng trường. Khi mở rộng sân khấu trung tâm, ta được hình $2$. Quá trình thiết kế sân khấu trung tâm được mô tả như sau:
	\begin{center}
		\begin{tikzpicture}[scale=1.0,font=\footnotesize,line join=round,line cap=round,>=stealth]
				\tikzstyle{ring segment}=[draw,fill=gray!25,line join=round] 
				 \tikzstyle{arrow}=[draw,fill=gray!60,line join=round,line cap=round]
				\begin{scope}
					\node[draw,minimum size=2.4cm,inner sep=0pt] (sq) at (0,0) {}; 
					\draw (sq.north  west) -- (sq.south east);
					 \draw (sq.south west) -- (sq.north east);
					\node[regular polygon,regular polygon sides=8,minimum size=1.7cm,rotate=22.5] at (0,0) {}; 
				\end{scope}
				\begin{scope}
					 \foreach \i in {0,...,9}{
						\pgfmathsetmacro{\startangle}{120 +  \i * 30}
						\pgfmathsetmacro{\endangle}{120 + (\i+1) * 30}
						\path[ring segment] (\startangle:1.8) arc (\startangle:\endangle:1.8) -- (\endangle:2.2) arc (\endangle:\startangle:2.2) -- cycle;
					} 
					\foreach \i  in {0,...,9}{
						 \pgfmathsetmacro{\startangle}{120 + \i * 30}
						\pgfmathsetmacro{\endangle}{120 + (\i+1) * 30}
						\draw (\startangle:2.5) arc (\startangle:\endangle:2.5) -- (\endangle:2.9) arc (\endangle:\startangle:2.9) -- cycle;
					 }
					\path[ring segment,fill=gray!40] (120:1.8) -- (120:2.9) -- (123:2.9) arc (123:120:2.9) -- (123:1.8) arc (123:120:1.8) -- cycle; 
					\path[ring segment,fill=gray!40] (60:1.8) -- (60:2.9) --  (57:2.9) arc (57:60:2.9) -- (57:1.8) arc (57:60:1.8) -- cycle;
				\end{scope}
				\node at  (1,-4)[above left]{Hình $1$};
			\end{tikzpicture}
		 \qquad
		\begin{tikzpicture}[scale=1.2,font=\footnotesize,line join=round,line cap=round,>=stealth]
				 \path
					(0,0)  coordinate (B)
					(0,4) coordinate (A) 
					 (4,0) coordinate (C)
					($(A)+(C)-(B)$) coordinate (D) 
					(0,2) coordinate  (E)
					(2,0)  coordinate (F)
					(4,2) coordinate  (G)
					(2,4) coordinate  (H);
				\draw (A) rectangle (C); 
				 \draw[smooth,samples=300,domain=0:4] plot (\x,{-1*(\x)*(\x-4)})
					plot (\x,{(\x-2)*(\x-2)}); 
				 \draw[smooth,samples=300,domain=0:4] plot (\x,{2+sqrt (4-\x)})
					 plot (\x,{2-sqrt (4-\x)});
				\draw[smooth,samples=300,domain=0:4] plot  (\x,{2+sqrt (\x)})
					plot (\x,{2-sqrt  (\x)});
				 \foreach \p/\q in {A/90,B/180,C/0,D/45,E/180,F/-90,G/0,H/90}
				\fill (\p) circle (1.0pt) 
					($(\p)+(\q:3mm)$)  node{$\p$};
				 \fill[draw,fill=gray!60] (0.586,2)--(1,1)--(2,0.586)--(3,1)--(3.414,2)--(3,3)--(2,3.414)--(1,3)--cycle;
				\node at (2.8,-1)[above left]{Hình  $2$};
			\end{tikzpicture}
	\end{center}
	\textbf{Bước $1$.} Vẽ hình vuông $ABCD$ có độ dài cạnh bằng $2$, và lấy trung điểm của bốn cạnh  lần lượt là $E$, $F$, $G$, $H$.\\
	\textbf{Bước $2$.} Vẽ đồ thị của các hàm bậc hai đi qua ba điểm $B$, $C$, $H$ và hàm  bậc hai đi qua ba điểm $F$, $D$, $A$.\\
	\textbf{Bước $3$.} Tương tự như bước $2$, vẽ đồ thị của các hàm  bậc hai đi qua ba điểm $A$, $B$, $G$ và ba điểm $C$, $D$, $E$.\\
	Biết rằng: Diện tích  phần tô đen trong hình $2$ được cho bởi công thức: $\dfrac{p\sqrt{2}+q}{3}$.\\
	Hãy tính giá trị của $p-q$, (với $p$, $q$ là các  số nguyên).
	\shortans{15} 
	 \loigiai{
		Chọn hệ  trục tọa độ $Oxy$ như hình vẽ:
		\begin{center}
			\begin{tikzpicture}[scale=1.2,font=\footnotesize,line join=round,line cap=round,>=stealth]
					 \path
						(0,0) coordinate (B) 
						(0,4)  coordinate (A)
						(4,0) coordinate (C) 
						($(A)+(C)-(B)$) coordinate  (D)
						(0,2)  coordinate (E)
						 (2,0) coordinate (F)
						(4,2)  coordinate (G)
						(2,4)  coordinate (H)
						 (intersection of A--C and B--D) coordinate (O)
						(3,1)  coordinate (N)
						(3,2) coordinate (M); 
					\draw  (A) rectangle (C);
					 \draw[smooth,samples=300,domain=0:4] plot (\x,{-1*(\x)*(\x-4)})
						plot (\x,{(\x-2)*(\x-2)}); 
					\draw[smooth,samples=300,domain=0:4] plot (\x,{2+sqrt  (4-\x)})
						plot  (\x,{2-sqrt (4-\x)});
					\draw[smooth,samples=300,domain=0:4] plot  (\x,{2+sqrt (\x)})
						 plot (\x,{2-sqrt (\x)});
					\foreach \p/\q in {A/90,B/180,C/0,D/45,E/150,F/-70,G/30,H/70,N/-60} 
					\fill (\p)  circle (1.0pt)
						($(\p)+(\q:3mm)$)  node{$\p$};
					\fill[draw,fill=gray!60] (0.586,2)--(1,1)--(2,0.586)--(3,1)--(3.414,2)--(3,3)--(2,3.414)--(1,3)--cycle; 
					\draw (N)--(M) node[above]{$\dfrac{1}{2}$} 
						 (A)--(C)
						 (B)--(D);
					\draw[->] ($(O)!1.5!(E)$) --($(O)!1.5!(G)$) node[below]{$x$}; 
					\draw[->] ($(O)!1.5!(F)$)  --($(O)!1.5!(H)$) node[left]{$y$};
					\node at (O)[above left]{$O$}; 
				\end{tikzpicture}
		\end{center}
		Gốc tọa độ $O$ là  giao điểm của hai đường chéo $AC$ và $BD$.\\
		Trục $Ox\equiv  EG$; $Oy\equiv HF$.\\
		Khi đó ta dễ suy ra được $AC\colon y=-x$ và đồ thị hàm số bậc hai đi qua $F$, $D$, $A$ là  $y=2x^2-1$.\\
		Xét phương trình hoành độ giao điểm của đường thẳng $AC$ và đồ thị hàm số bậc hai đi qua $F$, $D$, $A$ ta được $2x^2+x-1=0\Leftrightarrow\hoac{&x=-1\\& x=\dfrac{1}{2}.}$  \\
		Ta gọi $S_1$ là diện tích vùng được tô màu như  hình trên, khi đó:\\
		$$S_1=\displaystyle\int_{\frac{1}{2}}^{\frac{1}{\sqrt{2}}}\left(1-2 x^2\right)  \mathrm{d}x=\left(x-\dfrac{2}{3} x^3\right)\bigg|_{\frac{1}{2}}^{\frac{1}{\sqrt{2}}} =\dfrac{\sqrt{2}}{3}-\dfrac{5}{12}.$$
		Diện tích phần tô màu trong hình trên là $S_2=\dfrac{1}{8}$.\\ 
		Vậy diện tích tô màu như hình  $2$ là $S=8(S_1+S_2)=8\left(\dfrac{\sqrt{2}}{3}-\dfrac{5}{12}+\dfrac{1}{8}\right)=\dfrac{8\sqrt{2}-7}{3}\Rightarrow\heva{&p=8\\& q=7.}$ \\
	 Vậy $p-q=15$.}
\end{bt}

% nguồn: Phan2_Lop_12_Dai_so_Chuong_4
\begin{bt}%[2D4V3-5]%[ma:ade03dfe]
	Một bức tường hình chữ nhật $ABCD$ có kích thước $6$ (m),  $4$(m) được bạn Hà trang trí bằng cách vẽ hai đồ thị $f(x)=a^x$, $g(x)=\log_bx$ đối xứng nhau qua đường thẳng $d\colon y=x$ và chia thành ba phần (tham khảo hình vẽ bên). Phần $H_1$ được sơn màu xanh da trời, phần $H_2$ được sơn màu trắng, phần $H_3$ được sơn màu xanh lá cây. Biết rằng mỗi hộp sơn các màu chỉ sơn được $3\left(\text{m}^2\right)$ tường, đồng thời giá của hộp sơn màu xanh da trời là $100.000$ đồng/hộp, hộp sơn trắng là $140.000$ đồng/hộp, hộp sơn xanh lá cây là $130.000$ đồng/hộp. Tính giá tiền bạn Hà mua để sơn bức tường này? (đơn~vị là triệu đồng và cửa hàng sơn chỉ bán số nguyên của hộp).
	\begin{center}
		\begin{tikzpicture}[line cap=butt,line join=miter,>=stealth]
				\tikzset{ 
					 declare function={
					f (\x) =  ln (\x)/ln (sqrt (3));
					g  (\x) = 1*(sqrt (3))^(\x);
					h (\x) = \x; 
					 }}
				\begin{scope}
					\clip  (-3,-2) rectangle (3,2);
					\path[fill=blue!65!cyan,fill  opacity=0.3,draw]
						plot[domain=-3.1:1.5] (\x,{g  (\x)})--(-3.1,5)--cycle;
					 \path[fill=green!65!cyan,fill opacity=0.3,draw]
						 plot[domain=0.05:5.1] (\x,{f (\x)})--(5,-2)--cycle;
				\end{scope}
				\path  (-3,-2) node[below]{$A$}
					(3,-2) node[below right]{$B$} 
					 (3,2) node[below right]{$C$}
					 (-3,2) node[above]{$D$}
					 (-2,1) node[above]{$H_1$}
					 (-2,-1) node[above]{$H_2$}
					(2,-1) node[above]{$H_3$}; 
				\draw  (-3,-2)--(-3,2)--(3,2)--(3,-2)--cycle;
			\end{tikzpicture}
	\end{center}
	\shortans{1,15} 
	\loigiai{\begin{center}
		\begin{tikzpicture}[line cap=butt,line join=miter,>=stealth]
				\tikzset{ 
					declare  function={
					f (\x) = ln  (\x)/ln (sqrt (3));
					 g (\x) = 1*(sqrt (3))^(\x);
					h  (\x) = \x;
					 }}
				\begin{scope}
					\clip (-3,-2) rectangle (3,2); 
					\path[fill=blue!65!cyan,fill  opacity=0.3,draw]
						 plot[domain=-3.1:1.5] (\x,{g (\x)})--(-3.1,5)--cycle;
					\path[fill=green!65!cyan,fill opacity=0.3,draw] 
						plot[domain=0.05:5.1] (\x,{f (\x)})--(5,-2)--cycle; \draw[smooth,samples=100] plot[domain=-3:4] (\x,{h (\x)}); 
				\end{scope}
				 \draw[->] (-3,0)--(4,0)
					 node[shift={(-100:7pt)},font=\normalsize]{$x$};
				\draw[->]  (0,-2)--(0,2.5)
					 node[shift={(180:7pt)},font=\normalsize]{$y$};
				\draw  (5pt,0) |- (0,5pt)
					(0,0)  node[shift={(330:9pt)},font=\normalsize]{$O$};
				\foreach \x in  {1}{
					\draw (\x,1.5pt)--(\x,-1.5pt)  node[font=\scriptsize,shift={(-50:5pt)}]{\x};
				 }
				\foreach \y in {-1,1}{ 
					\draw (1.5pt,\y)--(-1.5pt,\y) node[font=\scriptsize,shift={(180:5pt)}]{\y}; 
				} 
				\path (-3,-2)  node[below right]{$A$}
					(3,-2)  node[below right]{$B$}
					(3,2)  node[below right]{$C$}
					(-3,2) node[above]{$D$} 
					(-2,1) node[above]{$H_1$} 
					(-2,-1) node[above]{$H_2$} 
					(2,-1)  node[above]{$H_3$};
				\draw  (-3,-2)--(-3,2)--(3,2)--(3,-2)--cycle;
			\end{tikzpicture}
	\end{center}
	Do đồ thị của $f(x)=a^x$ và  $g(x)=\log_bx$ đối xứng nhau qua đường thẳng $y=x$ nên $a=b$.\\
	Mặt khác đồ thị của  $y=\log_bx$ đi qua điểm $C(3;2)$ nên $b=\sqrt{3}$. \\
	Hai hàm số là  $f(x)=(\sqrt{3})^x$ và $g(x)=\log_{\sqrt{3}}x$.\\
	Đồ thị  của $f(x)$ cắt $y=2$ tại điểm $\left(\log_{\sqrt{3}}2;2\right)$ và đồ thị của $g(x)$ cắt $y=-2$ tại $\left(\dfrac{1}{3};-2\right)$. \\
	Diện tích của  miền $H_1$ là $S_1=\displaystyle\int^{\log_{\sqrt{3}}2}_{-3}\left(2-(\sqrt{3})^x\right)\,\mathrm{d}x\approx 5{,}233$. \\
	Diện tích của miền $H_3$ là $S_3=\displaystyle\int^3_{\tfrac{1}{3}}\left(\log_{\sqrt{3}}x+2\right)\,\mathrm{d}x\approx  7{,}145$. \\
	Diện tích của miền $H_2$ là $S_2=6\times 4 - 5{,}233  - 7{,}145=11{,}622$.\\
	Vậy cần dùng $2$ hộp sơn màu xanh da trời, $3$ hộp sơn màu xanh  lá cây và $4$ hộp sơn màu trắng.\\
	Số tiền bạn  Hà cần dùng mua sơn là\\ $T=100.000\times 2+130.000\times 3+140.000\times 4=1.150.000$ (đồng).
 }
\end{bt}

% nguồn: Phan2_Lop_12_Dai_so_Chuong_4
\begin{bt}%[2D4V3-5]%[ma:a8f71530]
	Một bồn hình trụ đang chứa dầu, được đặt nằm ngang, có chiều dài bồn là $5$ m, có bán kính đáy $1$  m. Người ta đã rút dầu trong bồn tương ứng với $0{,}5$ m của đường kính đáy. Tính thể tích của lượng dầu còn lại trong bồn (theo đơn vị m$^3$).
	\begin{center}
		\begin{tikzpicture}[font=\footnotesize,line join=round,line cap=round,>=stealth,declare function={a=1; b=2.5*a; d=3;}]
				 \path (d,-b) arc (-90:30:a and b) coordinate (B) arc (30:160:a and b) coordinate (A)
					(-d,-b) arc (-90:30:a and b) coordinate  (C) arc (30:160:a and b) coordinate (D)
					 (d,0) coordinate (O)
					(-d,0)  coordinate (O');
				\fill[gray!60](B)--(C)--(D) arc (160:270:a and b)--(d,-b) arc (-90:30:a and b); 
				\fill[gray!30] (B)--(A)--(D)--(C)--cycle; 
				\draw (-d,-b)--(d,-b) 
					(-d,b)--(d,b); 
				\draw (d,-b)  arc (-90:270:a and b)
					(-d,-b) arc (270:90:a and b); 
				 \draw [dashed](-d,-b) arc (-90:90:a and b);
				\path  (O)--++(0,3) coordinate (O1);
				\path (intersection of O--O1 and  A--B) coordinate (x);
				 \draw[<->](x)--(d,b) node[pos=0.5,rotate=90,above]{$0{,}5$ m};

			\end{tikzpicture}
	\end{center}
	\shortans{12,6} 
	 \loigiai{
		\begin{center}
			\begin{tikzpicture}[font=\footnotesize,line join=round,line cap=round,>=stealth,declare function={a=1; b=2.5*a; d=3;}]
					\path (d,-b) arc (-90:30:a and b) coordinate (B) arc (30:160:a and b)  coordinate (A)
						(-d,-b) arc (-90:30:a and b) coordinate (C) arc (30:160:a and b)  coordinate (D)
						(d,0) coordinate (O) 
						 (-d,0) coordinate (O');
					\fill[gray!60](B)--(C)--(D) arc  (160:270:a and b)--(d,-b) arc (-90:30:a and b);
					\fill[gray!30]  (B)--(A)--(D)--(C)--cycle;
					\draw  (-d,-b)--(d,-b)
						 (-d,b)--(d,b);
					\draw (d,-b) arc (-90:270:a  and b)
						(-d,-b) arc  (270:90:a and b);
					\draw [dashed](-d,-b) arc (-90:90:a and b); 
					\draw[->] (O)--++(0,3) coordinate (O1)  node[anchor=west]{$x$};
					\draw[->] (O)--($(O)!1.5!(O')$) node[anchor=south]{$y$}; 
					\path (intersection of  O--O1 and A--B) coordinate (x);
					\draw[<->](x)--(d,b)  node[pos=0.5,rotate=90,above]{$0{,}5$ m};
					\foreach \p/\g in  {B/45,A/135,C/45,D/135,O/0,O'/-90,x/45} \fill (\p) circle (1pt)+(\g:0.3) node{$\p$};
				\end{tikzpicture}
		\end{center}
		Thiết  diện cắt trục $Ox$ tại điểm có hoành độ là $x$.\\
		Khi  đó, ta có $AB=2\sqrt{1-x^2}$.\\
		Suy ra diện tích hìn chữ nhật  $ABCD$ bằng $5\cdot2\sqrt{1-x^2}=10\sqrt{1-x^2}$.\\
		Vậy thể tích lượng dầu còn lại trong bồn  là
		 $$V=10\int_{-1}^{0{,}5}\sqrt{1-x^2}\mathrm{~d}x=12{,}6\ \mathrm{~m}^3.$$
	 }
\end{bt}

% nguồn: Phan2_Lop_12_Dai_so_Chuong_4
\begin{bt}%[2D4V3-5]%[ma:10b3966c]
	 \immini{
	Một bình hoa có dạng khối tròn xoay với chiều cao là $25$ cm. Khi cắt bình hoa theo một mặt phẳng vuông góc với trục của nó thì ta luôn được thiết diện là một hình tròn có bán kính $R=\dfrac{4}{9} x^3-\dfrac{5}{3} x^2+\dfrac{4}{3} x+\dfrac{25}{36}$ (dm) với $x \in\left[0; \dfrac{5}{2}\right]$ là khoảng cách từ mặt cắt tới mặt đáy của bình hoa. Lượng nước cần đổ vào bình để mức nước trong bình cao bằng $\dfrac{2}{3}$ chiều cao của bình chiếm  tỉ lệ bao nhiêu phần trăm so với thể tích của bình hoa? (Kết quả làm tròn đến hàng đơn vị).}
	{\begin{tikzpicture}[scale=1.2,>=stealth,font=\footnotesize,line join=round,line cap=round]
		\draw[samples=100,domain=-3.5:0,rotate=270,red] plot (\x,{0.2*(\x)^3+(\x)^2+(\x)+1.2}); 
		 \draw[samples=100,domain=-3.5:0,rotate=270,red] plot (\x,{-0.2*(\x)^3-(\x)^2-(\x)+3.2});
		\draw[green!50!black,blue,pattern = north west  lines,fill=magenta] (2.2,3.5) ellipse (25pt and 5pt);
		 \draw[green!50!black,blue,pattern = north west lines,fill=orange] (2.2,0) ellipse (28pt and 5pt);
		\draw[rotate=270,pattern = north east lines,opacity=.2,line width = 1.2pt,fill=orange] 
			plot[domain=-2.5:0]  (\x,{0.2*(\x)^3+(\x)^2+(\x)+1.2})--(-1.5,3.15)--cycle;
		\draw[rotate=270,pattern = north east lines,opacity=.2,line width  = 1.2pt,fill=orange]
			plot[domain=-1.485:0] (\x,{-0.2*(\x)^3-(\x)^2-(\x)+3.2})--(0,1.2)--cycle; 
		 \draw[rotate=270,pattern = north east lines,opacity=.2,line width = 1.2pt,fill=orange]
			plot[domain=-3.5:-1.4] (\x,{-0.2*(\x)^3-(\x)^2-(\x)+3.2})--(-2.5,1.82)--cycle; 
		\draw[rotate=270,pattern = north east lines,opacity=.2,line width = 1.2pt,fill=orange] 
			plot[domain=-3.5:-2.5] (\x,{0.2*(\x)^3+(\x)^2+(\x)+1.2})--(-3.5,3)--cycle; 
	\end{tikzpicture}
	} 
	 \shortans{92}
	 \loigiai{
		Diện tích của mặt cắt vuông góc với trục tại vị trí cách mặt đáy của  bình hoa một khoảng $x$ (dm) là
		$$S(x)=\pi R^2=\pi\left(\dfrac{4}{9}x^3-\dfrac{5}{3} x^2+\dfrac{4}{3}x+\dfrac{25}{36}\right)^2 (\mathrm{dm}^2).$$ 
		Lượng nước đổ đầy bình hoa là \[V=\displaystyle\int_0^{\tfrac{5}{2}} S(x) \mathrm{\,d}x=\pi \displaystyle\int_0^{\tfrac{5}{2}}\left(\dfrac{4}{9} x^3-\dfrac{5}{3} x^2+\dfrac{4}{3}  x+\dfrac{25}{36}\right)^2\mathrm{\,d} x \approx 3{,}66~ (\mathrm{dm}^3).\]
		Lượng nước  đổ vào bình cao bằng $\dfrac{2}{3}$ chiều cao của bình là
		\[V_1=\displaystyle\int_0^{\tfrac{5}{3}} S(x)  \mathrm{\,d} x=\pi \displaystyle\int_0^{\tfrac{5}{3}}\left(\dfrac{4}{9} x^3-\dfrac{5}{3} x^2+\dfrac{4}{3} x+\dfrac{25}{36}\right)^2\mathrm{\,d}x \approx 3{,}37~ (\mathrm{dm}^3).\]
		Vậy tỉ lệ phần trăm lượng nước  cần đổ vào bình để mức nước trong bình cao bằng $\dfrac{2}{3}$ chiều cao của bình so với thể tích của bình hoa là $\dfrac{V_1}{V} \cdot 100\%=\dfrac{3{,}37}{3{,}66} \cdot 100\% \approx 92\%$.
	 }
\end{bt}

% nguồn: Phan2_Lop_12_Dai_so_Chuong_4
\begin{bt}%[2D4V3-5]%[ma:f7a7cb6d]
	Một chiếc lều vải du lịch dạng hình cong như hình vẽ. Khung chính bao gồm đáy là hình vuông cạnh $2$ m và hai xương dây $a$, $b$ nằm trên các đường parabol  đỉnh $S$. Biết chiều cao của lều là $SO=1{,}35$ m, $O$ là tâm của đáy. Tính thể tích chiếc lều theo đơn vị m$^3$. Viết kết quả dưới dạng số thập phân.
	\begin{center}
		\begin{tikzpicture}[scale=.8,>=stealth,line join=round,line cap=round,font=\footnotesize]
				\path (1.17,-.94)  coordinate (A)
					(3.39,.32) coordinate (B) 
					 (-1.17,.94) coordinate (C)
					(-3.39,-.32) coordinate (D) 
					 (0,3.45) coordinate (P);
				\draw  [fill=gray!60,draw,gray] (B).. controls +(-.75,1.58) and +(1.43,.1).. (P)
					.. controls +(-.48,.47) and  +(.16,1.09).. (C)--cycle;
				\draw [fill=gray!60,draw,gray] (P).. controls +(-.48,.47) and +(.16,1.09)..  (C)--(D)
					.. controls +(.59,1.62)  and +(-1.24,-.09).. cycle;
				\draw  [fill=gray!10,draw,gray] (D).. controls +(.59,1.62) and +(-1.24,-.09).. (P)
					.. controls +(.58,-.59) and  +(-.32,2.53).. (A)--cycle;
				\draw [fill=gray!30,draw,gray]  (P).. controls +(.58,-.59) and +(-.32,2.53).. (A)--(B)
					.. controls +(-.75,1.58) and  +(1.43,.1).. cycle;
			\end{tikzpicture}
		 \quad
		\begin{tikzpicture}[scale=1.5,>=stealth,line join=round,line cap=round,font=\footnotesize]
				 \def\g{-45}
				 \def\r{.95}
				\def\xmax{3} 
				\def\ymax{2} 
				 \def\zmax{3.5}
				 \pgfmathsetmacro\ra{sqrt (1-(\r)^2)}
				\pgfmathsetmacro\a{cos  (\g)}
				 \pgfmathsetmacro\b{sin (\g)}
				\pgfmathsetmacro\ax{\ra*cos  (\g)/sqrt ((\ra*\a)^2+(\b)^2)}
				\pgfmathsetmacro\bx{\ra*sin (\g)/sqrt  ((\ra*\a)^2+(\b)^2)}
				\def\f{-.5*(\x)^2+2*\r} 
				\path 
					(0,0) coordinate (O) 
					+({\ax},{\bx})  coordinate (i)
					 +({\bx/\ra},{-\ax*\ra}) coordinate (k)
					+(0,{\r}) coordinate  (j)
					+($2*(j)$) coordinate (P); 
				\path 
					(0,0) coordinate  (O)
					[samples at={2}]  plot ({\ax*\x},{\bx*\x+\r*\f}) coordinate (A)
					[samples  at={-2}] plot ({\ax*\x},{\bx*\x+\r*\f}) coordinate (C)
					[samples  at={2}] plot ({\bx/\ra*\x},{-\ax*\ra*\x+\r*\f}) coordinate (D)
					[samples at={-2}] plot  ({\bx/\ra*\x},{-\ax*\ra*\x+\r*\f}) coordinate (B);
				\draw  [smooth,samples=100] plot [domain=-.65:2] ({\ax*\x},{\bx*\x+\r*\f})
					plot  [domain=2:-2] ({\bx/\ra*\x},{-\ax*\ra*\x+\r*\f})
					 (D)--(A)--(B);
				\draw [dash pattern=on 1pt  off 2pt,samples=100] plot [domain=-.65:-2] ({\ax*\x},{\bx*\x+\r*\f});
				 \draw [dash pattern=on 1pt off 2pt] (D)--(C)--(B);
				\foreach \x/\y  in {A/-90,B/0,C/180,D/180}{
					 \draw (\x) circle (.5pt) ++(\y:0.2) node {$\x$};
				} 
				\draw  (P) circle (1pt) node [above] {$S$};
			\end{tikzpicture}
	\end{center}
	\shortans{2,7} 
	\loigiai{ 
		\begin{center}
			\begin{tikzpicture}[scale=1.5,>=stealth,line join=round,line cap=round,font=\footnotesize]
					 \def\g{-45}
					\def\r{.95} 
					 \def\xmax{3}
					\def\ymax{3} 
					 \def\zmax{2.5}
					\pgfmathsetmacro\ra{sqrt (1-(\r)^2)} 
					\pgfmathsetmacro\a{cos (\g)} 
					\pgfmathsetmacro\b{sin  (\g)}
					 \pgfmathsetmacro\ax{\ra*cos (\g)/sqrt ((\ra*\a)^2+(\b)^2)}
					\pgfmathsetmacro\bx{\ra*sin (\g)/sqrt  ((\ra*\a)^2+(\b)^2)}
					\def\f{-.5*(\x)^2+2*\r} 
					 \path
						(0,0) coordinate (O) 
						+({\ax},{\bx}) coordinate (i) 
						+({\bx/\ra},{-\ax*\ra}) coordinate  (k)
						+(0,{\r}) coordinate  (j)
						+($2*(j)$) coordinate  (P)
						[samples  at={2}] plot ({\ax*\x},{\bx*\x+\r*\f}) coordinate (A)
						[samples at={-2}] plot ({\ax*\x},{\bx*\x+\r*\f})  coordinate (C)
						 [samples at={2}] plot ({\bx/\ra*\x},{-\ax*\ra*\x+\r*\f}) coordinate (D)
						[samples  at={-2}] plot ({\bx/\ra*\x},{-\ax*\ra*\x+\r*\f}) coordinate (B)
						($(A)!.5!(C)$) coordinate  (I);
					 \foreach \i in {1.3}{
						 \fill [samples at={\i},cyan!90!white!60] plot ({\bx/\ra*\x},{-\ax*\ra*\x+\r*\f})--plot ({\ax*\x},{\bx*\x+\r*\f})-- plot ({-\bx/\ra*\x},{\ax*\ra*\x+\r*\f})
						--plot  ({-\ax*\x},{-\bx*\x+\r*\f})--cycle;
					 }
					\foreach \i in {1.3}{ 
						\draw [samples at={\i},gray] plot ({\bx/\ra*\x},{-\ax*\ra*\x+\r*\f}) coordinate (D')--plot ({\ax*\x},{\bx*\x+\r*\f})--plot ({-\bx/\ra*\x},{\ax*\ra*\x+\r*\f}) coordinate  (B');
						\draw [samples at={\i},gray,dash pattern=on 1pt off 2pt]  plot ({-\bx/\ra*\x},{\ax*\ra*\x+\r*\f})--plot ({-\ax*\x},{-\bx*\x+\r*\f})--plot ({\bx/\ra*\x},{-\ax*\ra*\x+\r*\f});
						\path ($(B')!.5!(D')$) coordinate  (I');
					 }
					\fill  [cyan!90!white!60] (A)--(B)--(C)--(D)--cycle;
					\draw [smooth,samples=100] 
						plot [domain=-.8:2]  ({\ax*\x},{\bx*\x+\r*\f})
						plot [domain=2:-2] ({\bx/\ra*\x},{-\ax*\ra*\x+\r*\f}) 
						(D)--(A)--(B); 
					\draw [dash pattern=on 1pt off  2pt,samples=100,gray] plot [domain=-.8:-2] ({\ax*\x},{\bx*\x+\r*\f});
					\draw [dash pattern=on 1pt off 2pt,gray] (D)--(C)--(B)--cycle  (A)--(C)
						(D')--(B'); 
					 \path ($(I)!1.3!(B)$) coordinate (t);
					 \draw [->,dashed] (I)--(t) node [above] {$x$};
					\draw [->,dash pattern=on  1pt off 2pt] (I)--++($.8*\ymax*(j)$) node [left] {$y$};
					\foreach \x/\y in {A/-90,B/-90,C/180,D/180}{ 
						\draw (\x) circle (.5pt) ++(\y:0.2) node {$\x$}; 
					} 
					\draw (P) circle (.5pt)  node [above left] {$S$};
					\pic [draw,angle radius=0.1cm] {right angle=P--I'--D'}; 
					\pic  [draw,angle radius=0.1cm] {right angle=P--I--D};
					\draw (I)  node[below left]{$O$};
					 \draw (D') node[left]{$M$}
						(B') node[right]{$P$} 
						(0.5,0.6) node{$N$} 
						 (-0.6,1.5) node{$Q$};
				\end{tikzpicture}
		\end{center}
		Ta tính được  $OD=OB=\sqrt{2}$.\\
		Gọi phương trình của đường  thẳng $a$ là $y=ax^2+bx+c$,\ $(a<0)$.\\
		Ta có $a$  đi qua các điểm $D\left(-\sqrt{2};0\right)$, $B\left(\sqrt{2};0\right)$, $S\left(0;\dfrac{27}{20}\right)$.\\
		Suy ra ta  có hệ $\heva{&2a-\sqrt{2}b+c=0\\&2a+\sqrt{2}b+c=0\\& c=\dfrac{27}{20}} \Leftrightarrow\heva{&a=-\dfrac{27}{40}\\&b=0\\& c=\dfrac{27}{20}} \Rightarrow y=-\dfrac{27}{40}x^2+\dfrac{27}{20}$.\\
		Gọi $I(0;y)$, $y\in\left[0;\dfrac{27}{20}\right)$.\\ 
		 Mặt phẳng vuông góc $Oy$ tại $I$ cắt hình đã cho theo $1$ thiết diện là hình vuông $MNPQ$ có diện tích $S(y)$.\\
		 Theo giả thiết trên các điểm $M$, $N$, $P$, $Q$ cùng có tung độ bằng $y$.\\
		Mà  hai điểm $M$, $P$ thuộc đường $a$ có phương trình $y=-\dfrac{27}{40}x^2+\dfrac{27}{20}$.\\
		Suy  ra
		$$\begin{aligned}&x^2=\dfrac{54-40y}{27} \Leftrightarrow x=\pm \dfrac{\sqrt{54-40y}}{3\sqrt{3}}\\
		&\Rightarrow M\left(-\dfrac{\sqrt{54-40y}}{3\sqrt{3}};y\right),\  P\left(\dfrac{\sqrt{54-40y}}{3\sqrt{3}};y\right)\Rightarrow MP=\dfrac{2\sqrt{54-40y}}{3\sqrt{3}}\\&\Rightarrow MN=\dfrac{\sqrt{108-80y}}{3\sqrt{3}} \Rightarrow S(y)=MN^2=\dfrac{108-80y}{27}.
	\end{aligned}$$
	Suy ra thể tích chiếc lều là $V=\displaystyle\int_0^{\tfrac{27}{20}}S(y) \mathrm{\,d}y=\displaystyle\int_0^{\tfrac{27}{20}}\dfrac{108-80y}{27}\mathrm{\,d}y=2{,}7$  (m$^3$).
 }
\end{bt}

% nguồn: Phan2_Lop_12_Dai_so_Chuong_4
\begin{bt}%[2D4V3-5]%[ma:b9b44e87]
	\immini{Một ly trà sữa dạng hình nón cụt, có đường kính đáy ly $6$\,(cm), đường  kính miệng ly $9$\,(cm), chiều cao $13{,}4$\,(cm), ở miệng ly có sử dụng một nắp đậy có hình dạng nửa mặt cầu và ở đỉnh của nửa mặt cầu này có một hình tròn có đường kính $2$\,(cm) để cắm ống hút, mặt phẳng chứa hình tròn này song song với mặt phẳng chứa miệng ly (tham khảo hình vẽ sau). Tính thể tích bên trong của ly bao gồm cả thể tích của nắp (\textit{làm tròn kết quả đến hàng đơn vị}).}
	{ 
	\begin{tikzpicture}[scale=0.3]
			\def\Rtop{4.5} 
			\def\Rbottom{3} 
			\def\Rhole{1} 
			 \def\hcone{9.9}
			\def\Rcap{4.5} 
			\draw (-\Rbottom,0)  arc (180:360:\Rbottom cm and 0.6cm);
			 \draw[dashed] (-\Rbottom,0) arc (180:-360:\Rbottom cm and 0.6cm);
			\draw (-\Rbottom,0) --  (-\Rtop,\hcone);
			\draw (\Rbottom,0) --  (\Rtop,\hcone);
			\draw (-\Rtop,\hcone) arc  (180:360:\Rtop cm and 0.9cm);
			\draw [dashed](-\Rtop,\hcone) arc (180:0:\Rtop cm and 0.9cm); 
			\draw  (-\Rtop,\hcone) arc (180:104:\Rtop);
			 \draw (\Rtop,\hcone) arc (0:78:\Rtop);
			\path  (-1,\hcone+\Rcap-0.1) coordinate (top);
			 \draw ([xshift=-\Rhole]top) arc (180:360:\Rhole cm and 0.25cm);
			\path (-1,\hcone+\Rcap-0.1) coordinate  (top);
			\draw ([xshift=-\Rhole]top) arc (180:-360:\Rhole cm  and 0.25cm);
			\draw[<->] (\Rtop+0.3,0) -- (\Rtop+0.3,\hcone) node[midway,right] {13,4 cm}; 
			\draw[<->] (-\Rbottom,-0.7) --  (\Rbottom,-0.7) node[midway,below] {6 cm};
			\draw[<->] (-\Rtop,\hcone) -- (\Rtop,\hcone) node[midway,above]  {9 cm};
			\draw[<->] (-\Rhole,\hcone+\Rcap+0.5) --  (\Rhole,\hcone+\Rcap+0.5) node[midway,above] {2 cm};
		\end{tikzpicture}
	} 
	 \shortans{791}
	\loigiai{ 
		Chọn hệ trục $Oxy$ (đơn vị trên trục là centimet) với trục $Ox$ đi qua tâm của  hai đáy hình nón cụt và gốc tọa độ $O$ trùng với tâm của đáy lớn như hình vẽ trên.
		\begin{center}
			\begin{tikzpicture}[scale=0.5,rotate=-90]
					\draw[->] (5,9.9) -- (-5.5,9.9)  node[left] {\(y\)};
					\draw[->] (0,-1) -- (0,16) node[below right]  {\(x\)};
					\draw[dashed](0,0) node[below]{$A$}--(-3,0) node[above]{$B$} 
						(0,9.9) node[below left]{$O$} 
						(-4.5,9.9) node[above  right]{$C$}
						(-1,14.4)  node[right]{$D$}--(0,14.4);
					 \def\Rtop{4.5}
					 \def\Rbottom{3}
					\def\Rhole{1} 
					 \def\hcone{9.9}
					\def\Rcap{4.5} 
					\draw (-\Rbottom,0)  arc (180:360:\Rbottom cm and 0.6cm);
					 \draw[dashed] (-\Rbottom,0) arc (180:-360:\Rbottom cm and 0.6cm);
					\draw (-\Rbottom,0) -- (-\Rtop,\hcone); 
					\draw (\Rbottom,0) --  (\Rtop,\hcone);
					\draw (-\Rtop,\hcone) arc (180:360:\Rtop cm and  0.9cm);
					 \draw [dashed](-\Rtop,\hcone) arc (180:0:\Rtop cm and 0.9cm);
					 \draw (-\Rtop,\hcone) arc (180:104:\Rtop);
					\draw (\Rtop,\hcone)  arc (0:78:\Rtop);
					\path  (-1,\hcone+\Rcap-0.1) coordinate (top);
					\draw ([xshift=-\Rhole]top) arc  (180:360:\Rhole cm and 0.25cm);
					\path (-1,\hcone+\Rcap-0.1)  coordinate (top);
					\draw ([xshift=-\Rhole]top) arc  (180:-360:\Rhole cm and 0.25cm);
				\end{tikzpicture}
		\end{center}
		Khi đó toạ độ các điểm $A$,  $B$, $C$ lần lượt là $A(-13{,}4;0)$, $B(-13{,}4;3)$, $C(0;4{,}5)$.\\
		Gọi phương trình đường thẳng đi qua $BC$ có  dạng $y=ax+b$.\\
		Ta có hệ phương trình $\heva{&-13{,}4a+b=3\\&0a+b=4{,}5}  \Leftrightarrow\heva{&a=\dfrac{15}{134} \\&b=\dfrac{9}{2}.}$\\
		Phương trình đường thẳng  $BC$ là $y=\dfrac{15}{134} x+\dfrac{9}{2}$.\\
		Khi đó thể tích của phần thân ly trà sữa chưa  bao gồm nắp là
		 \[
		V_1=\pi \displaystyle\int_{-13{,}4}^0\left(\dfrac{15}{134}  x+\dfrac{9}{2}\right)^2\,\mathrm{d}x=\dfrac{3\,819\pi}{20}.
		\] 
		Điểm $C$, $D$ thuộc đường tròn  tâm $O(0;0)$, bán kính $R=4{,}5$ có phương trình $x^2+y^2=\dfrac{81}{4}$. \\
		Suy ra $y=\sqrt{\dfrac{81}{4}-x^2}$ $(y\ge  0)$.\\
		Thay $y=1$, ta được $x=\sqrt{\dfrac{81}{4}-1}=\dfrac{\sqrt{77}}{2} \Rightarrow D\left(\dfrac{\sqrt{77}}{2};  1\right)$. \\
		Khi đó thể  tích nắp của ly trà sữa là $V_2=\pi \displaystyle\int_0^{\frac{\sqrt{77}}{2}}\left(\sqrt{\dfrac{81}{4}-x^2}\right)^2\,\mathrm{d}x=\dfrac{83\pi \sqrt{77}}{12}$.\\
		Từ $(1)$ và $(2)$  ta suy ra thể tích bên trong của ly bao gồm cả thể tích của nắp là
		\[ 
		V=V_1+V_2=\dfrac{3819\pi}{20}+\dfrac{83\pi \sqrt{77}}{12} \approx 791\,\left(\mathrm{cm}^3\right). 
		\] 
	 }
\end{bt}

% nguồn: Phan2_Lop_12_Dai_so_Chuong_4
\begin{bt}%[2D4V3-5]%[ma:09d04bb5]
	Cho khối $V$  có đáy là miền phẳng bị giới hạn bởi:\\
	Đồ thị hàm số $y=3x+\dfrac{2}{x}$, $x>0$. Trục $Ox$ và hai đường thẳng $x=1$ và $x=2$. Khi đó,  mỗi tiết diện của khối $V$ cắt bởi mặt phẳng vuông góc với trục $Ox$ đều là một tam giác đều. Thể tích của khối $V$ bằng bao nhiêu? (làm tròn kết quả đến hàng đơn vị).
	 \shortans{15}
	 \loigiai{
		Diện tích của tam giác đều  là $S(x)=\dfrac{\sqrt{3}}{4}\left(9 x^2+12+\dfrac{4}{x^2}\right)$.\\
		 Vậy $V=\displaystyle\int_1^2 S(x) \mathrm{\,d}x
	=\displaystyle\int_1^2\dfrac{\sqrt{3}}{4}\left(9x^2+12+\dfrac{4}{x^2}\right) \mathrm{\,d}x=\dfrac{\sqrt{3}}{4}\left(3x^3+12  x-\dfrac{4}{x}\right)\bigg|_1^2 =\dfrac{35\sqrt{3}}{4}\approx 15$.}
\end{bt}

% nguồn: Phan2_Lop_12_Dai_so_Chuong_4
\begin{bt}%[2D4V3-5]%[ma:ff8e28bf]
	\immini[thm]{Trên cánh cửa được trang trí các núm có dạng một khối tròn xoay được đúc bằng chất liệu đồng, khuôn đúc của nó được tạo thành khi quay miền $(H)$ (phần được tô màu trong hình vẽ bên) quanh trục $AB$. Miền $(H)$ được giới hạn bởi  các cạnh $AB$, $AD$ của hình vuông $ABCD$ và các cung phần tư của các đường tròn bán kính bằng $1$\,cm với tâm lần lượt là trung điểm của các cạnh $AD$, $AB$.
	 }{
	\begin{tikzpicture}[scale=1,font=\footnotesize,line join=round,line cap=round,>=stealth]
			\path (-2,0)  coordinate (A)
				(2,0)  coordinate (B)
				(2,4) coordinate  (C)
				 (-2,4) coordinate (D)
				($(A)!0.3!(C)$)  coordinate (H);
			\fill[black!30] (A) -- (B) arc (0:90:2)  -- cycle;
			\fill[black!30] (D) arc  (90:0:2) -- (B) -- (A) -- cycle;
			\draw[teal] (A)--(B)--(C)--(D)--cycle; 
			\draw[teal]  (D) arc (90:0:2);
			\draw[teal]  (0,2) arc (90:0:2);
			\foreach \i/\g in  {A/180,B/0,C/0,D/180}
			\fill (\i) circle (1pt)+(\g:3mm) node {$\i$}; 
			\node at (H) {$(H)$}; 
		\end{tikzpicture}
	\begin{tikzpicture}[scale=0.6,font=\footnotesize,line join=round,line cap=round,>=stealth]
			\definecolor{wood}{RGB}{138,81,56} 
			 \definecolor{copper}{RGB}{184,115,51}
			 \tikzset{num_dong/.pic={
			\fill[top color= copper!50,bottom color=  copper] (0,0) circle (1.5);
			 \shade[ball color=copper,pos=0.2] (0,0) circle (1);}}
			 \clip[rounded corners=0.3mm] (-0.5,-0.5) rectangle (5.5,7.5);
			 \fill[black!10!wood] (-0.5,-0.5) rectangle (5.5,7.5);
			\draw[teal,black!25!wood] (-0.5,-0.5) --  (0,0)
				 (5,7)--++(0.5,0.5)
				(-0.5,7.5)--++(0.5,-0.5) 
				(5,0)--++(0.5,-0.5); 
			 \fill[wood!80!white] (0,0) rectangle (5,7);
			\foreach \i in  {1,2,3,4,5,6,7,8,9}{\foreach \j in {0.5,1.5,2.5,3.5,5,6.5} \path ({0.5*\i},{\j}) pic[scale=0.1]{num_dong};}
		\end{tikzpicture}
	} 
	\noindent Biết công thức tính khối lượng của một vật là $P=VD$, trong đó $P$ là khối lượng của vật (đơn vị g), $V$ là thể tích của vật (đơn vị cm$^3$) và $D$ là khối lượng riêng của vật (đơn vị g/cm$^3$), khối  lượng riêng của đồng là $D=8{,}96$\,g/cm$^3$. Giá đồng trên thị trường là $200\,000$ đồng/kg. Giá tiền vật liệu để đúc một núm đồng trên là bao nhiêu nghìn đồng? (làm tròn kết quả đến hàng đơn vị).
	\shortans{22} 
	 \loigiai{
		Chọn hệ trục $Oxy$ như  hình vẽ sau
		\begin{center}
			\begin{tikzpicture}[scale=1,font=\footnotesize,line join=round,line cap=round,>=stealth]
					\path  (0,0) coordinate (O)
						 (-2,0) coordinate (A)
						(2,0)  coordinate (B)
						(2,4) coordinate (C) 
						(-2,4) coordinate (D) 
						(-2,2) coordinate (I) 
						($(A)!0.3!(C)$)  coordinate (H);
					 \fill[black!30] (A) -- (B) arc (0:90:2) -- cycle;
					\fill[black!30] (D) arc (90:0:2) -- (B)  -- (A) -- cycle;
					\draw[teal] (B)--(C)--(D)--(A); 
					\draw[teal] (D) arc  (90:0:2);
					\draw[teal] (0,2) arc (90:0:2); 
					 \draw[teal,->] (-2.5,0)--(3,0) node[below]{$x$};
					\draw[teal,->]  (0,-0.5)--(0,4.5) node[left]{$y$};
					\foreach \i/\g in {A/145,B/45,C/0,D/180,O/-135,I/180} 
					\fill (\i) circle  (1pt)+(\g:3mm) node {$\i$};
					\fill (0,2) circle (1pt)+(-45:3mm) node  {$1$};
					\fill (0,4) circle (1pt)+(-45:3mm) node {$2$}; 
					\node at (H) {$(H)$}; 
					\node at (-0.4,3.7)  {$(C_1)$};
					\node at (1,2.1)  {$(C_2)$};
					\node  at ($(A)+(-90:3mm)$){$-1$};
					\node at  ($(B)+(-90:3mm)$){$1$};
				\end{tikzpicture}
		\end{center}
		Đường tròn $(C_1)$ có tâm $I(-1;1)$ và  bán kính $R_1=1$.\\
		Suy ra $(C_1)\colon (x+1)^2+(y-1)^2=1$. Do đó $y=1+\sqrt{1-(x+1)^2}$ (vì $y>1$).  \\
		Đường tròn $(C_2)$  có tâm $O(0;0)$ và bán kính $R_2=1$. \\
		Suy ra $(C_2)\colon x^2+y^2=1$. Do đó $y=\sqrt{1-x^2}$ (vì $y>0$).  \\
		Thể tích một núm đồng là  $V=\pi\displaystyle\int_{-1}^0\left[1+\sqrt{1-(x+1)^2}\right]^2\mathrm{d}x+\pi\displaystyle\int_0^1\left(\sqrt{1-x^2}\right)^2\mathrm{d}x\ \left(\text{cm}^3\right)$. \\
		Giá tiền vật liệu để đúc một  núm đồng trên là $V\cdot 8{,}96\cdot 10^{-3}\cdot 200\approx 22$ (nghìn đồng). \\
		Vậy giá  tiền vật liệu để đúc một núm đồng trên là $22$ nghìn đồng.
	 }
\end{bt}

% nguồn: Phan2_Lop_12_Dai_so_Chuong_4
\begin{bt}%[2D4V3-5]%[ma:03a2f6bd]
	Bình chứa nhiên liệu của một chiếc máy bay phản lực được mô tả trong không gian là một khối tròn xoay thu được khi quay hình phẳng giới hạn bởi đồ thị hàm số $y=\dfrac{1}{8}x^2\sqrt{2-x}$ với $(0\le x\le 2)$ và trục $Ox$ quanh trục $Ox$ ($x$ tính theo đơn  vị m). Thể tích của bình chứa bằng $\dfrac{a\pi}{b}$ m$^3$ với $\dfrac{a}{b}$ là phân số tối giản. Khi đó $a+b$ bằng
	\begin{center}
		\begin{tikzpicture}[line join=round,line cap=round,>=stealth,scale=2.5]
				\draw[->]  (-0.5,0)--(2.7,0) node[below left] {$x$};
				\draw[->] (0,-0.5)--(0,1) node[below left] {$y$}; 
				 \draw (0,0) node [below left] {$O$};
				\draw[smooth,samples=200,domain=0:2,variable=\x]plot (\x,{0.125*(\x)^(2)*sqrt  (2-(\x))});
				\draw[smooth,samples=200,domain=0:2,variable=\x]plot  (\x,{-0.125*(\x)^(2)*sqrt (2-(\x))});
				 \draw (1,1) node[above] {\tiny$\boxed{y=\dfrac{1}{8}x^2\sqrt{2-x}}$};
				\draw[->] (1.9,1)--(3,1.5); 
				\tikzset{maybay/.pic={ 

				\def\D{(-2,-.45)--(-2.36,.24)--(-2.2,.3)--(-1.35,-.25)--cycle;} 
				\draw  \D;
				\fill[brown!40!]  \D;

				\def\Q{(.15,-.1) 
					..controls +(-140:0.1) and  +(140:0.1).. (.2,-.35)--(.8,-.17)--(.7,.07)--cycle;}
				 \draw \Q;
				\draw[xshift=.5cm] \Q; 
				\fill[brown!40!,xshift=.5cm] \Q; 

				\draw[rotate=-70,xshift=-.23cm,yshift=1.27cm] (.7,-.1) ellipse (.12cm and 0.07cm); 

				 \def\T{(-2.1,-.5)
					..controls +(40:0) and +(-170:0.7).. (2,.74) 
					..controls +(-40:0)  and +(170:0.3).. (2.55,.7)
					..controls  +(-150:0) and +(30:.65).. (2.1,.3)
					 ..controls +(-150:0) and +(-18:1.25).. (-2.1,-.6)
					..controls +(90:0) and +(-90:0)..  (-2.1,-.51)
					 -- (-2.5,-.52)--cycle;
				} 
				\draw  \T;
				 \fill[brown!70!] \T;

				 \def\C{(.5,.05)
					 ..controls +(170:0) and +(-10:0).. (-1.48,.12)
					..controls +(-80:0)  and +(170:0.02).. (-1.65,.3)
					..controls +(180:0) and +(0:0).. (-1.77,.3) 
					..controls +(-80:0) and  +(110:0).. (-1.64,.12)
					..controls  +(-35:0) and +(145:0).. (-.2,-.17)
					 --cycle;}
				 \draw \C;
				\fill[brown!40!]  \C;

				\fill[brown!40!]  \Q;

				\draw[rotate=-70,xshift=-.4cm,yshift=.8cm]  (.7,-.1) ellipse (.12cm and 0.07cm);

				\draw (2.18,0.73)--(2.5,0.7)--(2.14,0.58)--cycle; 
				 \fill[brown!40!](2.18,0.73)--(2.51,0.71)--(2.14,0.57)--cycle;
					 }}
				\path (3,1.6)  pic[scale=0.4]{maybay};
			\end{tikzpicture}
	\end{center}
	\shortans{31} 
	 \loigiai{
		Thể tích của bình chứa  là
		 $$V=\displaystyle\int_0^2\pi\left(\dfrac{1}{8}x^2\sqrt{2-x}\right)^2\mathrm{\,d}x=\dfrac{\pi}{30}\ \left(\mathrm{m}^3\right).$$
		Do đó, $a=1$ và $b=30$ nên  $a+b=31$.
	 }
\end{bt}

% nguồn: Phan2_Lop_12_Dai_so_Chuong_4
\begin{bt}%[2D4V3-5]%[ma:37d7867a]
	Một ly thủy tinh có hình dạng phần chứa nước là một hình parabol tròn xoay. Hình dạng này được tạo ra bằng cách quay một phần của đường parabol quanh trục đối xứng của nó. Biết phần chứa nước của ly  có chiều cao tính từ đáy ly lên đến miệng ly là $10$ cm , đường kính miệng ly là $8$ cm (chỉ tính phần chứa nước, không tính phần thủy tinh).\\
	\begin{center}
		\begin{tikzpicture}[scale=.5,font=\footnotesize,line join=round,line cap=round,>=stealth]
				\draw[domain=-4:4,samples=1000]  plot (\x,{((\x)^2)/2});
				\draw (0,8) ellipse  (4cm and 2cm);
				\draw[dashed] (-4,8) node[above left]{$A$}  circle (1pt)--(4,8) node[right]{$B$} circle (1pt);
				 \draw (0,8) node[above]{$O$} circle (1pt);
				\draw[dashed] (0,0)  node[below]{$I$} circle (1pt)--(0,8);
				 \draw[dashed](-4,8)--(-5.5,8)
					 (0,0)--(-5.5,0);
				\draw[<->](-5,0)--(-5,8); \node[left,rotate=90] at (-5.6,6){$10$ cm}; 
				 \draw[<->](-4,10.5)--(0,10.5) node[above]{$8$ cm}--(4,10.5);
			\end{tikzpicture}
		\begin{tikzpicture}[scale=1,font=\footnotesize,line join=round,line cap=round,>=stealth]
				 \draw[domain=-2:2,samples=1000] plot (\x,{(2*(\x)^2)/2});
				\draw  (0,4) ellipse (2cm and 1cm);
				 \draw (0,2) ellipse (1.4cm and 0.5cm);
				\draw[fill=blue!50!white]  (0,-3) ellipse (2cm and 0.5cm);
				\draw  (0,-3) node[below]{$I$} circle (1pt)--(0,0);
			\end{tikzpicture}
		\begin{tikzpicture}[scale=1,font=\footnotesize,line join=round,line cap=round,>=stealth]
				 \draw[domain=-2:2,samples=1000] plot (\x,{(2*(\x)^2)/2});
				\draw (0,4)  ellipse (2cm and 1cm);
				\draw  (0,3) ellipse (1.7cm and 0.6cm);
				\draw[fill=blue!50!white] (0,-3) ellipse (2cm and 0.5cm); 
				\draw (0,-3)  node[below]{$I$} circle (1pt)--(0,0);
			\end{tikzpicture}
	\end{center}
	Ban đầu, người ta đổ vào ly một lượng nước có thể tích bằng $\dfrac{1}{4}$ thể tích của ly khi nó chứa đầy nước. Sau đó, người ta đổ thêm vào ly một lượng nước có thể tích bằng với lượng nước đã đổ ban đầu.  Hỏi sau khi đổ thêm, chiều cao của mực nước trong ly đã tăng thêm bao nhiêu centimet so với lúc ban đầu (làm tròn đến hàng phần trăm)?
	 \shortans{2,07}
	\loigiai{ 
		\begin{center}
			\begin{tikzpicture}[scale=1,font=\footnotesize,line join=round,line cap=round,>=stealth]
					\begin{scope}[shift={(6,0)}]
						\fill[blue!20,domain=0:4,variable=\x,samples=1000] 
						(0,0)  -- plot ({\x},{sqrt (\x)}) -- (4,0) -- cycle;
						\fill[blue!20,domain=0:4,variable=\x,samples=1000] 
						(0,0)  -- plot ({\x},{-sqrt (\x)}) -- (4,0) -- cycle;
						 \fill[yellow!20] (2,0) ellipse (0.35cm and 1.41cm);
						\fill[yellow!20] (1,0) ellipse (0.15cm and  1cm);
						\fill[yellow!20] (4,0) ellipse  (0.25cm and 2cm);
						 \draw[->] (2,0) -- (4.5,0) node[right] {$x$};
						\foreach \x/\xtext in {1/,2/} 
						\draw (\x  cm,1pt) -- (\x cm,-1pt) node[anchor=north] {$\xtext$};
						\draw[->] (0,-2.5) -- (0,2.5)  node[above] {$y$};
						\foreach \y/\ytext in {} 
						\draw (1pt,\y cm) -- (-1pt,\y  cm) node[anchor=east] {$\ytext$};
						\draw[domain=0:4,smooth,variable=\x,blue,samples=1000] plot ({\x},{sqrt (\x)}); 
						\draw[domain=0:4,smooth,variable=\x,blue,samples=1000] plot  ({\x},{-sqrt (\x)});
						\draw (2,0) ellipse (0.35cm and  1.41cm);
						\draw (4,0) ellipse (0.25cm  and 2cm);
						\draw  (2,1.41)--(2,-1.41)
						(-0.5,0)--(0,0); 
						\draw[dashed] (0,0)--(2,0); 
						\fill (0,0) circle (1pt) node[below left]{$O$}; 
						\draw[dashed,rotate=-90] (1,1) arc (0:180:1cm and  0.15cm);
						\draw[rotate=90] (1,-1) arc  (0:180:1cm and 0.15cm);
						\draw[dashed] (4,-2)--(0,-2) node [left]  {$-4$};
						\draw[dashed] (4,2)--(0,2) node [left] {$4$}; 
					\end{scope}
				\end{tikzpicture}
		\end{center}
		Ta có $(P): y^{2}=2 p  x$\\
		$x=10$, $y=4 \Rightarrow p=\dfrac{4}{5} \Rightarrow(P): y^{2}=\dfrac{8}{5}  x$\\
		Thể tích nước trong ly $V=\pi \displaystyle\int_{0}^{10} y^{2} \mathrm{\,d} x=\pi \displaystyle\int_{0}^{10} \dfrac{8}{5} x \mathrm{\,d} x=\left.\dfrac{4  \pi x^{2}}{5}\right|_{0}^{10}=80 \pi$.\\
		Chiều  cao của nước ban đầu $\pi \displaystyle\int_{0}^{h_{1}} \dfrac{8}{5} x \mathrm{\,d} x=\dfrac{1}{4} \cdot 80 \pi \Leftrightarrow \dfrac{4 h_{1}^{2}}{5}=20 \Leftrightarrow h_{1}=5$.\\
		Chiều cao của mực nước  sau khi đổ thêm $\pi \displaystyle\int_{0}^{h_{2}} \dfrac{8}{5} x \mathrm{\,d} x=2 \cdot \dfrac{1}{4} \cdot 80 \pi \Leftrightarrow \dfrac{4 h_{2}^{2}}{5}=40 \Leftrightarrow h_{2}=5 \sqrt{2}$.\\
		Chiều cao  của mực nước tăng thêm so với ban đầu: $h=h_{2}-h_{1}=5 \sqrt{2}-5 \approx 2{,}07$ (cm).\\
	 }
\end{bt}

% nguồn: Phan2_Lop_12_Dai_so_Chuong_4
\begin{bt}%[2D4V3-5]%[ma:67709e99]
	\immini{Một khối tròn xoay được tạo thành khi quay hình phẳng $(H)$ (phần tô đậm trong hình vẽ) quanh trục $AB$. Miền $(H)$ được giới hạn bởi đường tròn đường kính $AB$ và cung tròn tâm $A$. Biết $AB=8$ (cm) và điểm $K$ trong hình vẽ thỏa mãn $AK=3$ (cm). Hỏi thể tích của khối tròn xoay đó bằng bao nhiêu cm$^3$ (\textit{làm tròn  kết quả đến hàng đơn vị})?}
	{ 
	\begin{tikzpicture}[
				 scale=1,
				 point/.style={circle,fill=red,inner sep=1.5pt}]
			 \def\R{1.5}
			 \def\xK{-0.5}
			\coordinate (O) at (0,0); 
			\coordinate (A) at  (-\R,0);
			\coordinate (B) at (\R,0); 
			\coordinate (K)  at (\xK,0);
			\pgfmathsetmacro{\yC}{sqrt (\R*\R - \xK*\xK)} 
			\coordinate  (C) at (\xK,\yC);
			\coordinate (D) at (\xK,-\yC); 
			\pgfmathsetmacro{\startAngleO}{atan2  (\yC,\xK)}
		\pgfmathsetmacro{\endAngleO}{atan2 (-\yC,\xK)}
		\pgfmathsetmacro{\RarcA}{veclen  (\xK+\R,\yC)}
		 \pgfmathsetmacro{\startAngleA}{atan2 (\yC,\xK+\R)}
	\pgfmathsetmacro{\endAngleA}{atan2 (-\yC,\xK+\R)}
	\fill[cyan,fill opacity=0.7] (C) arc (\startAngleO:\endAngleO:\R) -- (D) arc (\endAngleA:\startAngleA:\RarcA) -- cycle;
	\draw[green!50!black] (O) circle (\R); 
	\draw[green!50!black] (C) arc (\startAngleA:\endAngleA:\RarcA);
	 \draw[blue!50!black] (A) -- (B);
	 \draw[blue!50!black,dashed] (C) -- (D);
	\node[point,label=left:$A$] at (A) {}; 
	\node[point,label=right:$B$] at (B)  {};
	 \node[point,label=left:$C$] at (C) {};
	\node[point,label={[label distance=-2pt]240:$K$}] at (K) {}; 
	\end{tikzpicture}
	} 
	\shortans{135} 
	 \loigiai{
	\immini{Chọn hệ trục tọa độ $Oxy$ như  hình vẽ với $A(-4;0)$, $B(4;0)$, $K(-1;0)$.\\
		Gọi $C$ là điểm trên cung  tròn tâm $A$ sao cho $KC \perp AB$, $OK=1$.\\
		Khi  đó
		\begin{eqnarray*}
			CK&=&\sqrt{OC^2-OK^2}=\sqrt{16-1}=\sqrt{15}\\ \Rightarrow  AC&=&\sqrt{AK^2+CK^2}=\sqrt{3^2+15}=
			 \sqrt{24}=2\sqrt{6}.
		\end{eqnarray*}
		Ta có $OD = 2\sqrt{6}-4 \Rightarrow  D\left(2\sqrt{6}-4; 0\right)$.\\
		}{ 
		\begin{tikzpicture}[
					 scale=.6,
					point/.style={circle,fill=red,inner  sep=1.5pt}]
				 \def\R{4}
				\def\xK{-1} 
				\coordinate (O) at (0,0); 
				\coordinate (A) at  (-\R,0);
				\coordinate  (B) at (\R,0);
				\coordinate (K)  at (\xK,0);
				\pgfmathsetmacro{\yC}{sqrt (\R*\R  - \xK*\xK)}
				\coordinate (C) at (\xK,\yC); 
				\coordinate (Cprime) at  (\xK,-\yC);
				\pgfmathsetmacro{\startAngleO}{atan2 (\yC,\xK)} 
			\pgfmathsetmacro{\endAngleO}{atan2 (-\yC,\xK)}
			 \pgfmathsetmacro{\RarcA}{veclen (\xK+\R,\yC)}
			\pgfmathsetmacro{\startAngleA}{atan2 (\yC,\xK+\R)} 
		\pgfmathsetmacro{\endAngleA}{atan2 (-\yC,\xK+\R)}
		\coordinate (D) at  ({-\R+\RarcA},0);
		\fill[cyan,fill opacity=0.7] (C) arc (\startAngleO:\endAngleO:\R) -- (Cprime) arc (\endAngleA:\startAngleA:\RarcA) -- cycle;
		\draw[->] (-5.5,0) -- (5.5,0)  node[below] {$x$};
		\draw[->]  (0,-4.5) -- (0,4.5) node[left] {$y$};
		\draw[green!50!black]  (O) circle (\R);
		\draw[green!50!black] (C) arc (\startAngleA:\endAngleA:\RarcA);
		\draw[blue!50!black]  (A) -- (B);
		\draw[blue!50!black,dashed] (C) -- (Cprime); 
		\node[point,label=above left:$A$]  at (A) {};
		 \node[point,label=above left:$B$] at (B) {};
		 \node[point,label=above left:$K$] at (K) {};
		 \node[point,label=above left:$C$] at (C) {};
		\node[point,label=above:$D$] at (D) {}; 
		\node[point,label=above left:$O$] at  (O) {};
		 \foreach \val in {-4,-1,4}{
			\draw (\val,2pt) -- (\val,-2pt) node[below=3pt]  {$\val$};
		} 
		\end{tikzpicture}
		 }
		\noindent 
		Phương trình đường tròn đường kính $AB$ có dạng $x^2+y^2=16$.\\ Suy ra nửa  trên trục $Ox$ có phương trình dạng $y=\sqrt{16-x^2}$.\\
		Phương trình  đường tròn tâm $A$, bán kính $AC$ có dạng $\left(x+4\right)^2+y^2=24 \Rightarrow y = \sqrt{24-\left(x+4\right)^2}$.\\
		Thể tích khối tròn xoay được tạo thành  là
		\[V=\pi \left[\displaystyle\int_{-1}^{4} \left(\sqrt{16-x^2}\right)^2\,\mathrm{d}x- \displaystyle\int_{-1}^{2\sqrt{6}-4}  \left(\sqrt{24-(x+4)^2}\right)^2\,\mathrm{d}x\right] \approx 135\,\,(\text{cm}^3).\]
	 }
\end{bt}

% nguồn: Phan2_Lop_12_Dai_so_Chuong_4
\begin{bt}%[2D4V3-5]%[ma:43672b29]
	\immini[thm]{Bạn An thực hiện thiết kế một logo hình con mắt cho một phòng khám nhãn khoa. Logo là một hình phẳng giới hạn bởi hai parabol $y=f(x)$ và $y=g(x)$ có các kích thước như hình vẽ dưới đây (phần được tô màu đen) và một hình tròn có bán kính bằng $0{,}5$\,dm  ở giữa là phần con ngươi (phần được tô màu xanh), đơn vị trên mỗi trục tọa độ là decimet. Biết rằng chi phí để sơn phần con ngươi hình tròn màu xanh là $25\,000$ đồng /dm$^2$ và chi phí để sơn phần còn lại màu đen là $20\,000$ đồng/dm$^2$. Chi phí để sơn logo trên là bao nhiêu nghìn đồng? (Kết quả làm tròn đến hàng đơn vị).
	 }{
	\begin{tikzpicture}[line join=round,line cap=round,>=stealth]
			 \tikzset{every node/.style={scale=0.9}}
			\begin{scope}
				\clip (-3,-1.5) rectangle  (3,2.5);
				 \fill[gray!35]plot[samples=200,domain=-2.449489743:2.449489743,smooth,variable=\x] (\x,{-0.25*(\x)^2+2})--plot[samples=200,domain=2.449489743:-2.449489743,smooth,variable=\x] (\x,{0.25*(\x)^2-1})--cycle;
				\draw plot[samples=200,domain=-2.449489743:2.449489743,smooth,variable=\x]  (\x,{-0.25*(\x)^2+2});
				 \draw plot[samples=200,domain=-2.449489743:2.449489743,smooth,variable=\x] (\x,{0.25*(\x)^2-1});
				\draw[fill=blue!35] (0,0.5) circle  (0.5);
			\end{scope}
			\draw[->]  (-3.1,0)--(3.1,0) node[below left] {$x$};
			\draw[->] (0,-1.6)--(0,3) node[below left]  {$y$};
			\draw (0,0) node [below  left] {$O$};
			\draw (-2,1pt)--(-2,-1pt) node [below left]  {$-2$};
			\draw (2,1pt)--(2,-1pt)  node [below] {$2$};
			\draw (1pt,-1)--(-1pt,-1)  node [below left] {$-1$};
			\draw (1pt,1)--(-1pt,1) node [above  left] {$1$};
			\draw (1pt,2)--(-1pt,2) node [above left] {$2$}; 
			\draw[dashed](-2,0)--(-2,1)--(0,1) 
				(2,0)--(2,1)--(0,1); 
			\draw (-2,1) circle  (1.0pt)
				(2,1) circle (1.0pt) 
				 (0,0.5) circle (1.0pt);
		\end{tikzpicture}
	 }
	 \shortans{200}
	 \loigiai{
		Gọi phương trình parabol $y=f(x)$ là $f(x)=ax^2+bx+c$ với $a\neq 0$ đi qua các điểm có tọa độ lần  lượt là $(2;1)$, $(0;2)$, $(-2;1)$.\\
		 Lập hệ phương trình ba ẩn, từ đó suy ra được $\heva{&a=-\dfrac{1}{4}\\&b=0\\&c=2}$ nên $f(x)=-\dfrac{1}{4}x^2+2$.\\
		Gọi phương trình parabol $y=g(x)$ là $g(x)=mx^2+n x+p$ với $m \neq 0$  đi qua các điểm có tọa độ lần lượt là $(-2;0)$, $(0;-1)$, $(2;0)$.\\
		Lập hệ phương trình ba ẩn, từ đó  suy ra được $\heva{&m=\dfrac{1}{4}\\&n=0\\&p=-1}$ nên $g(x)=\dfrac{1}{4}x^2-1$.\\
		Phương trình  hoành độ giao điểm
		 \allowdisplaybreaks
		\begin{eqnarray*}
			 f(x)=g(x)&\Leftrightarrow&-\dfrac{1}{4}x^2+2=\dfrac{1}{4}x^2-1\\
			 &\Leftrightarrow&\dfrac{1}{2}x^2=3\\
			&\Leftrightarrow&x=\pm \sqrt  6.
		\end{eqnarray*}
		 Diện tích hình tròn được sơn màu xanh là \[S_1=\pi r^2=\pi\cdot 0{,}5^2=0{,}25 \pi\,(\text{dm}^2).\]
		Diện tích hình phẳng giới hạn bởi hai  đường thẳng $y=f(x)$ và $y=g(x)$ là
		\allowdisplaybreaks 
		\begin{eqnarray*}
			 S_2&=&\displaystyle\int_{-\sqrt 6}^{\sqrt 6}|f(x)-g(x)|\mathrm{\,d}x\\
			&=&\displaystyle\int_{-\sqrt  6}^{\sqrt 6}\left|\left(-\dfrac{1}{4}x^2+2\right)-\left(\dfrac{1}{4}x^2-1\right)\right|\mathrm{\,d}x.
		\end{eqnarray*}
		Diện tích phần  sơn được tô màu đen là
		 \allowdisplaybreaks
		\begin{eqnarray*}
			S_3&=&S_2-S_1\\ 
			&=&\displaystyle\int_{-\sqrt 6}^{\sqrt 6}\left|\left(-\dfrac{1}{4}x^2+2\right)-\left(\dfrac{1}{4}x^2-1\right)\right|  \mathrm{\,d}x-0{,}25 \pi.
		\end{eqnarray*}
		 Chi phí để sơn logo là
		 \allowdisplaybreaks
		\begin{eqnarray*}
			T&=&0{,}25 \pi \cdot 25\,000+\left(\displaystyle\int_{-\sqrt 6}^{\sqrt 6}\left|\left(-\dfrac{1}{4}x^2+2\right)-\left(\dfrac{1}{4}x^2-1\right)\right|  \mathrm{\,d}x-0{,}25 \pi\right)\cdot 20\,000\\
			&\approx& 200\,\text{ (nghìn  đồng).}
		\end{eqnarray*}
	} 
\end{bt}

% nguồn: Phan2_Lop_12_Dai_so_Chuong_4
\begin{bt}%[2D4V3-5]%[ma:373a26fc]
	\immini[thm]{Một hoa văn hình tròn tâm $O$, ngoại tiếp tam giác đều $ABC$ có cạnh $AB=2\sqrt 3$\,cm. Đường cong qua ba điểm $A$, $B$, $C$ là một phần của parabol (xem hình vẽ). Tính diện tích của phần hình phẳng giới hạn bởi đường tròn và parabol (phần không gạch)  theo đơn vị cm$^2$ (kết quả làm tròn đến chữ số thập phân thứ hai).
	}{ 
	\begin{tikzpicture}[line join=round,line cap=round,>=stealth]
			\tikzset{every node/.style={scale=0.9}} 
			\begin{scope}
				\clip (-2,-1)  rectangle (2,4);
				\fill[pattern=north east lines,pattern color=gray]plot[samples=200,domain=-1.732050808:1.732050808,smooth,variable=\x] (\x,{-1*(\x)^2+3})--plot[samples=200,domain=1.732050808:-1.732050808,smooth,variable=\x] (\x,{1-sqrt (4-(\x)^2)})--cycle; 
				\draw(-1.732050808,0) circle (1.0pt)  node[below]{$A$}
					(1.732050808,0) circle  (1.0pt) node[below]{$B$}
					(0,3) circle  (1.0pt) node[above left]{$C$};
				\draw plot[samples=200,domain=-1.732050808:1.732050808,smooth,variable=\x]  (\x,{-1*(\x)^2+3});
				\draw  (0,1) circle (2);
				 \draw (0,1) circle (1.0pt) node[above left]{$O$};
				\draw (-1.732050808,0)--(0,3)--(1.732050808,0)--cycle  (-1.732050808,0)--(0,1)--(1.732050808,0);
			\end{scope}
		\end{tikzpicture}
	} 
	\shortans{3,18} 
	 \loigiai{
		Gắn hệ trục tọa độ $Oxy$ như  hình vẽ.\\
		\begin{center}
			\begin{tikzpicture}[line join=round,line cap=round,>=stealth]
					\tikzset{every node/.style={scale=0.9}} 
					\begin{scope}
						\clip  (-2,-1) rectangle (2,4);
						\fill[pattern=north east lines,pattern color=gray]plot[samples=200,domain=-1.732050808:1.732050808,smooth,variable=\x] (\x,{-1*(\x)^2+3})--plot[samples=200,domain=1.732050808:-1.732050808,smooth,variable=\x] (\x,{1-sqrt  (4-(\x)^2)})--cycle;
						\draw(-1.732050808,0)  circle (1.0pt) node[below]{$A$}
						(1.732050808,0) circle (1.0pt) node[below]{$B$} 
						(0,3) circle  (1.0pt) node[above left]{$C$};
						\draw plot[samples=200,domain=-1.732050808:1.732050808,smooth,variable=\x]  (\x,{-1*(\x)^2+3});
						\draw (0,1)  circle (2);
						\draw (0,1) circle  (1.0pt) node[above left]{$O$};
						\draw  (-1.732050808,0)--(0,3)--(1.732050808,0)
						 (-1.732050808,0)--(0,1)--(1.732050808,0);
					\end{scope}
					\draw[->]  (-2.1,0)--(2.1,0) node[below] {$x$};
					\draw[->] (0,-1.1)--(0,4.1) node[below left]  {$y$};
					\draw (0,0) node [below left] {$O'$}; 
				\end{tikzpicture}
		\end{center}
		Khi đó $A(-\sqrt 3;0)$, $B(\sqrt 3;0)$, $C(0;3)$.\\ 
		Parabol đi qua ba điểm $A(-\sqrt 3;0)$, $B(\sqrt 3;0)$, $C(0;3)$ nên  parabol có phương trình là \[y=-x^2+3.\]
		Đường tròn ngoại tiếp tam giác $ABC$ có tâm $O(0;1)$ và bán kính $R=2$ nên có  phương trình là \[x^2+(y-1)^2=4.\]
		Suy ra $y=1-\sqrt{4-x^2}$ (phần nằm dưới trục hoành).\\ 
		Diện tích phần gạch  là $S=\displaystyle\int_{-\sqrt 3}^{\sqrt 3}-x^2+3-\left(1-\sqrt{4-x^2}\right)\mathrm{\,d}x$.\\
		Do đó diện tích  phần không gạch là $S'=\pi \cdot 2^2-S=3{,}18$\,(cm$^2$).
	 }
\end{bt}

% nguồn: Phan2_Lop_12_Dai_so_Chuong_4
\begin{bt}%[2D4V3-5]%[ma:5217af45]
	\immini{Cho các đồ thị hàm số $y=a^x$, $y=a^{-x}$, $y=Ax^2+Bx+C$, $1\neq a > 0$, $A$, $B$, $C\in \mathbb{R}$ và đường tròn $(C)$. Gọi $(D)$ là miền phẳng được tô đậm (hình vẽ). Quay miền $(D)$ quanh trục $Ox$  ta được một vật thể tròn xoay có thể tích bằng bao nhiêu (kết quả làm tròn đến hàng đơn vị)?
	\shortans{652}}{ 
	\begin{tikzpicture}[line cap=round,line join=round,font=\footnotesize,>=stealth,scale=0.7]
			\tikzset{label  style/.style={font=\footnotesize}}
			\pgfmathsetmacro{\a}{2} 
			\pgfmathsetmacro{\b}{0.5} 
			\draw[fill=blue!15](0,4) circle  (3);
			 \draw[->] (-3.5,0)--(4,0) node[below]{$x$};
			 \draw[->] (0,-2)--(0,8) node[left]{$y$};
			 \path
				(0,0)  coordinate (O)
				(0,4) coordinate  (I)
				(0,1) coordinate  (A);
			\draw[smooth,samples=100]  plot[domain=-3.5:3] (\x,{(\a)^(\x)});
			 \draw[smooth,samples=100] plot[domain=-3:3.5] (\x,{(\b)^(\x)});
			\draw[smooth,samples=100] plot[domain=-1.7:1.7] (\x,{-2*(\x)^2+4}); 
			\draw (I) let\p1=($(I)-(A)$) in circle  ({veclen (\x1,\y1});
			 \draw[fill=gray!1](-1,2)--plot[smooth,samples=200,domain=-1:0](\x,{-2*(\x)^2+4})--plot[smooth,samples=200,domain=0:-1](\x,{(\b)^(\x)});
			 \draw[fill=gray!1](0,4)--plot[smooth,samples=200,domain=0:1](\x,{-2*(\x)^2+4})--plot[smooth,samples=200,domain=1:0](\x,{(\a)^(\x)});
			\fill (0,2) circle  (1.5pt) node[left]{$2$}
				(0,4) circle  (1.5pt) node[above left]{$4$}
				(0,1) circle (1.5pt) node[below  left]{$1$}
				(-3,0)  circle (1.5pt) node[below]{$-3$}
				 (-2,0) circle (1.5pt) node[below]{$-2$}
				(-1,0) circle (1.5pt) node[below]{$-1$} 
				(1,0) circle (1.5pt)  node[below]{$1$}
				(2,0) circle (1.5pt)  node[below]{$2$}
				(3,0) circle (1.5pt)  node[below]{$3$}
				(1,2) circle (1.5pt); 
			\foreach\x/\y in {O/-135}  \draw (\x) circle (1pt)+(\y:0.3) node{$\x$};
			\draw [dashed]  (1,0)--(1,2)--(0,2);
			\draw  (0,1)--(0,4);
		\end{tikzpicture}
	 }
	\loigiai{ 
		Dựa vào các điểm đặc  biệt của đồ thị hàm số, ta tìm được:
		\begin{itemize}
			\item  Hàm số $y=a^x$ đi qua điểm $(1;2)$ nên $a=2$.
			\item Hàm số $y=Ax^2+Bx+C$ có đỉnh  là $(0;4)$ nên $C=4$, $B=0$; và có đồ thị đi qua điểm $(1;2)$ nên $A=-2$. Khi đó hàm số parabol là $y=-2x^2+4$.
			\item Đường tròn có tâm $(0;4)$ và bán kính $R=3$ nên có phương trình $(C)\colon x^2+(y-4)^2=9$, từ đây suy ra phương trình đường tròn này gồm hai  hàm số là $y=4+\sqrt{9-x^2}$ và $y=4-\sqrt{9-x^2}$.
		\end{itemize}
		Trước hết ta tính thể tích giới hạn bởi đường tròn $(C)$ khi  quay quanh trục $Ox$, tức là hình phẳng giới hạn bởi $y=4+\sqrt{9-x^2}$ và $y=4-\sqrt{9-x^2}$ khi quay quanh trục $Ox$ là
		 $$V_1=\pi \displaystyle\int_{-3}^3\left[\left(4+\sqrt{9-x^2}\right)^2-\left(4-\sqrt{9-x^2}\right)^2\right]\mathrm{\,d} x =72\pi^2.$$
		Tiếp theo ta tính thể tích khi  quay phần màu trắng giới hạn bởi ba hàm số $y=2^x$, $y=2^{-x}$, $y=-2x^2+4$ quanh trục $Ox$.\\
		Từ đồ thị trên ta thấy 
		$$V_2=2 \pi \displaystyle\int_0^1\left[\left(-2 x^2+4\right)^2-(2^x)^2\right] \mathrm{\,d} x\approx 58{,}45.$$ 
		Vậy thể tích cần tính là $V=V_1-V_2\approx 652$ (đvtt). 
	} 
\end{bt}

% nguồn: Phan2_Lop_12_Dai_so_Chuong_4
\begin{bt}%[2D4V3-5]%[ma:0ca9dba3]
	Một thùng rượu  (xét phần bên trong) có $2$ đáy là các hình tròn với bán kính là $30$ cm, thiết diện $(P)$ vuông góc với trục nối tâm của $2$ đáy và cách đều $2$ đáy có bán kính là $40$ cm (bên trong), chiều cao thùng rượu là $1$ m (hình vẽ). Biết rằng mặt phẳng chứa trục cắt mặt xung quanh (bên trong) thùng rượu theo các đường parabol có đỉnh nằm trên mặt phẳng $(P)$, hỏi dung tích của thùng rượu (đơn vị: lít) là bao nhiêu (làm tròn kết quả đến hàng đơn vị)?\\
	\begin{center}
		\begin{tikzpicture}[scale=0.4,font=\footnotesize,line join=round,line cap=round,>=stealth]
				\begin{scope}[rotate=90]
					\fill[orange!70!black] (-5,-3) arc (-90:-270: 0.4 and  3)--(-5,3)--plot[domain=-5:5](\x,{-1/25*(\x)^2+4})--(5,3) arc (90:-90: 0.4 and 3)--(5,-3)--plot[domain=5:-5](\x,{1/25*(\x)^2-4})--cycle;
					\fill[orange!50!black] (5,0)  ellipse (0.4 and 3);
					\fill[orange!80!black] (5,0) ellipse (0.3  and 2.8);
					\draw[color=orange!50!black,line width=3pt] (4.9,3)  arc (90:270: 0.4 and 3);
					\draw[color=orange!50!black,line  width=3pt] (4,3.3) arc (90:270: 0.4 and 3.3);
					\draw[color=orange!50!black,line width=3pt]  (2.5,3.7) arc (90:270: 0.4 and 3.7);
					\draw[color=orange!50!black,line width=3pt] (-2.5,3.7)  arc (90:270: 0.4 and 3.7);
					\draw[color=orange!50!black,line  width=3pt] (-4,3.3) arc (90:270: 0.4 and 3.3);
					\draw[color=orange!50!black,line width=3pt] (-4.9,3) arc (90:270: 0.4 and  3);
				\end{scope}
				\node[below] at  (current bounding box.south){$a)$};
			\end{tikzpicture}
	\end{center}
	\shortans{425} 
	\loigiai{ 
		\begin{center}
			\begin{tikzpicture}[scale=0.4,font=\footnotesize,line join=round,line cap=round,>=stealth]
					\begin{scope}[rotate=90]
						\fill[orange!70!black] (-5,-3) arc (-90:-270: 0.4 and 3)--(-5,3)--plot[domain=-5:5](\x,{-1/25*(\x)^2+4})--(5,3) arc  (90:-90: 0.4 and 3)--(5,-3)--plot[domain=5:-5](\x,{1/25*(\x)^2-4})--cycle;
						 \fill[orange!50!black] (5,0) ellipse (0.4 and 3);
						\fill[orange!80!black] (5,0) ellipse (0.3  and 2.8);
						\draw[color=orange!50!black,line width=3pt] (4.9,3) arc (90:270:  0.4 and 3);
						\draw[color=orange!50!black,line width=3pt] (4,3.3) arc (90:270: 0.4 and 3.3); 
						 \draw[color=orange!50!black,line width=3pt] (2.5,3.7) arc (90:270: 0.4 and 3.7);
						\draw[color=orange!50!black,line width=3pt] (-2.5,3.7) arc  (90:270: 0.4 and 3.7);
						\draw[color=orange!50!black,line width=3pt] (-4,3.3) arc (90:270: 0.4  and 3.3);
						\draw[color=orange!50!black,line width=3pt] (-4.9,3) arc (90:270: 0.4 and  3);
					\end{scope}
					\node[below] at (current bounding  box.south){$a)$};
				\end{tikzpicture}
			\begin{tikzpicture}[scale=0.4,font=\footnotesize,line join=round,line cap=round,>=stealth]
					 \fill[color=orange!30] (-5,-3)--(-5,3)--plot[domain=-5:5](\x,{-1/25*(\x)^2+4})--(5,3)--(5,-3)--plot[domain=5:-5](\x,{1/25*(\x)^2-4})--cycle;
					\draw[->] (-7,0) --(7,0) node[below]{$x$}; 
					\draw[->] (0,-5)  --(0,5) node[left]{$y$};
					\draw (0,0) node[above right]{$O$}; 
					\draw[smooth,samples=200,domain=-5:5]  plot (\x,{-1/25*(\x)^2+4});
					\draw[smooth,samples=200,domain=-5:5]  plot (\x,{1/25*(\x)^2-4});
					\draw  (5,0) ellipse (0.4 and 3);
					\draw (-5,3) arc (90:270: 0.4  and 3);
					\draw[dashed] (-5,-3) arc (-90:90: 0.4 and 3); 
					 \draw[dashed] (-5,-3)--(-5,3)--(5,3)--(5,-3)--cycle;
					\fill[color=black] 
						(-5,3) node[above left] {$A$} circle (1.25pt) 
						(0,4) node[above  right] {$B$} circle (1.25pt)
						(5,3)  node[above right] {$C$} circle (1.25pt)
						(2,3.8) node[above]  {$(P)$};
					\draw (-5,0)  node[below left]{$-0.5$}
						(5,0) node[below right]{$0.5$}; 
					\node[below] at (current bounding box.south){$b)$}; 
				\end{tikzpicture}
		\end{center}
		Gọi $(P)\colon y=a{x^2}+bx+c$ là parabol đi qua điểm $A\left(0{,}5; 0{,}3\right)$ và có đỉnh $S\left(0; 0{,}4\right)$ (hình vẽ). Khi đó, dung  tích thùng rượu bằng thể tích khối tròn xoay khi cho hình phẳng giới hạn bởi $(P)$, trục hoành và hai đường thẳng $x=0{,}5$; $x=-0{,}5$ quay quanh trục $Ox$.\\
		Tìm được $(P)\colon y=-\dfrac{2}{5}{x^2}+0{,}4$. Suy ra $V=\pi\displaystyle\int_{-0{,}5}^{0{,}5}{\left(-\dfrac{2}{5}{x^2}+0{,}4\right)^2\,\mathrm{d}x}=\dfrac{203\pi}{1\,500}$ m$^3\approx 425$  lít.
	} 
\end{bt}

% nguồn: Phan2_Lop_12_Dai_so_Chuong_4
\begin{bt}%[2D4V3-5]%[ma:1cd67678]
	Có một vật thể là hình tròn xoay có dạng giống như một cái ly như hình vẽ dưới đây. Người ta đo được đường kính của miệng ly là $4$ cm và chiều cao là $6$ cm. Biết rằng thiết diện của chiếc ly  cắt bởi mặt phẳng đối xứng là một parabol. Tính thể tích $V$ (cm$^3$) của vật thể đã cho.\\
	\begin{center}
		\begin{tikzpicture}[scale=0.8,rotate=180]
				 \draw[domain=-2:2,smooth,variable=\x]
					plot ({\x},{-1.5*(\x)^2 +  6});
				\draw[dashed]  (-2,0) arc (180:360:2 cm and 0.6cm);
				\draw (-2,0) arc (180:-360:2  cm and 0.6cm);
				\filldraw  (-2,0) circle (1pt) node[above] {$A$};
				\filldraw (2,0) circle (1pt) node[above] {$B$}; 
				 \filldraw (0,6) circle (1pt) node[below] {$I$};
				\draw[dashed] (-2,0) -- (2,0); 
				\draw[dashed]  (0,6) -- (0,0);
				 \draw (0,6)--(3,6)
					 (2,0)--(3,0);
				\draw[<->] (3,0) -- (3,6); 
				 \node at (3.7,3) {6 cm};
				\draw[<->] (-2,-1) -- (2,-1); 
				\node at (0,-1.3) {4  cm};
			\end{tikzpicture}
	\end{center}
	 \shortans{37,7}
	\loigiai{ 
		\begin{center}
			\begin{tikzpicture}[scale=0.8,rotate=90]
					 \draw[domain=-2:2,smooth,variable=\x]
						plot ({\x},{-1.5*(\x)^2  + 6});
					 \draw[dashed] (-2,0) arc (180:360:2 cm and 0.6cm);
					\draw (-2,0)  arc (180:-360:2 cm and 0.6cm);
					\filldraw  (-2,0) circle (1pt) node[below] {$A$};
					\filldraw (2,0) circle (1pt)  node[above] {$B$};
					\filldraw (0,6) circle (1pt) node[below  left] {$O$};
					\draw  (0,6.5)--(0,6);
					\draw[dashed](0,6)--(0,0); 
					 \draw[->](0,0)--(0,-1.5) node[above]{$x$};
					\draw[->](-2,6)--(2,6) node[above]{$y$}; 
					\draw (2,0) --  (0,0);
					\draw (0,0) node[below right]{$6$}; 
				\end{tikzpicture}
		\end{center}
		Chọn hệ trục tọa độ $Oxy$ thích hợp, ta được  phương trình parabol là $(P):{y^2}=\dfrac{2}{3}x$ với $0\le x\le 6$\\
		Thể tích  vật thể đã cho là $V=\pi\displaystyle\int_0^6\dfrac{2}{3}x\,\mathrm{d}x=12\pi\approx 37{,}7$ cm$^3$.
	} 
\end{bt}

% nguồn: Phan2_Lop_12_Dai_so_Chuong_4
\begin{bt}%[2D4V3-5]%[ma:1f65e07e]
	Một khối bê tông có chiều cao $2$ mét được đặt trên mặt đất. Nếu cắt khối bê tông bởi một mặt phẳng nằm ngang và cách mặt  đất $x$ mét $(0\le x\le 2)$ thì được mặt cắt là một hình chữ nhật có chiều dài $3$ m và chiều rộng bằng $(0{,}8)^x$ mét (tham khảo hình vẽ).
	\begin{center}
		\begin{tikzpicture}[line join=round,line cap=round,>=stealth,font=\footnotesize,scale=.9]
				\path 
					(1.46,6.34) coordinate  (A)
					(2.3,6.97) coordinate  (B)
					(6.87,5.6) coordinate  (C)
					(6,5) coordinate (D) 
					 (6.85,3.28) coordinate (A')
					(5.53,2.52) coordinate  (B')
					(0.99,3.86) coordinate (C') 
					(2.31,4.66) coordinate (D') 
					(0.19,1.66) coordinate  (K)
					(2.3,2.91) coordinate (L) 
					(6.85,1.55)  coordinate (M)
					(4.68,0.23)  coordinate (N)
					 (7.77,6.14) coordinate (P)
					(7.77,3.87) coordinate  (Q)
					(7.77,2.11) coordinate  (S)
					(7.44,5.95) coordinate (R) 
					(7.44,1.89) coordinate (T) 
					($(Q)!0.5!(S)$)  node[right]{$x$(m)}
					 ($(R)!0.4!(T)$) node[right]{$2$ m}
					 ($(D')!0.48!(A')$) node[above]{$3$ m};
				 \fill[pattern=north east lines,pattern color=black] (D')--(A')--(B')--(C')--cycle;
				\draw [dashed,semithick] (B)--(L) 
					(M)--(L)--(K) 
					(A')--(D')--(C') 
					 (C)--(P)
					(A')--(Q) 
					(M)--(S); 
				\draw [semithick]  (A)--(B)--(C)--(D)--cycle (A')--(B')--(C')
					(K)--(N)--(M) 
					(M)--(C); 
				 \draw[<->] (Q)--(S);
				\draw[<->]  (R)--(T);
				\draw (D)  to [out=-95,in=75] (B')
					(B') to  [out=-105,in=70] (N);
				\draw  (A) to [out=-95,in=75] (C')
					(C') to [out=-105,in=70] (K); 
				\draw[<-] (6.19,2.9)--(6.5,0.96) node[right,xshift=-1mm]  {\scriptsize $(0{,}8)^x$(m)};
			\end{tikzpicture}
	\end{center}
	Tính thể tích (đơn vị m$^3$) của  khối bê tông đó (làm tròn kết quả đến hàng phần trăm).
	\shortans{4,84} 
	\loigiai{ 
		\begin{center}
			\begin{tikzpicture}[line join=round,line cap=round,>=stealth,font=\footnotesize,scale=.9]
					\path 
						(1.46,6.34) coordinate  (A)
						 (2.3,6.97) coordinate (B)
						(6.87,5.6) coordinate  (C) node[right]{$2$}
						(6.87,6.6)  coordinate (C'') node[right]{$x$}
						(6,5) coordinate (D) 
						(6.85,3.28) coordinate  (A') node[right]{$x$}
						(5.53,2.52)  coordinate (B')
						(0.99,3.86) coordinate  (C')
						 (2.31,4.66) coordinate (D')
						 (0.19,1.66) coordinate (K)
						(2.3,2.91) coordinate  (L)
						(6.85,1.55) coordinate  (M) node[right]{$O$}
						(4.68,0.23) coordinate (N) 
						(7.77,6.14) coordinate (P) 
						 (7.77,3.87) coordinate (Q)
						(7.77,2.11) coordinate (S) 
						(7.44,5.95) coordinate (R) 
						(7.44,1.89) coordinate  (T)
						($(D')!0.48!(A')$) node[above]{$3$} 
						($(B')!0.48!(A')$)  node[below,rotate={atan (0.5)}]{$(0{,}8)^x$};
					\fill[pattern=north east lines,pattern  color=black] (D')--(A')--(B')--(C')--cycle;
					\draw[->] (C)--(C''); 
					\draw  [dashed,semithick] (B)--(L)
						(M)--(L)--(K) 
						 (A')--(D')--(C');
					\draw [semithick] (A)--(B)--(C)--(D)--cycle (A')--(B')--(C') 
						(K)--(N)--(M) 
						(M)--(C); 
					\foreach  \p/\r in {C/0,A'/0,M/0}
					\fill (\p) circle  (1.7pt);
					\draw (D) to [out=-95,in=75] (B') 
						(B')  to [out=-105,in=70] (N);
					\draw  (A) to [out=-95,in=75] (C')
						(C') to [out=-105,in=70]  (K);
				\end{tikzpicture}
		\end{center}
		Ta có diện tích của mặt cắt là  $S(x)=3\cdot 0,8^x\,\left( \text{ m}^2\right)$.\\
		Coi trục $Ox$  là trục thẳng đứng vuông góc với mặt đất.\\
		Khi đó, ta thấy khối bê tông bị giới hạn bởi hai mặt  phẳng vuông góc với trục $Ox$ tại các điểm $x=0;\,x=2$.\\
		$S(x)$ là diện tích thiết diện của vật thể  bị cắt bởi mặt phẳng vuông góc với trục $Ox$ tại điểm $x$, $0\le x\le 2$ và $S(x)$ là hàm số liên tục trên đoạn $[0; 2]$.\\
		Vậy $V=\displaystyle\int_0^2{S(x)\mathrm{\,d}x}=\displaystyle\int_0^2{3\cdot 0{,}8^x\mathrm{\,d}x}\approx 4{,}84$  m$^3$.
	 }
\end{bt}

% nguồn: Phan2_Lop_12_Dai_so_Chuong_4
\begin{bt}%[2D4C3-5]%[ma:6bf8e13e]
	\immini{ 
	Một kiến trúc sư chịu trách nhiệm thiết kế một tòa nhà cao $30$ (m). Mặt cắt ngang tại mọi độ cao, vuông góc với trục thẳng đứng, luôn là một hình vuông (xem hình vẽ). Mặt đáy tòa nhà là hình vuông có cạnh $L_{0}=26$ (m), mặt đỉnh là hình vuông có cạnh $L_{30}=20$ (m). Mặt cắt ngang tại vị trí hẹp nhất của tòa nhà là hình vuông có cạnh $L_{\min}=13{,}75$ (m). Mặt cắt của tòa nhà theo mặt phẳng đứng chứa đường chéo đáy có dạng là hình phẳng giới hạn bởi hai đường cong parabol đối xứng với nhau qua trục thẳng đứng đi  qua tâm đáy của tòa nhà. Tính thể tích của tòa nhà đó (\textit{làm tròn đến hàng đơn vị, đơn vị tính: mét khối}).
	 }{
	\begin{tikzpicture}[line join=round,line cap=round,scale=1,transform shape,yscale=.65,>=stealth]
			\clip (-3,-7) rectangle  (3,6.2);
			\tikzset{hinh/.pic={ 
			 \def\x{2.25}
			\def\y{4} 
			 \path
				(0,0) coordinate  (O)
				(-\x,\y)  coordinate (A)
				 (-\x,-\y) coordinate (M)
				(\x,\y) coordinate (C) 
				(\x,-\y) coordinate  (P)
				(-\x+1,\y-1.8) coordinate  (B)
				($(A)+(C)-(B)$)  coordinate (D)
				(-\x+1,-\y-1.8) coordinate  (N)
				($(M)+(P)-(N)$) coordinate (Q) 
				($(A)!.5!(C)$)  coordinate (G)
				($(M)!.5!(P)$) coordinate  (I);
			\draw (A)..controls  +(-70:3.5) and +(70:3.5)..(M)
				 coordinate[pos=1/2](O1);
			\draw (C)..controls +(-110:3.5) and  +(110:3.5)..(P)
				 coordinate[pos=1/2](O3);
			\path ($(O1)!5!(O3)$) coordinate (T); 
			\draw[name path=bn](B)..controls +(-70:3.7)  and +(50:1.3)..(N);
			 \draw[dashed](D)..controls +(-110:3) and +(120:2.5)..(Q)
				coordinate[pos=2.9/5](O4); 
			\path ($(O4)!3!(O)$) coordinate  (T);
			\path[name path=d]  (O4)--(T);
			\draw[name intersections ={of= bn and d}]  (intersection-1) coordinate (O2);
			 \draw (A)--(B)--(C)--(D)--cycle (M)--(N)--(P);
			\draw[dashed] (M)--(Q)--(P); 
				 }}
			\path (0,0)  pic[scale=1]{hinh};
		\end{tikzpicture}
	} 
	 \shortans{8976}
	\loigiai{ 
		Chọn hệ trục toạ  độ $Oxy$ như hình vẽ.\\
		Gọi $L(x)$ là hàm biến thiên của độ dài đường chéo mặt cắt của toà nhà tại  độ cao $x$.\\
		Theo đề ta có, $L(x)$ là một parabol đi qua ba điểm $\left(0;13\sqrt{2}\right)$, $\left(30;10\sqrt{2}\right)$, $\left(x_{0};\dfrac{55\sqrt{2}}{8}\right)$, trong đó $x_{0}$ là vị trí toàn nhà có  cạnh $L_{\min}=13{,}75$ (m).\\
		 Phương trình parabol có dạng $L(x)=a\left(x-x_{0}\right)^2+\dfrac{55\sqrt{2}}{8}$, với $0<x_{0}<30$.\\
	\immini{Ta có 
		\begin{eqnarray*}
			 &&\heva{&L(0)=a(0-x_0)^2+\dfrac{55\sqrt{2}}{8}=13\sqrt{2}\\&L(30)=a(30-x_0)^2+\dfrac{55\sqrt{2}}{8}=10\sqrt{2}}\\
			 &\Rightarrow&\heva{&ax_0^2=\dfrac{49\sqrt{2}}{8}\\&a(30-x_0)^2=\dfrac{25\sqrt{2}}{8}}\\
			&\Rightarrow&\dfrac{x_0^2}{(x_0-30)^2}=\dfrac{49}{25}\Rightarrow \hoac{&x_0=105\text{ (loại)}\\&x_0=17{,}5\text{  (nhận).}}\\
			&\Rightarrow&  a=\dfrac{49\sqrt{2}}{8x_0^2}=\dfrac{\sqrt{2}}{50}\\
			&\Rightarrow&L(x)=\dfrac{\sqrt{2}}{50}\left(x-17{,}5\right)^2+\dfrac{55\sqrt{2}}{8}. 
		\end{eqnarray*}}{	\begin{tikzpicture}[line join=round,line cap=round,scale=1,transform shape,yscale=.75,>=stealth]
		 \clip (-3.5,-7) rectangle (3.5,6.2);
		 \tikzset{hinh/.pic={
		\def\x{2.25} 
		\def\y{4} 
		\path 
			 (0,0) coordinate (I)
			(-\x,\y) coordinate (A) 
			(-\x,-\y)  coordinate (M)
			(\x,\y) coordinate  (C)
			(\x,-\y) coordinate  (P)
			(-\x+1,\y-1.8)  coordinate (B)
			($(A)+(C)-(B)$) coordinate  (D)
			 (-\x+1,-\y-1.8) coordinate (N)
			 ($(M)+(P)-(N)$) coordinate (Q)
			($(A)!.5!(C)$)  coordinate (G)
			($(M)!.5!(P)$)  coordinate (O);
		\draw (A)..controls +(-70:3.5) and +(70:3.5)..(M) 
			 coordinate[pos=1/2](O1);
		 \draw (C)..controls +(-110:3.5) and +(110:3.5)..(P)
			coordinate[pos=1/2](O3); 
		\path ($(O1)!5!(O3)$) coordinate (T); 
		\draw[name path=bn](B)..controls  +(-70:3.7) and +(50:1.3)..(N);
		\draw[dashed]  (D)..controls +(-110:3) and +(120:2.5)..(Q)
			coordinate[pos=2.9/5](O4); 
		\path ($(O4)!3!(I)$)  coordinate (T);
		\path[name path=d]  (O4)--(T);
		\draw[name intersections ={of= bn and d}] (intersection-1) coordinate (O2) circle  (1pt);
		\draw  (A)--(C)
			(B)--(D) 
			(A)--(B)--(C)--(D)--cycle  (M)--(N)--(P)
			 (O1)--(O2)--(O3);
		 \draw[dashed] (P)--(Q)--(M)--cycle (Q)--(N)
			 (O3)--(O4)--(O1)
			 (O2)--(O4)
			 (O1)--(O3);
		\draw[->](O)--++(-3,0) node[below]{$y$}; 
		 \draw[->](I)--++(0,5.5) node[right]{$x$};
		\draw  (O)--(I);
			}} 
		\path (0,0) pic[scale=1]{hinh}; 
	\end{tikzpicture}
		 }
		Diện  tích mặt cắt của vật thể là $S\left(x\right)=\left[\sqrt{2}L(x)\right]^2=2\left[\dfrac{\sqrt{2}}{50}\left(x-17{,}5\right)^2+\dfrac{55\sqrt{2}}{8}\right]^2$.\\
		Thể  tích của toà nhà là
		\[V=\displaystyle\int_{0}^{30}S(x)\,\mathrm{d}x=\displaystyle\int_{0}^{30}  2\left[\dfrac{\sqrt{2}}{50}\left(x-17{,}5\right)^2+\dfrac{55\sqrt{2}}{8}\right]^2\,\mathrm{d}x=8\,976\,(\mathrm{m}^{3}).\]
	 }
\end{bt}

% nguồn: Phan2_Lop_12_Dai_so_Chuong_4
\begin{bt}%[2D4C3-5]%[ma:809efb61]
	Xét một Lavabo (bồn rửa) làm bằng sứ đặc hình dạng là một nửa khối elip tròn xoay có thông số kĩ thuật mặt trên của Lavabo là: dài $\times$ rộng: $660 \times 380$ mm (tham khảo hình vẽ bên dưới), Lavabo có độ dày đều là $20$ mm. Thể tích chứa nước của  Lavabo bằng bao~nhiêu dm$^3$? (kết quả làm tròn đến hàng phần chục)
	 \shortans{18,8}
	\loigiai{ 
		Giả sử mặt trên của Lavabo được biểu diễn như hình vẽ bên dưới. Gọi hệ trục tọa độ $Oxy$ như hình vẽ. Gọi  $(E)$ là elip nhỏ bên trong.
		\begin{center}
			\resizebox{!}{5cm}{\begin{tikzpicture}[line join=round,line cap=round,>=stealth]
		\draw[->] (-6,0)--(0,0) node[below right]{$O$}--(6,0)  node[below]{$x$};
		\draw[->]  (0,-3)--(0,5) node[right]{$y$};
		\draw (0,0) ellipse ({4} and  {2});
		\draw (0,0) ellipse  ({3.7} and {1.7});
		 \draw (2,1) node{$(E)$};
		\draw [pattern = north west lines,thin,domain=-3.7:3.7,samples=100] plot (\x,{(2.89*(1 - (\x)^2/13.69))^0.5}); 
	\end{tikzpicture}}
		\end{center}
		Độ  dài trục lớn của $(E)$ là $2a=660-40=620$ (mm) $=\dfrac{31}{5}$ (dm), suy ra $a= \dfrac{31}{10}$ dm.\\
		Độ dài trục bé của $(E)$ là $2b=380-40=340$ (mm) $=\dfrac{17}{5}$ (dm), suy  ra $a= \dfrac{17}{10}$ dm.\\
		Vậy phương trình của  $(E)$ là
		$$ \dfrac{x^2}{\left( \dfrac{31}{10} \right)^2} + \dfrac{y^2}{\left( \dfrac{17}{10} \right)^2}  = 1 \Rightarrow y^2 = \dfrac{289}{100} \left( 1 - \dfrac{100x^2}{961} \right).$$
		Thể tích khối tròn xoay khi quay miền giới hạn bởi $(E)$, trục $Ox$ và $x=-\dfrac{31}{10}$,  $x=\dfrac{31}{10}$ (phần gạch chéo trong hình) quanh trục $Ox$ là
		$$ V  = \pi \displaystyle\int_{-\frac{31}{10}}^{\frac{31}{10}} \dfrac{289}{100} \left( 1 - \dfrac{100x^2}{961} \right) \mathrm{\,d}x
		=\dfrac{8959\pi}{750} \text{ (dm}^3).$$ 
		Vậy thể tích chứa nước của Lavabo là  $\dfrac{V}{2} \approx 18{,}8$ (dm$^3$).
	 }
\end{bt}

% nguồn: Phan2_Lop_12_Dai_so_Chuong_4
\begin{bt}%[2D4C3-5]%[ma:38f8e5b8]
	Cho một thanh gỗ hình trụ $A$ có đáy là hình tròn bán kính $1$, và một thanh gỗ hình lăng  trụ đứng $B$ có đáy là hình vuông cạnh $\sqrt{2}$ như hình vẽ.
	\begin{center}
		\begin{tikzpicture}[scale=1,font=\footnotesize,line join=round,line cap=round,>=stealth]
				 \path
					(0,0) coordinate  (O)
					(2,0) coordinate (A); 
				\draw  (A)--(O) node[midway,above]{$1$} circle (2cm);
			\end{tikzpicture}
		 \qquad
		\begin{tikzpicture}[scale=1,font=\footnotesize,line join=round,line cap=round,>=stealth]
				 \path
					(0,0) coordinate (O) 
					 (2,0) coordinate (A)
					 (0,-2) coordinate (B)
					(-2,0)  coordinate (C)
					(0,2) coordinate (D); 
				\draw (A)--(D)  node[midway,above]{$\sqrt{2}$}--(C)--(B)--cycle
					 (A)--(C)
					 (B)--(D);
				\draw  pic[draw,angle radius=0.2cm] {right angle = D--O--A};
			\end{tikzpicture}
	\end{center}
	Hai thanh gỗ có trục chính cắt nhau và tạo với nhau  một góc $30^{\circ}$.
	\begin{center}
		\begin{tikzpicture}[scale=0.8,every node/.style={scale=1}]
				 \def\rA{0.6}
				 \def\lA{5.5}
				\def\wA{0.2} 
				\shade[left color=gray!30,right color=gray!40] (-\lA,-\rA) rectangle  (\lA,\rA);
				\draw  (-\lA,\rA) -- (\lA,\rA);
				\draw (-\lA,-\rA)  -- (\lA,-\rA);
				\draw[fill=gray!30] (-\lA,0) ellipse  (0.2 and 0.6);
				 \coordinate (xa) at ($(\lA,0)+(0,0.6)$);
				\draw[fill=gray!50]  (xa) arc (90:-90:0.2 and 0.6);
				\begin{scope}[rotate=30,xshift=-1cm,yshift=0.3cm]
					\coordinate (TL) at (-3.5,0.7); 
					\coordinate  (TR) at (3.5,0.7);
					 \coordinate (ML) at (-3.2,0.7);
					\coordinate  (MR) at (4.8,0.7);
					\coordinate (BL)  at (-3.2,-0.2);
					\coordinate (BR) at  (4.8,-0.1);
					\coordinate  (EndLTop) at (-3.8,0.5);
					\coordinate (EndLBot) at (-3.8,-0.45); 
					\coordinate (1) at  ($(EndLBot)+(ML)-(BL)$);
					\coordinate (2) at  ($(EndLBot)+(BR)-(BL)$);
					\coordinate (x1) at ($(EndLBot)!0.5!(ML)$); 
					\coordinate (y1)  at ($(BR)!0.5!(MR)$);
					\fill[fill=gray!80,draw] (EndLBot) -- (BL) -- (BR) -- (MR)--(ML) --  (1) -- cycle (EndLBot) -- (BL) -- (BR) -- (2)--cycle;
					 \draw (BL)--(ML)
						 ($(y1)!1.25!(x1)$)--(x1)
						($(x1)!1.25!(y1)$)--(y1); 
					\draw[dashed] (x1)  -- (y1);
				\end{scope}
				\draw (-0.3,0) arc (0:30:0.8); 
				\node[above] at (0.1,0.2)  {$30^\circ$};
			\draw[dashed] (-6.5,0) -- (6.5,0) coordinate (axisA_end);
			\node at (6.0,-0.5)  {A};
		\draw (axisA_end) -- (7.5,0);
		\node[above] at  (7.6,0) {trục chính};
		\node at  (4.0,2.6) {B};
		\node[above] at (6.3,3.0)  {trục chính};
		\end{tikzpicture}
	\end{center}
	Trong hai đường chéo của tiết diện vuông góc với trục của thanh  gỗ $B$, có một đường chéo vuông góc với mặt phẳng chứa trục của cả hai thanh gỗ, còn đường chéo còn lại thì nằm trong mặt phẳng đó. Khi đó, phần thể tích chung của hai thanh gỗ là $a\pi+b$. Hãy tính giá trị của biểu thức $10 a+3 b$, biết rằng $a$ và $b$ là các số hữu tỉ.
	\shortans{24} 
	 \loigiai{
		Đầu tiên, ta sẽ thực hiện $1$ lát cắt ngang vuông  góc với \lq\lq đường chéo vuông góc với mặt phẳng chứa trục của cả hai thanh gỗ\rq\rq.
		\begin{center}
			\begin{tikzpicture}[scale=1,every node/.style={scale=1}]
					 \def\rA{0.6}
					 \def\lA{5.5}
					\def\wA{0.2} 
					\shade[left color=gray!30,right color=gray!40] (-\lA,-\rA)  rectangle (\lA,\rA);
					\draw (-\lA,\rA)  -- (\lA,\rA);
					\draw  (-\lA,-\rA) -- (\lA,-\rA);
					 \draw[fill=gray!30] (-\lA,0) ellipse (0.2 and 0.6);
					\coordinate (xa) at  ($(\lA,0)+(0,0.6)$);
					 \draw[fill=gray!50] (xa) arc (90:-90:0.2 and 0.6);
					\begin{scope}[rotate=30,xshift=-1cm,yshift=0.3cm]
						\coordinate  (TL) at (-3.5,0.7);
						\coordinate (TR) at  (3.5,0.7);
						\coordinate (ML) at (-3.2,0.7); 
						\coordinate  (MR) at (4.8,0.7);
						 \coordinate (BL) at (-3.2,-0.2);
						\coordinate (BR) at (4.8,-0.1); 
						 \coordinate (EndLTop) at (-3.8,0.5);
						 \coordinate (EndLBot) at (-3.8,-0.45);
						\coordinate (1) at  ($(EndLBot)+(ML)-(BL)$);
						\coordinate (2) at ($(EndLBot)+(BR)-(BL)$); 
						\coordinate (x1) at ($(EndLBot)!0.5!(ML)$); 
						\coordinate  (y1) at ($(BR)!0.5!(MR)$);
						\fill[fill=gray!80,draw] (EndLBot) -- (BL) -- (BR) -- (MR)--(ML) -- (1) -- cycle (EndLBot) -- (BL) --  (BR) -- (2)--cycle;
						\draw (BL)--(ML) 
						 ($(y1)!1.25!(x1)$)--(x1)
						($(x1)!1.25!(y1)$)--(y1); 
						\draw[dashed] (x1) -- (y1); 
					\end{scope}
					 \draw (-0.3,0) arc (0:30:0.8);
					\node[above]  at (0.1,0.2) {$30^\circ$};
				\draw[dashed] (-6.5,0) -- (6.5,0) coordinate (axisA_end);
				\node at (6.0,-0.5) {A}; 
			\draw (axisA_end) -- (7.5,0);
			 \node[above] at (7.6,0) {trục chính};
			\node at  (4.0,2.6) {B};
			 \node[above] at (6.3,3.0) {trục chính};
			\coordinate (B)  at (4,2.6);
			\coordinate  (w1) at (-6.6,1.4);
			\coordinate (w2) at (-2.4,-4.0); 
			\coordinate  (I1) at ($(w1)!0.5!(w2)$);
			\coordinate (w3)  at ($(w1)+(B)-(I1)$);
			 \coordinate (w4) at ($(w2)+(B)-(I1)$);
			\draw (w1)--(w2); 
			 \draw[->] (w1)--(w3);
			 \draw[->] (w2)--($(w2)!0.7!(w4)$);
			 \node at (-6.5,1)[left]{cắt};
			\end{tikzpicture}
		\end{center}
		Ta gọi $O$ là giao của $2$ trục chính.  Giả sử lát cắt ngang này, cắt trục $Ox$ tại điểm có hoành độ bằng $x$ với $x\in[0; 1]$.\\
		Khi đó ta được hình sau: 
		\begin{center}
			\begin{tikzpicture}[scale=0.8,font=\footnotesize,line join=round,line cap=round,>=stealth]
					 \path
						 (0,0) coordinate (O)
						(2,0) coordinate  (A)
						(5,0) coordinate (O1) 
						(7,0) coordinate  (A1)
						(5,-2) coordinate  (B)
						(3,0) coordinate (C) 
						(5,2)  coordinate (D);
					\draw (A)--(O)  node[midway,above]{$1$} circle (2cm);
					\draw[->] (O)--(90:2.5) node[left]{$x$}; 
					\draw (A1)--(D) node[midway,above]{$\sqrt{2}$}--(C)--(B)--cycle 
						 (A1)--(C)
						 (B)--(D);
					\draw pic[draw,angle radius=0.2cm]  {right angle = D--O1--A1};
					\draw[dashed] (-3,1) node[above]{Cắt ngang}--(7.5,1); 
					\node at  (O)[below left]{$O$};
					\node  at (0.8,1)[above]{$\sqrt{1-x^2}$};
					\node at  (O1)[below left]{$O$};
					\node at (6,1)[below left]{$1-x$}; 
					\draw[->] (O1)--($(O1)!1.2!(D)$) node[left]{$x$}; 
					\node  at (0,-4){};
				\end{tikzpicture}
			\quad 
			\begin{tikzpicture}[scale=0.8,font=\footnotesize,line join=round,line cap=round,>=stealth]
					\path 
						(0,-1.2) coordinate (x1) 
						(0,1.2) coordinate (x2) 
						(9,1.2) coordinate  (x3)
						(9,-1.2) coordinate  (x4)
						(0,-2.2)  coordinate (y2)
						(2,-4.2) coordinate (y1) 
						 (7,4.2) coordinate (y3)
						($(y1)+(y3)-(y2)$) coordinate (y4) 
						(intersection of  x1--x4 and y2--y3) coordinate (B)
						(intersection of x1--x4 and y1--y4) coordinate (C) 
						(intersection of x2--x3 and y2--y3) coordinate  (A)
						(intersection of x2--x3 and y1--y4)  coordinate (D);
					\draw[dashed]  (x1)--(x2)--(x3)--(x4)--cycle;
					\draw[dashed] (y1)--(y2)--(y3)--(y4)--cycle; 
					\foreach \p/\q in  {A/90,B/100,C/-80,D/90}
					\fill (\p)  circle (1.0pt)
						 ($(\p)+(\q:2.5mm)$) node{$\p$};
					\node at ($(D)!0.5!(C)$)[right]{$4\sqrt{1-x^2}$}; 
					\node  at ($(x1)!0.5!(x2)$)[right]{$2\sqrt{1-x^2}$};
					\fill[pattern=north  east lines,opacity=0.5] (A)--(B)--(C)--(D)--cycle;
					\node at ($(B)!0.5!(C)$)[above]{$4 (1-x)$}; 
				\end{tikzpicture}
		\end{center}
		Khi đó, ta cắt thanh gỗ $A$ có chiều rộng  là $2\sqrt{1-x^2}$ và thanh gỗ $B$ có chiều rộng là $2(1-x)$.\\
		Còn $ABCD$ là tứ giác là hình bình hành do lát cắt tạo thành vào  phần chung của $A$ và $B$.\\
		Dễ dàng tìm được diện tích tứ giác này  là $S(x)=\dfrac{1}{2}\cdot 4(1-x)\cdot 4\sqrt{1-x^2}$.\\
		Vậy thể  tích cần tìm là
		\begin{eqnarray*}
			V&=&2 \displaystyle\int_0^1\left[\dfrac{1}{2}  \cdot 4(1-x) \cdot 4 \sqrt{1-x^2}\right] \mathrm{\,d}x\\
			&=&16 \displaystyle\int_0^1\left[(1-x) \sqrt{1-x^2}\right] \mathrm{\,d}x \\ 
			&=&16\left(\displaystyle\int_0^1 \sqrt{1-x^2} \mathrm{\,d}x -\displaystyle\int_0^1 x \sqrt{1-x^2}  \mathrm{\,d} x\right)\\
			&=&16\left(\dfrac{\pi}{4}-\dfrac{1}{3}\right)=4 \pi-\dfrac{16}{3}\\ &\Rightarrow& \hoac{&a=4 \\&b=-\dfrac{16}{3}} \Rightarrow 10a+3b=24. 
		\end{eqnarray*}
	 }
\end{bt}

% nguồn: Phan2_Lop_12_Dai_so_Chuong_4
\begin{bt}%[2D4C3-5]%[ma:f48999d1]
	 \immini{
	Một cốc thuỷ tinh có dạng hình trụ đang đựng một lượng nước, bán kính trong lòng đáy cốc là $8$\,(cm) và chiều cao trong lòng cốc là $10$\,(cm). Tính thể tích lượng nước trong cốc (đơn vị cm$^3$),  biết khi nghiêng cốc nước vừa lúc khi nước chạm miệng cốc thì ở đáy cốc mực nước trùng với đường kính của đáy cốc (\textit{làm tròn kết quả đến hàng đơn vị}).
	 }
	{ 
	\begin{tikzpicture}[scale=0.6,>=stealth,line join=round,font=\scriptsize]
			 \path (0,0) coordinate (coc);
			\fill [gray!60]  (coc)++(-.79,-1.4)++(-30:3.28)
				.. controls +(-3.7,.2)  and +(-4.84,-.2).. cycle;
			 \draw [gray,fill=gray!5] (coc)++(-.79,-1.4) coordinate (A)--++(-30:3.28) coordinate (B)
				.. controls +(.62,-.24)  and +(.52,-.42).. ++(1.61,2.79) coordinate (B') coordinate [pos=.5] (M)
				--++(150:3.28) coordinate (A') coordinate [pos=1/5] (q1) coordinate [pos=4/5] (q2) coordinate [pos=1/4] (p1) coordinate [pos=3/4]  (p2)
				.. controls  +(.52,-.42) and +(.62,-.24).. ++(-1.61,-2.79);
			\draw [gray!70,fill=gray!60] (p1).. controls +(1.9,.47) and  +(-.54,1.88).. (p2)--(q2)
				.. controls +(-.07,2.05) and +(1.81,.96)..  (q1);
			\path  ($(A)!.5!(A')$) coordinate (O)
				($(B)!.5!(B')$) coordinate (O') 
				($2*(O')-(M)$) coordinate (N); 
			\begin{scope}
				\clip  (M)--(N)--(A)--++(-30:3.8);
				\fill [gray!40] (coc)++(-.79,-1.4)++(-30:3.28) 
					.. controls +(.62,-.24) and +(.52,-.42).. ++(1.61,2.79) 
					..  controls +(-.52,.42) and +(-.62,.24).. ++(-1.61,-2.79);
			\end{scope}
			\fill  [gray!40] (A)--(B)--(O');
			\fill [gray!60] (coc)++(-.79,-1.4).. controls  +(.29,-.25) and +(-.8,-.12).. (M)--(N)
				.. controls +(-.8,.2) and  +(.95,.1).. cycle;
			\draw [gray!80] (A).. controls +(-.62,.24) and  +(-.52,.42).. ++(1.61,2.79)
				.. controls +(.52,-.42) and +(.62,-.24).. ++(-1.61,-2.79); 
			\draw  [dashed,dash pattern=on 1pt off 1pt,gray] (B).. controls +(-.62,.24) and +(-.52,.42).. ++(1.61,2.79)
				(B)--(B') 
				 (A)--(O')
				 (M)--(N);
			\draw [|<->|] (B)++(-30:1)--++($(B')-(B)$) node [midway,sloped,fill=white,inner sep=1pt]  {$16$ cm};
			\draw [|<->|] (B')++(60:2)--++($(A)-(B)$) node [midway,sloped,fill=white,inner sep=1pt] {$10$ cm}; 
		\end{tikzpicture}
	} 
	\shortans{427} 
	\loigiai{ 
	\immini{ 
		Chọn hệ trục tọa  độ $Oxyz$ như hình bên.\\
		Coi phần nước còn lại trong cốc  như phần vật thể nằm giữa hai mặt phẳng $x=-8$ và $x=8$.\\
		Cắt phần nước theo một mặt phẳng vuông  góc với trục $Ox$ tại điểm có hoành độ $x$ ($-8\leq x\leq 8$) ta được thiết diện là tam giác $HKL$ vuông tại $K$.
		}{ 
		\begin{tikzpicture}[scale=0.7,>=stealth,line join=round,font=\scriptsize]
				\path (0,0) coordinate (coc)  coordinate (O);
				\begin{scope}[rotate=-60]
					\draw [gray,fill=gray!5] (coc)++(-.79,-1.4) coordinate (A)--++(-30:3.28)  coordinate (B)
						.. controls +(.62,-.24) and +(.52,-.42).. ++(1.61,2.79) coordinate (B') coordinate [pos=.18] (K) coordinate [pos=.45]  (M)
						--++(150:3.28) coordinate (A') coordinate [pos=1/5] (q1) coordinate  [pos=4/5] (q2) coordinate [pos=1/4] (p1) coordinate [pos=3/4] (p2)
						.. controls +(.52,-.42) and  +(.62,-.24).. ++(-1.61,-2.79);
					\draw [gray!70,fill=gray!60] (p1)..  controls +(1.9,.47) and +(-.54,1.88).. (p2)--(q2)
						.. controls +(-.07,2.05) and +(1.81,.96).. (q1); 
					 \path
						(A)++(0:{3.28/cos (30)}) coordinate (I) 
						++($(A)-(B)$)  coordinate (I')
						($(B)!.5!(B')$) coordinate (O') 
						(B)++(180:{3.28*cos  (30)}) coordinate (B1)
						(B).. controls +(-.62,.24) and  +(-.52,.42).. ++(1.61,2.79) coordinate [pos=.55] (N);
					\path [name path=l] (coc)++(-.79,-1.4).. controls  +(.29,-.25) and +(-.8,-.12).. (M);
					\path [name path=h] (K)--++($(A)-(B)$) coordinate  (h);
					\path  [name intersections={of=l and h,by={H}}];
					\path (intersection cs: first line={(K)--($(K)+(B')-(B)$)},second line={(M)--(N)}) coordinate  (L);
					\begin{scope}
						 \clip (M)--(N)--(A)--++(-30:3.8);
						\fill  [gray!40] (B)
						.. controls +(.62,-.24) and  +(.52,-.42).. ++(1.61,2.79)
						.. controls +(-.52,.42)  and +(-.62,.24).. ++(-1.61,-2.79);
					\end{scope}
					\fill  [gray!40] (A)--(B)--(O');
					\fill  [gray!60] (A).. controls +(.29,-.25) and +(-.8,-.12).. (M)--(N)
						..  controls +(-.8,.1) and +(.95,.2).. cycle;
					\draw [gray!80] (A).. controls  +(-.62,.24) and +(-.52,.42).. ++(1.61,2.79)
						.. controls +(.52,-.42) and +(.62,-.24)..  ++(-1.61,-2.79);
					\draw [dashed,dash pattern=on  1pt off 1pt,gray] (B).. controls +(-.62,.24) and +(-.52,.42).. ++(1.61,2.79)
						 (B)--(B')
						(A)--(O') 
						 (M)--(N);
					\foreach \x in  {A,B,M,N,H,K,L,O'}{
						\fill[gray] (\x) circle  (1pt);
					} 
					\draw [gray]  (H)--(K);
					\draw [dashed,dash pattern=on  1pt off 1pt,gray] (H)--(L)--(K)
						 (O)--(O')
						 (N)--++($(N)-(O')$);
					\fill  [pattern=north east lines,pattern color=gray] (H)--(K)--(L);
					\draw [->,gray] (O)--++(150:1.5) node  [left] {$z$};
					 \draw [->,gray] (M)--++($(M)-(N)$) node [below] {$x$};
					\draw [->,gray]  (B)--++($.5*(B)-.5*(I)$) node [below] {$y$};;
					\draw (O')  node [xshift=7pt,yshift=4pt]{$O$}
						(L) node [xshift=4pt,yshift=-2.5pt] {$x$} 
						(M) node [xshift=4pt,yshift=-4pt]  {$8$}
						(N) node  [xshift=10pt,yshift=3pt] {$-8$}
						(H)  node [yshift=4pt]{$H$}
						(K) node [xshift=-4pt,yshift=-4pt]{$K$} 
						(L) node [xshift=5pt,yshift=2.5pt]{$L$}; 
				\end{scope}
				 \fill (-1.5,-3.35) node[above left]{$B$};
				\fill (-1.5,0) node[left]{$A$}; 
			\end{tikzpicture}
		} 
		\noindent 
		Ta có 
		\begin{itemize}
			\item $KL=\sqrt{OK^2-OL^2}=\sqrt{8^2-x^2}=\sqrt{64-x^2}$. 
			\item $\widehat{AOB}=\widehat{HLK}$ (góc tạo bởi mặt  nước và mặt đáy cốc).
			 \item $HK=KL\cdot\tan\widehat{HLK}=KL\cdot\tan\widehat{AOB}=\sqrt{64-x^2}\cdot\dfrac{10}{8}=\dfrac{5\sqrt{64-x^2}}{4}$.
		\end{itemize}
		 Diện tích tam giác $HKL$ là
		\[S(x)=\dfrac{1}{2}\cdot HK\cdot KL  =\dfrac{1}{2}\cdot\dfrac{5\sqrt{64-x^2}}{4}\cdot\sqrt{64-x^2}=\dfrac{5}{8}\left(64-x^2\right).
		 \]
		Thể tích lượng nước trong cốc  là
		 \[V=\displaystyle\int_{-8}^{8}\dfrac{5}{8}\left(64-x^2\right)\,\mathrm{d}x=\dfrac{1\,280}{3}\approx427\,\left(\mathrm{cm^3}\right).\]
	} 
\end{bt}

% nguồn: Phan2_Lop_12_Dai_so_Chuong_4
\begin{bt}%[2D4C3-5]%[ma:ba120548]
	Hoạ sĩ thiết kế một micro có dạng khối tròn xoay, mặt cắt đứng chứa trục của khối tròn xoay có dạng như hình sau, trong đó $OA=OB=OI=2$ cm, $MC=MD=1$ cm, đường thẳng $OM$ là đường trung trực  của đoạn thẳng $CD$, $OM=20$ cm, $\widehat{AOB}=90^{\circ}$.
	\begin{center}
		\begin{tikzpicture}[>=stealth,line join=round,line cap=round,font=\footnotesize.0]
				 \def\r{1};
				 \path
					(0,0)  coordinate (O)
					(45:\r) coordinate (A) 
					 (-45:\r) coordinate (B)
					 (8,0.59) coordinate (C)
					(8,-0.59) coordinate (D) 
					 ($(C)!.5!(D)$) coordinate (M)
					(0:\r)  coordinate (I);
				\draw[dashed] (A) arc (45:-45:1); 
				 \draw (A)--(C)--(D)--(B)
					(A) arc  (45:315:1);
				 \draw[dashed] (B)--(O)--(A)
					 (O)--(M);
				\foreach \x/\g in  {C/60,B/-70,O/180,A/70,D/-90,I/60,M/0}
				\fill(\x) circle (1pt) +(\g:3mm)  node {$\x$};
			\end{tikzpicture}
	\end{center}
	Thể tích của micro này là  bao nhiêu cm$^3$? (Làm tròn kết quả đến hàng đơn~vị)
	\shortans{117} 
	\loigiai{ 
		\begin{center}
			\begin{tikzpicture}[>=stealth,line join=round,line cap=round,font=\footnotesize.0]
					\def\r{1}; 
					\path 
						 (0,0) coordinate (O)
						 (45:\r) coordinate (A)
						(-45:\r) coordinate  (B)
						(8,0.59) coordinate (C) 
						(8,-0.59) coordinate (D) 
						($(C)!.5!(D)$) coordinate  (M)
						(0:\r)  coordinate (I)
						($(A)!.5!(B)$)  coordinate (H);
					\draw[->] (-2,0) -- (9,0) node[below] {\scriptsize $x$}; 
					\draw[->] (0,-2) --  (0,2) node[left] {\scriptsize $y$};
					\draw (0.1,0.1)  node[below left]{\scriptsize $ $};
					\draw[dashed] (A) arc (45:-45:1); 
					 \draw (A)--(C)--(D)--(B)
						 (A) arc (45:315:1);
					\draw[dashed]  (B)--(O)--(A)
						(O)--(M) 
						 (A)--(B);
					\foreach \x/\g in {C/60,B/-70,O/220,A/70,D/-90,I/60,M/40,H/-120} 
					\fill(\x) circle (1pt)  +(\g:3mm) node {$\x$};
					\draw (M) node[below right]{\scriptsize $20$} 
						(0.12,1.12)  node{\scriptsize $2$}
						(0.12,-1.12) node{\scriptsize $-2$} 
						 (-1.2,-0.12) node{\scriptsize $-2$}
						(1.12,-0.12) node{\scriptsize $2$} 
						 (-1.2,0.12) node{\scriptsize $K$};
					\draw (-1,0) circle (1.0pt) 
						(1,0)  circle (1.0pt);
					 \draw pic[draw,angle radius=2mm] {right angle = A--H--O};
				\end{tikzpicture}
		\end{center}
		Đặt hệ  trục $Oxy$ như hình vẽ, đơn~vị của mỗi trục là $1$ cm.\\
		Gọi $y=f(x)$ là đồ thị bao gồm cung nhỏ $KA$  và đoạn thẳng $AC$.\\
		Gọi  $H$ là trung điểm $AB$.\\
		Vì  $\triangle OAH$ vuông cân nên $OH=AH=\sqrt{2}$ suy ra $A\left(\sqrt{2};\sqrt{2}\right)$.
		 $$V=\pi\displaystyle\int^{20}_{-2}\left[f(x)\right]^2\,\mathrm{\,d}x=\pi\displaystyle\int^{\sqrt{2}}_{-2}\left[f(x)\right]^2\,\mathrm{\,d}x+\pi\displaystyle\int^{20}_{\sqrt{2}}\left[f(x)\right]^2\,\mathrm{\,d}x.$$
		Vì $A\left(\sqrt{2};\sqrt{2}\right)$, $C(20;1)$  nên phương trình đường thẳng $AC$ là
		\[\dfrac{x-20}{\sqrt{2}-20}=\dfrac{y-1}{\sqrt{2}-1}\Leftrightarrow  y-1=(\sqrt{2}-1)\dfrac{x-20}{\sqrt{2}-20}\Leftrightarrow y=\dfrac{(1-\sqrt{2})x+19\sqrt{2}}{20-\sqrt{2}}.\]
		Phương trình đường tròn là $x^2+y^2=4$, hàm số có  đồ thị cung nhỏ $KA$ là
		$$y=\sqrt{4-x^2},\,x\in  [-2;\sqrt{2}].$$
		Khí đó 
		 \allowdisplaybreaks
		\begin{eqnarray*}
			V&=&\pi\displaystyle\int^{\sqrt{2}}_{-2}\ {\left(\sqrt{4-x^2}\right)}^2\,\mathrm{\,d}x+\pi\displaystyle\int^{20}_{\sqrt{2}}\  {\left[\dfrac{(1-\sqrt{2})x+19\sqrt{2}}{20-\sqrt{2}}\right]}^2\,\mathrm{\,d}x\\
			&=&\pi\displaystyle\int^{\sqrt{2}}_{-2}\ \left(4-x^2\right)\,\mathrm{\,d}x+\pi\displaystyle\int^{20}_{\sqrt{2}}\ {\left[\dfrac{(1-\sqrt{2})x+19\sqrt{2}}{20-\sqrt{2}}\right]}^2\,\mathrm{\,d}x\approx  117\,\left(\text{cm}^3\right).
		\end{eqnarray*}
	 }
\end{bt}

% nguồn: Phan2_Lop_12_Dai_so_Chuong_4
\begin{bt}%[2D4C3-5]%[ma:d8dffd65]
	Lượng mưa theo giờ, tính bằng inch/giờ, ở hai địa điểm khác nhau sau $t$ giờ khi bão đổ bộ, được cho bởi các hàm số $f(t) = 0{,}73t^3-2t^2+t+0{,}6$ và $g(t)=0{,}17t^2-0{,}5t+1{,}1$. Tính diện tích hình phẳng giới hạn bởi đồ thị hàm số của hai hàm số đã cho và hai đường thẳng $t=0$, $t=2$, từ đó tính sự chênh lệch lượng mưa ở hai địa điểm khác nhau sau $2$ giờ. (Làm tròn đến hàng phần~trăm). 
	\shortans{0,87} 
	\loigiai{ 
	Diện tích hình phẳng  giới hạn bởi đồ thị hàm số $f(t) = 0{,}73t^3-2t^2+t+0{,}6$; $g(t)=0{,}17t^2-0{,}5t+1{,}1$ và hai đường thẳng $t = 0$, $t=2$ là
	\begin{align*}
		\displaystyle \int_{0}^{2} \left|f(t)-g(t)\right| \mathrm{\,d}t = \int_{0}^{2} \left|0{,}73t^2 + 1{,}5t - 0{,}5\right| \mathrm{\,d}t  = \dfrac{13}{15} \approx 0{,}87.
	\end{align*}
	 }
\end{bt}

% nguồn: Phan2_Lop_12_Dai_so_Chuong_4
\begin{bt}%[2D4C3-5]%[ma:7de07636]
	\immini{Một chiếc đồng hồ cát như hình vẽ gồm hai phần đối xứng nhau qua mặt phẳng nằm ngang và đặt trong một hình trụ. Thiết diện thẳng đứng qua trục của nó là hai parabol chung đỉnh và đối xứng nhau qua mặt phẳng nằm ngang. Ban đầu lượng cát dồn hết ở phần trên của đồng hồ thì chiều cao của mực cát bằng $\dfrac{2}{3}$ chiều cao của bên đó (xem hình vẽ). Cát chảy từ trên xuống dưới với tốc độ $v(t)=0{,}2t+13$ (cm$^3$/phút). Biết sau $20$ (phút) thì cát chảy hết xuống phần bên dưới của đồng hồ. Khi chiều cao của cát còn $4$\,(cm) thì bề mặt trên cùng của cát tạo thành một đường tròn có chu vi bằng $8\pi$\,(cm). Hỏi chiều cao của khối trụ bên ngoài bằng bao nhiêu centimet (\textit{nếu kết quả là số thập phân  thì làm tròn đến hàng đơn vị})?}{
	\begin{tikzpicture}[line cap=butt,line join=miter,>=stealth,scale=.7,yscale=.7]
			\tikzset{declare function={a=1;  b=0; f (\x)=a*(\x)^2;}}
			\begin{scope}
				\foreach[count=\i]  \goc/\t in {0/A,180/B}{
					\begin{scope}[rotate=\goc]
					\draw[samples=100] plot[domain=-2:2]  (\x,{f (\x)});
					\draw (0,4) ellipse (2cm and 0.6cm); 
					\end{scope}
				 }
			\end{scope}
			\draw[dashed] (1,1)  arc (0:180:1cm and 0.4cm);
			\draw (1,1) arc  (0:-180:1cm and 0.4cm);
			\draw[dashed] (1.5,{f (1.5)}) arc (0:180:1.5cm and 0.4cm); 
			 \draw (1.5,{f (1.5)}) arc (0:-180:1.5cm and 0.4cm);
			\draw  (-2,4)--(-2,-4)
				(2,4)--(2,-4); 
			 \draw[<->](0,0)--(0,1);
			\draw[<->](0.5,0)--(.5,{f  (1.5)});
			\draw ($(0.5,0)!2/3!(.5,{f (1.5)})$) node[right]{\small{$h$}} 
				 ($(0,0)!0.5!(0,1)$) node[right]{\small{$4$}};
			 \draw[->](-0.2,{1-.4})--(-0.9,0.1) node[below]{\small{$8\pi$}};
		\end{tikzpicture}
	 }
	 \shortans{21}
	\loigiai{ 
		\begin{center}
			\begin{tikzpicture}[scale=1,>=stealth,font=\scriptsize,line join=round,line cap=round]
					\draw[-stealth] (-5,0) --  (5,0) node[right] {$x$};
					\draw[-stealth] (0,-0.5) --  (0,5) node[above] {$y$};
					\node[below right] at (0,0)  {$O$};
					\draw[samples=200,domain=-4.5:4.5] plot (\x,{0.25*\x*\x}); 
					\fill  (2.8,2) circle (1pt) node[right] {$A$};
					\fill (-2.8,2) circle (1pt) node[left]  {$B$};
					\draw[dashed] (2.8,0) --  (2.8,2);
					\draw[dashed]  (-2.8,0) -- (-2.8,2);
					 \node[below] at (2.8,0) {$4$};
					\node[below] at (-2.8,0) {$-4$}; 
					\node at (0,2)  [left] {$4$};
					\draw[xscale=0.8,blue,dashed] plot[samples=200,domain=-3.5:3.5] (\x,{2  + 0.5*sqrt (1-(\x/3.5)^2)});
					 \draw[xscale=0.8,blue] plot[samples=200,domain=-3.5:3.5] (\x,{2 - 0.5*sqrt (1-(\x/3.5)^2)});
					\fill[blue]  (0,2) circle (1pt);
					\fill[blue]  (2.8,2) circle (1pt);
					\draw[red,dashed] plot[samples=200,domain=-4:4] (\x,{4 + 0.6*sqrt (1-(\x/4)^2)}); 
					\draw[red] plot [samples=200,domain=-4:4] (\x,{4 -  0.6*sqrt (1-(\x/4)^2)});
					\fill[red] (0,4)  circle (1pt);
					\fill[red] (4,4) circle (1pt); 
				\end{tikzpicture}
		\end{center}
		Gọi $(P)$ là parabol phía trên.\\ Chọn hệ trục tọa độ $Oxy$ sao cho đỉnh $O$ trùng với đỉnh của parabol và trục $Oy$ là trục đối xứng của parabol.\\ Khi chiều cao của cát là $4$ (cm), chu vi đường tròn bề mặt cát là $8\pi$ (cm) nên bán kính đường tròn là $r=\dfrac{8\pi}{2\pi}=4$ (cm).\\ 
		Vậy  parabol $(P)$ đi qua điểm $A(4;4)$, thay vào phương trình $$y=ax^2\Rightarrow 4=a\cdot4^2\Rightarrow a=\dfrac{1}{4}.$$
		Khi đó phương trình  của parabol $(P)\colon y=\dfrac{x^2}{4}$.\\
		Thể tích  cát ban đầu bằng thể tích khối tròn xoay sinh ra khi quay nhánh phải của parabol $(P)$ quanh trục $Oy$. Thể tích này cũng bằng lượng cát chảy xuống trong $20$ phút.\\
		Thể tích  ban đầu là $V=\pi\displaystyle\int_{0}^{h}x^2\,\mathrm{d}y=\pi\displaystyle\int_{0}^{h}4y\,\mathrm{d}y=2\pi h^2$.\\
		Lượng cát chảy ra trong $20$ phút bằng  $V=\displaystyle\int_{0}^{20}v(t)\,\mathrm{d}t=\displaystyle\int_{0}^{20}(0{,}2t+13)\,\mathrm{d}t=300$.\\
		Do đó $2\pi h^2=300\Leftrightarrow h^2=\dfrac{300}{2\pi}\Rightarrow  h=\sqrt{\dfrac{150}{\pi}}$.\\
		Chiều cao của mực cát ban đầu bằng $\dfrac{2}{3}$ chiều  cao của một bên parabol, tức là $h=\dfrac{2}{3}H$ với $H$ là chiều cao một bên của parabol.\\
		Vậy chiều cao của khối  trụ bên ngoài là $l=2\cdot\dfrac{3}{2}\cdot\sqrt{\dfrac{150}{\pi}}\approx21$\,(cm).
	} 
\end{bt}

% nguồn: Phan2_Lop_12_Dai_so_Chuong_4
\begin{bt}%[2D4C3-5]%[ma:fe1d541a]
	Xét phần đựng nước bên trong một lục bình có dạng mặt tròn xoay trong đó đường sinh là đường cong $(C)$ quay quanh trục $(d)$ như hình bên. Biết rằng $(C)$ là đường cong bậc ba; chiều cao phần đựng nước là 35 cm; bán kính đáy phần đựng nước là 10 cm; bán kính miệng phần đựng nước là $14{,}375$ cm; chỗ rộng nhất của phần đựng nước có bán kính $30$ cm; độ cao của phần đựng nước tính từ đáy đến chỗ rộng nhất  là $10$ cm. Thể tích tối đa mà chiếc bình này đựng được là bao nhiêu lít? (kết quả làm tròn đến hàng phần chục)
	\begin{center}
		\begin{tikzpicture}
				\path (-2.5,5) coordinate (At) 
					(-2.5,-1) coordinate  (Ad)
					 (2.5,5) coordinate (Bt)
					(2.5,-1) coordinate  (Bd)
					($(0,4)+(205:1.1 and  0.33)$) coordinate (Md)
					($(0,4)+(25:1.1 and 0.33)$) coordinate (Mt) 
					(205:1 and 0.3) coordinate  (Nd)
					(25:1 and 0.3) coordinate  (Nt)
					($(Nt)!-0.25cm!(Nd)$)  coordinate (Nt)
					(intersection of Md--Mt and At--Ad)  coordinate (A)
					(intersection of Md--Mt  and Bt--Bd) coordinate (B)
					(intersection of Nd--Nt and At--Ad)  coordinate (C)
					(intersection of Nd--Nt and Bt--Bd) coordinate  (D);
				\draw[dashed] (1,0) arc (0:180:1 and  0.3)
					($(0,1.5)+(0:1.5)$) arc (0:205:1.5 and 0.4) 
					(-1.1,4).. controls +(-90:0.4)  and +(87:0.6)..
					 (-0.6,3).. controls +(-94:0.3) and +(90:0.6)..
					 ($(0,1.5)+(180:1.52)$).. controls +(-90:0.7) and +(114:0.5)..
					 (-1,0)
					 (Nd)--(Nt)
					 (0,-0.25)--(0,4);
				\draw 
					(1.1,4)..  controls +(-90:0.4) and +(93:0.6)..
					 (0.6,3).. controls +(-86:0.3) and +(90:0.6)..
					($(0,1.5)+(0:1.55)$)..  controls +(-90:0.7) and +(66:0.5)..
					 (1,0)
					(A)--(C)--(Nd) 
					 (Nt)--(D)--(B)--(A)
					($(0,1.5)+(205:{1.5 and  0.4})$) arc (200:360:1.5 and 0.4)
					(-1,0) arc  (180:360:1 and 0.3)
					(0,4) ellipse  (1.1 and 0.33)
					(Md).. controls +(-40:1.7) and +(83:1).. (-1.21,1.6).. controls +(-97:0.7)  and +(0:0).. (Nd)
					(0,-.32)--+(0,-.3) 
					 (0,4)--+(0,0.7) node[right]{$d$};
				\fill (0,0) circle  (1pt)
					(0,4)  circle (1pt);
			\end{tikzpicture}
	\end{center}
	\shortans{49,1} 
	 \loigiai{
		Đổi tất  cả sang đơn~vị về đềximét.
		\begin{center}
			\begin{tikzpicture}
					\path (-2.5,5) coordinate  (At)
						(-2.5,-1) coordinate  (Ad)
						(2.5,5) coordinate (Bt) 
						(2.5,-1) coordinate (Bd) 
						($(0,4)+(205:1.1 and 0.33)$) coordinate (Md) 
						($(0,4)+(25:1.1 and 0.33)$) coordinate (Mt) 
						(205:1 and 0.3) coordinate  (Nd)
						(25:1 and  0.3) coordinate (Nt)
						 ($(Nt)!-0.25cm!(Nd)$) coordinate (Nt)
						(intersection of Md--Mt and At--Ad) coordinate  (A)
						(intersection of Md--Mt and Bt--Bd) coordinate (B) 
						(intersection of  Nd--Nt and At--Ad) coordinate (C)
						(intersection of Nd--Nt and Bt--Bd)  coordinate (D);
					\draw (0,0)--(0,4)--+(0,0.7) node[right]{$d$}; 
					\draw[<->] (0,4)--(1.1,4); 
					 \draw[<->] (0,1.5)--(1.55,1.5);
					 \draw[<->] (0,0)--(1.0,0);
					 \draw[<->] (-0.5,0)--(-0.5,4);
					 \draw[<->] (-0.2,0)--(-0.2,1.5);
					\draw 
						(1.1,4).. controls  +(-90:0.4) and +(93:0.6)..
						(0.6,3).. controls  +(-86:0.3) and +(90:0.6)..
						($(0,1.5)+(0:1.55)$).. controls  +(-90:0.7) and +(66:0.5)..
						 (1,0);
					\fill (0,0)  circle (1pt)
						(0,4) circle  (1pt);
					\draw (-0.5,2)  node[left]{\scriptsize $3{,}5$}
						(0.5,4) node[above]{\scriptsize $1{,}4375$} 
						(0.75,1.45)  node[above]{\scriptsize $3$}
						(0.5,0) node[below]{\scriptsize $1$} 
						(0.1,0.75) node[left]{\scriptsize  $1$}
						 (1,2.5) node[right]{\scriptsize $(C)$};
				\end{tikzpicture}
		\end{center}
		Ta có thể minh hoạ đường sinh $(C)$  và trục quay $d$ như hình bên.\\
		Gắn hệ trục như hình  vẽ (trục quay $d$ là trục hoành, trục tung cắt đường sinh tại điểm thuộc phần rộng nhất).
		\begin{center}
			\begin{tikzpicture}[scale=1.0,font=\footnotesize,line join=round,line cap=round,>=stealth];
					\draw[->]  (-3,0) -- (2,0) node[below] {\scriptsize $x$};
					\draw[->] (0,-0.5) -- (0,4) node[left]  {\scriptsize $y$};
					 \draw (0.1,0.1) node[below left]{\scriptsize $O$};
					\draw[thin,samples=150,smooth,domain=-2.5:1] plot (\x,{-0.5*((\x)^3)-1.5*((\x)^2)-0*(\x)+3}); 
					\draw[dashed] (-2.5,0) node[below]{$-2{,}5$} --  (-2.5,1.4375)--(0,1.4375) node[right]{$1{,}4375$};
					\draw[dashed] (1,0) node[below]{$1$} -- (1,1)--(0,1)  node[left]{$1$}
						(0,3) node[above left]{$3$}; 
					\draw (-2.5,1.4375)  circle (1.2pt)
						(1,1) circle  (1.2pt)
						 (0,3) circle (1.2pt);
				\end{tikzpicture}
		\end{center}
		Đặt $(C)\colon y=f(x)=ax^3+bx^2+cx+d$. Dựa vào đồ thị  ta thấy
		\[\heva{&f'(0)=0\\&f(0)=3\\&f(1)=1\\&f(-2{,}5)=1{,}4375}\Leftrightarrow \heva{&c=0\\&d=3\\&a+b+c+d=1\\&-15{,}625a+6{,}25b-2{,}5c+d=1{,}4375}\Leftrightarrow  \heva{&a=-0{,}5\\&b=-1{,}5\\&c=0\\&d=3.}\]
		Suy  ra $y=f(x)=-0{,}5x^3-1{,}5x^2+3$.\\
		Khi đó thể tích phần chứa nước chính là thể tích khối tròn xoay tạo thành khi quay phần hình phẳng bị giới hạn bởi $(C)\colon y=f(x)$, trục hoành, các đường thẳng  $x=-2{,}5$ và $x=1$ quanh trục hoành.\\
		Do đó thể tích  này bằng $V=\pi\displaystyle\int^1_{-2{,}5}{\left(-0{,}5x^3-1{,}5x^2+3\right)}^2\,\mathrm{\,d}x\approx 49{,}1$ (lít).
	 }
\end{bt}


%%%%%%%%%%%%%%%% TỰ LUẬN %%%%%%%%%%%%%%%%

% nguồn: Phan2_Lop_12_Dai_so_Chuong_4
\begin{bt}%[2D4H3-5]%[ma:ead322b2]
	Một bình chứa nước  có hình dạng như hình sau. Biết rằng khi nước trong bình có chiều cao $x$ (dm) ($0 \le x \le 4$) thì mặt nước là hình vuông có cạnh $\sqrt{2+\dfrac{x^2}{4}}$ (dm). Tính dung tích của bình.
	\begin{center}
		\foreach \m in  {70}{
		\begin{tikzpicture}[scale=0.8]

			\def\m{70} 
			 \tdplotsetmaincoords{70}{\m}

			 \clip (-4,-1) rectangle (6,6);
			\pgfmathsetmacro{\l}{sqrt (2)} 
			\begin{scope}[tdplot_main_coords,line join=round,line cap=round,magenta]
				 \def\xx{2.5}
				\path (0,{\l+sqrt (2+(\xx)^2/4)},\xx) coordinate (A) 
					 ({\l+sqrt (2+(\xx)^2/4)},0,\xx) coordinate (B)
					(0,{-\l-sqrt  (2+(\xx)^2/4)},\xx) coordinate (C)
					({-\l-sqrt  (2+(\xx)^2/4)},0,\xx) coordinate (D);
				\draw[samples=100,domain=0:4]plot (0,{\l+sqrt (2+(\x)^2/4)},\x)  coordinate (a);
				\draw[samples=100,domain=0:4]plot (0,{-\l-sqrt  (2+(\x)^2/4)},\x) coordinate (b);
				\draw[dashed,samples=100,domain=0:\xx]plot ({-\l-sqrt  (2+(\x)^2/4)},0,\x);
				\draw[samples=100,domain=\xx:4]plot  ({-\l-sqrt (2+(\x)^2/4)},0,\x) coordinate (c);
				 \draw[samples=100,domain=0:4]plot ({\l+sqrt (2+(\x)^2/4)},0,\x) coordinate (d);
				\draw  (a)--(c)--(b)--(d)--cycle;
				\fill[cyan,opacity=0.25](A)--(D)--(C)--plot[samples=100,domain=\xx:0](0,{-\l-sqrt (2+(\x)^2/4)},\x) coordinate (m)--({\l+sqrt (2+(0)^2/4)},0,0) coordinate (n)--(0,{\l+sqrt (2+(0)^2/4)},0) coordinate  (p)--plot[samples=100,domain=0:\xx](0,{\l+sqrt (2+(\x)^2/4)},\x);
				 \draw[cyan](A)--(B)--(C)--(D)--cycle ({-\l-sqrt (2+(0)^2/4)},0,0) coordinate (q);
				\draw (m)--(n)--(p); 
				 \draw[dashed](m)--(q)--(p);
				 \draw[dashed](0,0,0) coordinate (I)--({sqrt (2)},{sqrt (2)},0)
					(0,0,\xx)--($(A)!0.5!(B)$)  coordinate (N);
				\draw (N)--(4,4,\xx); 
				 \path (0,0,4) coordinate (I')
					(intersection of I--I'  and d--b) coordinate (M);
				\draw[dashed](I)--(M); 
				\draw (M)--(I')--(4,4,4) 
					 ($(n)!0.5!(p)$)--(4,4,0);
				\draw[<-stealth](3.5,3.5,0)--(3.5,3.5,4) node[scale=0.8,right,pos=0.5]{$4$ dm}; 
				\draw[<-stealth](2.5,2.5,0)--(2.5,2.5,\xx)  node[scale=0.8,right,pos=0.5]{$x$ dm};
			\end{scope}
		\end{tikzpicture}
		} 
	\end{center}
	 \loigiai{
		Tại độ cao $x$, mặt nước là một hình vuông có độ dài  cạnh là $\sqrt{2+\dfrac{x^2}{4}}$.
		Diện tích của mặt cắt ngang tại độ cao $x$ là 
		$$ A(x) = \left(\sqrt{2+\frac{x^2}{4}}\right)^2 = 2+\frac{x^2}{4} \quad (\text{dm}^2)  .$$
		Bình có chiều cao từ $x=0$  đến $x=4$, do đó dung tích của bình là
		$$ V =  \int_{0}^{4} A(x) \mathrm{d}x = \int_{0}^{4} \left(2+\frac{x^2}{4}\right) \mathrm{d}x = \frac{40}{3} .$$
		Vậy, dung tích  của bình là \textbf{$\dfrac{40}{3}$} dm$^3$ (tức $\dfrac{40}{3}$ lít).
	} 
\end{bt}

% nguồn: Phan2_Lop_12_Dai_so_Chuong_4
\begin{bt}%[2D4H3-5]%[ma:4eb900a6]
	Một cốc rượu có hình dạng tròn xoay và kích thước như hình vẽ, thiết diện dọc của cốc bổ dọc cốc thành 2 phần bằng nhau là một đường Parabol. Tính thể tích tối đa mà cốc có thể chứa được (làm tròn  2 chữ số thập phân).
	\begin{center}
		\begin{tikzpicture}[scale=0.5]
				\draw[domain=-3:3,smooth,variable=\x] plot  ({\x},{0.625*\x*\x});
				\draw (0,-2) ellipse (1.2cm and  0.2cm);
				\draw (0,{0.625*(-3)*(-3)}) ellipse (3cm and 0.8cm); 
				\draw (-0.2,-2)  -- (-0.2,{0.625*(-0.2)*(-0.2)});
				\draw (0.2,-2) -- (0.2,{0.625*(0.2)*(0.2)}); 
				\draw[<->] (-3,7.5) -- (3,7.5)  node[midway,above] {8 cm};
				\draw[<->] (4.5,0)  -- (4.5,{0.625*(-3)*(-3)}) node[midway,right] {10 cm};
			\end{tikzpicture}
	\end{center}
	 \loigiai{
		\begin{center}
			\begin{tikzpicture}[scale=0.4]
					\draw[domain=-4.2:4.2,smooth,variable=\x,blue] plot  ({\x},{0.625*\x*\x});
					\draw[->] (-4.5,0) -- (4.5,0)  node[right] {$x$};
					\draw[->] (0,-1) -- (0,11) node[above]  {$y$};
					 \foreach \x/\point in {-4/M,4/N}{
						\filldraw (\x,{0.625*\x*\x}) circle (2pt) node[above right]  {$\point$};
						 \ifnum\x=0
						\node[below left] at (\x,{0.625*\x*\x}) {$\point$}; 
						\else 
						\node[above right] at  (\x,{0.625*\x*\x}) {$\point$};
						 \fi
						\draw[dashed] (-4,0) -- (-4,{0.625*(-4)*(-4)}) -- (0,{0.625*(-4)*(-4)}); 
						\draw[dashed] (4,0) --  (4,{0.625*(-4)*(-4)}) -- (0,{0.625*(-4)*(-4)});
					} 
					\filldraw (0,0) circle (2pt) node[below left]  {$O$};
					\filldraw (-4,0) circle  (2pt) node[below] {$-4$};
					\filldraw (4,0)  circle (2pt) node[below] {$4$};
					\filldraw (0,{0.625*(-4)*(-4)}) circle (2pt)  node[below left] {$10$};
				\end{tikzpicture}
		\end{center}
		Parabol có  phương trình $y=\dfrac{5}{8}x^2\Leftrightarrow x^2=\dfrac{8}{5}y$. \\
		Thể tích tối đa  cốc $V=\pi\displaystyle\int_0^{10} \dfrac{8}{5}y\mathrm{\,d}y=\dfrac{8\pi y^2}{10}\big|_0^{10}\approx 251{,}3 $.
	 }
\end{bt}

% nguồn: Phan2_Lop_12_Dai_so_Chuong_4
\begin{bt}%[2D4V3-5]%[ma:e7a7e2ba]
	Một Chi đoàn thanh niên đi dự trại ở một đơn vị bạn, họ dự định dựng một lều trại có dạng parabol (nhìn từ mặt trước, lều trại được căng thẳng từ trước ra sau, mặt sau trại cũng là parabol có kích thước giống như mặt trước) với kích thước: nền trại là một hình chữ nhật có chiều rộng là 3 mét, chiều sâu là $6$ mét, đỉnh của parabol cách mặt đất là $3$ mét. Hãy tính thể tích (đơn vị: m$^3$) phần không gian phía  trong trại để cử số lượng người tham dự trại cho phù hợp.
	\begin{center}
		\begin{tikzpicture}[scale=0.65,font=\footnotesize,line join=round,line cap=round,>=stealth]
				 \def\f{-7/6*(\x)^2+5/6*(\x)+3}
				\def\g{-7/6*(\x-6)^2+5/6*(\x-6)+5} 
				 \def\s{-7/6*(\x-3)^2+5/6*(\x-3)+4}
				\coordinate (A)  at (-1,1);
				\coordinate (B)  at (2,0);
				\coordinate (C) at (5,3); 
				\coordinate (D) at (8,2); 
				\coordinate  (I) at (5/14,529/168);
				\coordinate (J)  at (89/14,865/168);
				\draw[smooth,samples=200,domain=-1:2] plot  (\x,{\f});
				\draw[smooth,samples=200,domain=5:8]  plot (\x,{\g});
				\draw  (A)--(B)--(D)
					 (I)--(J);
				\draw[dashed] (A)--(C)--(D); 
				 \draw[<->](I)--(0.5,0.5);
				\draw  (0.5,0.5) node[below]{$3$}
					 (0.5,2) node[left]{$3$};
				\draw  (5,1) node[below]{$6$};
			\end{tikzpicture}
	\end{center}
	\loigiai{ 
		Gắn hệ trục tọa độ $Oxy$ vào parabol sao cho đỉnh của parabol là $I(0; 3)$, parabol cắt trục $Ox$ tại $2$ điểm là $A\left(-\dfrac{3}{2};  0\right)$, $B\left(\dfrac{3}{2}; 0\right)$. Tìm được parabol $(P)\colon y=-\dfrac{4}{3}{x^2}+3$.\\
		Cắt lều trại bằng mặt phẳng song song với mặt phẳng chứa parabol ta được thiết diện có diện tích $S$  bằng diện tích của hình phẳng giới hạn bởi $(P)$ và trục $Ox$ ta có $S=\displaystyle\int_{-\tfrac{3}{2}}^{\tfrac{3}{2}}{\left|-\dfrac{4}{3}{x^2}+3\right|\,\mathrm{d}x}=6$.\\
		Khi đó thể tích phần không gian phía trong trại là $V=\displaystyle\int_0^66\,\mathrm{d}x=36$  m$^3$.
	 }
\end{bt}

% nguồn: Phan2_Lop_12_Dai_so_Chuong_4
\begin{bt}%[2D4V3-5]%[ma:d3d2dfd8]
	Nghệ thuật làm gốm có lịch sử  phát triển lâu đời và vẫn còn tồn tại cho đến ngày nay. Giả sử một bình gốm có mặt trong của bình là một mặt tròn xoay sinh ra khi cho phần đồ thị của hàm số $y=\dfrac{1}{175} x^2+\dfrac{3}{35} x+5$ $\left(0 \leq x \leq 30\right)$ ($x$, $y$ tính theo cm) quay tròn quanh bệ gốm có trục trùng với trục hoành $O x$. Hỏi để hoàn thành bình gốm đó ta cần sử dụng bao nhiêu cm$^3$ đất sét, biết rằng bình gốm đó có độ dày không đổi là $1$ cm.
	\loigiai{ 
		Thể tích phần  bên trong của bình gốm là
		\[V_1=\pi \displaystyle \int_0^{30} \left(\dfrac{1}{175} x^2+\dfrac{3}{35} x+5\right)^2 \mathrm{\,d}x=240\pi  \, (\text{cm}^3).\]
		Thể tích của bình gốm bao gồm cả  độ dày là
		\[V_2=\pi \displaystyle \int_0^{30} \left(\dfrac{1}{175} x^2+\dfrac{3}{35} x+5+1\right)^2 \mathrm{\,d}x=270\pi  \, (\text{cm}^3).\]
		Thể tích cần dùng  để hoàn thành bình gốm là
		\[V=V_2-V_1=270 \pi-240\pi=30 \pi  \, (\text{cm}^3).\]
	 }
\end{bt}
